% Microorganismes dans notre environnement — SVTEEHB 3ème — Cours signé OnBuch+
% Programme officiel MINESEC. Compiler avec : tectonic N04-microorganismes-notre-environnement.tex
\newcommand{\DOCMATIERE}{SVTEEHB}
\newcommand{\DOCNIVEAU}{3ème}
\newcommand{\DOCMODULE}{SVTEEHB – classe de 3ème}
\newcommand{\DOCLECON}{N04}
\newcommand{\DOCTITRE}{Microorganismes dans notre environnement}
% OnBuch+ — préambule « cours signés » ENRICHI (compilé avec tectonic / XeLaTeX)
% Système visuel « Pop » d'OnBuch+ : fond crème (jamais blanc), bordures encre
% épaisses, ombres « dures » décalées, titres Archivo Black en capitales.
% Version MATHÉMATIQUES : graphes (pgfplots/tikz), boîtes pédagogiques
% (définition, propriété, méthode, attention, le savais-tu, exercice résolu,
% exercice, TP, bilan), page de garde et signature OnBuch+ ancrée en bas.
%
% Macros attendues dans main.tex AVANT \input{preamble} :
%   \DOCMATIERE  \DOCNIVEAU  \DOCMODULE  \DOCLECON (numéro)  \DOCTITRE
\documentclass[11pt]{article}

\usepackage{fontspec}
\usepackage[french]{babel}
\usepackage[a4paper,margin=2cm,top=3.4cm,bottom=2.8cm,headheight=1.4cm,footskip=1.3cm]{geometry}
\usepackage{amsmath,amssymb,amsfonts}
\usepackage{mathtools} % \xmapsto, \coloneqq... souvent utilisés par les modèles
\usepackage{cancel} % simplifications barrées (bilans de forces)
% \abs{z} (module, valeur absolue) et \norm{u} (norme), utilisés par les modèles
\providecommand{\abs}{}\let\abs\relax\DeclarePairedDelimiter{\abs}{\lvert}{\rvert}
\providecommand{\norm}{}\let\norm\relax\DeclarePairedDelimiter{\norm}{\lVert}{\rVert}
\usepackage[table]{xcolor} % [table] : fournit aussi \rowcolors (alternance de lignes)
\usepackage{pagecolor}
\usepackage{tikz}
\usetikzlibrary{arrows.meta,positioning,calc,shapes.geometric,shapes.misc,shapes.arrows,decorations.pathmorphing,decorations.pathreplacing,decorations.markings,patterns,shadows,mindmap,trees,fit,backgrounds,matrix,intersections,angles,quotes}
\usepackage{pgfplots}
\usepackage[version=4]{mhchem} % équations nucléaires \ce{^{235}_{92}U}
\usepackage[european,siunitx]{circuitikz} % schémas électriques
\pgfplotsset{compat=1.18}
\usepgfplotslibrary{fillbetween,groupplots} % aires entre courbes (calcul intégral)
\usepackage{enumitem}
\usepackage{fancyhdr}
% [most] charge toutes les bibliothèques tcolorbox (listings, minted, raster,
% documentation...) et gonfle inutilement la mémoire TeX sur un document long
% (~300 boîtes) : on ne charge que ce qui est réellement utilisé.
\usepackage[skins,breakable]{tcolorbox}
\usepackage{graphicx}
\usepackage{colortbl}
\usepackage{booktabs}
\usepackage{tabularx}
\usepackage{multirow}
\usepackage{diagbox} % case d'en-tête diagonale des tableaux à double entrée
% Colonnes extensibles centrée / gauche / droite (C, L, R), souvent utilisées par les modèles
\newcolumntype{C}{>{\centering\arraybackslash}X}
\newcolumntype{L}{>{\raggedright\arraybackslash}X}
\newcolumntype{R}{>{\raggedleft\arraybackslash}X}
\usepackage{array}
\usepackage{multicol}
\usepackage{siunitx}
% parse-numbers=false : accepte aussi \SI{2,5 \times 10^{-3}}{...}
\sisetup{locale=FR,output-decimal-marker={,},group-minimum-digits=4,parse-numbers=false}
\DeclareSIUnit\ohm{\ensuremath{\symup{\Omega}}} % U+2126 absent de la police
\DeclareSIUnit\division{div} % réglages d'oscilloscope (ms/div, V/div)
\DeclareSIUnit\tr{tr} % tours (vitesse de rotation, ex. \SI{1500}{\tr\per\minute})
\DeclareSIUnit\molar{M} % concentration molaire (ex. \SI{3}{\molar}, \SI{1}{\micro\molar})
\DeclareSIUnit\an{an} % année (le modèle confond parfois avec \year, un compteur TeX) — ex. \SI{5}{\kilo\meter\per\an}
\DeclareSIUnit\FCFA{FCFA} % franc CFA (ex. \SI{500}{\FCFA\per\kilo\gram})
\DeclareSIUnit\degreeBrix{\textdegree Brix} % teneur en sucre d'un sirop/jus (réfractométrie)
\DeclareSIUnit\month{mois} % le modèle utilise parfois \month, un compteur TeX, comme unité de durée
\DeclareSIUnit\microliter{\micro\liter} % le modèle écrit parfois \microliter en un seul token
\DeclareSIUnit\million{millions} % dénombrement (ex. \SI{15}{\million\per\mL} spermatozoïdes)
\usepackage{qrcode}
% \degree : commande fréquemment utilisée par les modèles pour un angle ou une
% température en dehors de \SI{}{} (non fournie par les paquets chargés).
\providecommand{\degree}{^{\circ}}
% \qed (fin de démonstration), utilisé par les modèles sans amsthm
\providecommand{\qed}{\hfill$\square$}
% \barre : double barre des valeurs interdites dans un tableau de variations (tkz-tab)
\providecommand{\barre}{\ensuremath{\|}}
% environnement proof (amsthm non chargé) utilisé par les modèles
\ifdefined\proof\else\newenvironment{proof}[1][Démonstration]{\par\noindent\textit{#1.}\ }{\qed\par}\fi
\usepackage{lastpage}
\usepackage{hyperref}
\hypersetup{hidelinks}

% ============================================================
% Palette « Pop » OnBuch+ — mêmes hex que la classe P (plus_ui.dart)
% ============================================================
\definecolor{poporange}{HTML}{F59321}
\definecolor{popdark}{HTML}{E07A0C}
\definecolor{popcream}{HTML}{FFF6E8}
\definecolor{popink}{HTML}{1C1714}
\definecolor{poppurple}{HTML}{7A5AE0}
\definecolor{popgreen}{HTML}{2E9E6A}
\definecolor{poppink}{HTML}{E0457B}
\definecolor{popblue}{HTML}{2F7DE1}
\definecolor{popgold}{HTML}{C98A12}
\definecolor{popmuted}{HTML}{8A7F76}
\definecolor{popline}{HTML}{EADFD2}
\definecolor{popgray}{HTML}{8A7F76}
\colorlet{popgrey}{popgray}
% Alias tolérants (le modèle nomme parfois les couleurs différemment)
\colorlet{popwhite}{white}
\colorlet{popblack}{popink}
\colorlet{popred}{poppink}
\colorlet{popyellow}{popgold}
% Teintes claires (fonds de tableaux/encadrés) — le modèle nomme parfois
% les couleurs avec un suffixe L (ex. popblueL) sans les avoir définies.
\colorlet{poporangeL}{poporange!20}
\colorlet{popdarkL}{popdark!20}
\colorlet{poppurpleL}{poppurple!20}
\colorlet{popgreenL}{popgreen!20}
\colorlet{poppinkL}{poppink!20}
\colorlet{popblueL}{popblue!20}
\colorlet{popgoldL}{popgold!20}
\colorlet{popredL}{popred!20}
\colorlet{popgrayL}{popgray!20}
% Teintes claires pour fonds de tableaux / zones
\colorlet{popblueL}{popblue!10!white}
\colorlet{poporangeL}{poporange!14!white}
\colorlet{popgreenL}{popgreen!12!white}
\colorlet{poppurpleL}{poppurple!12!white}
\colorlet{poppinkL}{poppink!12!white}
\colorlet{popgoldL}{popgold!14!white}

\pagecolor{popcream}

% ============================================================
% Typographie
% ============================================================
\defaultfontfeatures{Ligatures=TeX}
\setmainfont{PlusJakartaSans-Medium.ttf}[Path=../../../fonts/,BoldFont=PlusJakartaSans-Bold.ttf]
% Maths en sans-serif (Fira Math) avec chiffres et lettres droites de Plus
% Jakarta, pour que les formules se fondent dans le texte.
\usepackage[math-style=ISO,bold-style=ISO]{unicode-math}
\setmathfont{FiraMath-Regular.otf}[Path=../../../fonts/]
\setmathfont{PlusJakartaSans-Medium.ttf}[Path=../../../fonts/,range={up/num,up/{Latin,latin},\mathup}]
\usepackage{icomma} % virgule décimale sans espace en mode math
% Caractères mathématiques Unicode absents des polices de texte (Plus Jakarta,
% Archivo Black) : rendus par leur équivalent mathématique, en texte comme en
% formule — sinon ils s'affichent en carré vide (ex. « Arithmétique dans ℤ »).
\usepackage{newunicodechar}
\newunicodechar{ℝ}{\ensuremath{\mathbb{R}}}
\newunicodechar{ℤ}{\ensuremath{\mathbb{Z}}}
\newunicodechar{ℂ}{\ensuremath{\mathbb{C}}}
\newunicodechar{ℕ}{\ensuremath{\mathbb{N}}}
\newunicodechar{ℚ}{\ensuremath{\mathbb{Q}}}
\newunicodechar{𝔻}{\ensuremath{\mathbb{D}}}
\newunicodechar{α}{\ensuremath{\alpha}}
\newunicodechar{β}{\ensuremath{\beta}}
\newunicodechar{γ}{\ensuremath{\gamma}}
\newunicodechar{δ}{\ensuremath{\delta}}
\newunicodechar{ε}{\ensuremath{\varepsilon}}
\newunicodechar{θ}{\ensuremath{\theta}}
\newunicodechar{λ}{\ensuremath{\lambda}}
\newunicodechar{μ}{\ensuremath{\mu}}
\newunicodechar{σ}{\ensuremath{\sigma}}
\newunicodechar{τ}{\ensuremath{\tau}}
\newunicodechar{φ}{\ensuremath{\varphi}}
\newunicodechar{ψ}{\ensuremath{\psi}}
\newunicodechar{ω}{\ensuremath{\omega}}
\newunicodechar{Σ}{\ensuremath{\Sigma}}
% majuscules produites par \MakeUppercase dans les titres (α-aminés -> Α)
\newunicodechar{Α}{\ensuremath{\alpha}}
\newunicodechar{Β}{\ensuremath{\beta}}
\newunicodechar{Γ}{\ensuremath{\Gamma}}
\newunicodechar{Δ}{\ensuremath{\Delta}}
\newunicodechar{Ε}{\ensuremath{\varepsilon}}
\newunicodechar{Θ}{\ensuremath{\Theta}}
\newunicodechar{Λ}{\ensuremath{\Lambda}}
\newunicodechar{Μ}{\ensuremath{\mu}}
\newunicodechar{Τ}{\ensuremath{\tau}}
\newunicodechar{Φ}{\ensuremath{\Phi}}
\newunicodechar{Ψ}{\ensuremath{\Psi}}
\newunicodechar{Ω}{\ensuremath{\Omega}}
\newunicodechar{↦}{\ensuremath{\mapsto}}
\newunicodechar{∈}{\ensuremath{\in}}
\newunicodechar{∉}{\ensuremath{\notin}}
\newunicodechar{∧}{\ensuremath{\wedge}}
\newunicodechar{⇔}{\ensuremath{\Leftrightarrow}}
\newunicodechar{⟺}{\ensuremath{\Longleftrightarrow}}
\newunicodechar{⇒}{\ensuremath{\Rightarrow}}
\newunicodechar{⟹}{\ensuremath{\Longrightarrow}}
\newunicodechar{∘}{\ensuremath{\circ}}
\newunicodechar{∪}{\ensuremath{\cup}}
\newunicodechar{∩}{\ensuremath{\cap}}
\newunicodechar{≡}{\ensuremath{\equiv}}
\newunicodechar{⊂}{\ensuremath{\subset}}
\newunicodechar{⊥}{\ensuremath{\perp}}
\newunicodechar{∃}{\ensuremath{\exists}}
\newunicodechar{∀}{\ensuremath{\forall}}
\newunicodechar{∛}{\ensuremath{\sqrt[3]{\,}}}
\newunicodechar{∜}{\ensuremath{\sqrt[4]{\,}}}
\newunicodechar{⃗}{\ensuremath{{}^{\to}}}
\newunicodechar{ⁿ}{\ifmmode^{n}\else\textsuperscript{\ensuremath{n}}\fi}
\newunicodechar{ˣ}{\ifmmode^{x}\else\textsuperscript{\ensuremath{x}}\fi}
\newunicodechar{ᵇ}{\ifmmode^{b}\else\textsuperscript{\ensuremath{b}}\fi}
\newunicodechar{ʳ}{\ifmmode^{r}\else\textsuperscript{\ensuremath{r}}\fi}
\newunicodechar{ᵘ}{\ifmmode^{u}\else\textsuperscript{\ensuremath{u}}\fi}
\newunicodechar{ʸ}{\ifmmode^{y}\else\textsuperscript{\ensuremath{y}}\fi}
\newunicodechar{ᵃ}{\ifmmode^{a}\else\textsuperscript{\ensuremath{a}}\fi}
\newunicodechar{ᵉ}{\ifmmode^{e}\else\textsuperscript{\ensuremath{e}}\fi}
\newunicodechar{ᵏ}{\ifmmode^{k}\else\textsuperscript{\ensuremath{k}}\fi}
\newunicodechar{ᵖ}{\ifmmode^{p}\else\textsuperscript{\ensuremath{p}}\fi}
\newunicodechar{ᴺ}{\ifmmode^{N}\else\textsuperscript{\ensuremath{N}}\fi}
\newunicodechar{ᶜ}{\ifmmode^{c}\else\textsuperscript{\ensuremath{c}}\fi}
\newunicodechar{ʰ}{\ifmmode^{h}\else\textsuperscript{\ensuremath{h}}\fi}
\newunicodechar{ᵠ}{\ifmmode^{q}\else\textsuperscript{\ensuremath{q}}\fi}
\newunicodechar{ₙ}{\ifmmode_{n}\else\textsubscript{\ensuremath{n}}\fi}
\newunicodechar{ₐ}{\ifmmode_{a}\else\textsubscript{\ensuremath{a}}\fi}
\newunicodechar{ᵢ}{\ifmmode_{i}\else\textsubscript{\ensuremath{i}}\fi}
\newunicodechar{ᵣ}{\ifmmode_{r}\else\textsubscript{\ensuremath{r}}\fi}
\newunicodechar{ₖ}{\ifmmode_{k}\else\textsubscript{\ensuremath{k}}\fi}
\newunicodechar{ᵦ}{\ifmmode_{\beta}\else\textsubscript{\ensuremath{\beta}}\fi}
\newunicodechar{ₓ}{\ifmmode_{x}\else\textsubscript{\ensuremath{x}}\fi}
\newfontfamily\popdisplay{ArchivoBlack-Regular.ttf}[Path=../../../fonts/]
\newfontfamily\poplogo{SpaceGrotesk-Bold.ttf}[Path=../../../fonts/]
\newfontfamily\popsemi{PlusJakartaSans-SemiBold.ttf}[Path=../../../fonts/]
\newfontfamily\popxbold{PlusJakartaSans-ExtraBold.ttf}[Path=../../../fonts/]
\newfontfamily\popxboldls{PlusJakartaSans-ExtraBold.ttf}[Path=../../../fonts/,LetterSpace=3.0]

\setlist[enumerate]{leftmargin=1.7em,itemsep=5pt,topsep=4pt}
\setlist[itemize]{leftmargin=1.7em,itemsep=4pt,topsep=4pt,label={\color{poporange}\textbullet}}
\setlength{\parindent}{0pt}
\setlength{\parskip}{5pt}
\renewcommand{\arraystretch}{1.25}

% pgfplots : style Pop par défaut
\pgfplotsset{
  every axis/.append style={axis line style={popink,line width=1pt},tick style={popink},
    grid style={popline,line width=0.5pt},label style={font=\small},tick label style={font=\footnotesize}},
  popcurve/.style={poporange,line width=1.8pt,no markers,smooth},
}
% Même style disponible dans un \draw TikZ simple (hors pgfplots)
\tikzset{popcurve/.style={poporange,line width=1.8pt}}
\tikzset{popgrid/.style={popline,very thin}}
\colorlet{popcurve}{poporange} % le nom du style est parfois utilisé comme couleur

% ============================================================
% Pill / squiggle
% ============================================================
\newcommand{\poppill}[3]{%
  \tikz[baseline=-0.55ex]\node[draw=popink,line width=1.2pt,rounded corners=2.6pt,%
    fill=#2,inner xsep=6.5pt,inner ysep=3pt]{%
    \popxboldls\fontsize{7}{7}\selectfont\color{#3}\MakeUppercase{#1}};%
}
\newcommand{\popsquiggle}[1]{%
  \tikz[baseline=-0.4ex]{
    \draw[#1,line width=2.6pt,line cap=round]
      plot [smooth] coordinates {(0,0.09) (0.34,0.01) (0.68,0.21) (1.02,0.01) (1.36,0.10)};
  }%
}
% Mot-clé surligné (marqueur orange sous le texte)
\newcommand{\cle}[1]{{\popxbold\color{popdark}#1}}

% ============================================================
% En-tête / pied
% ============================================================
\pagestyle{fancy}
\fancyhf{}
\renewcommand{\headrulewidth}{0pt}
\renewcommand{\footrulewidth}{0pt}
\lhead{\raisebox{-0.15ex}{\poplogo\fontsize{13}{13}\selectfont\color{popink}OnBuch\color{poporange}+}}
\rhead{\poppill{\DOCMATIERE\ \textbullet\ \DOCNIVEAU\ \textbullet\ Leçon \DOCLECON}{white}{popink}}
\lfoot{\popxbold\fontsize{7.6}{7.6}\selectfont\color{popmuted}\MakeUppercase{Cours signé OnBuch+}\ \textbullet\ {\popsemi\DOCTITRE}}
\rfoot{\tikz[baseline=-0.6ex]\node[rounded corners=8.5pt,fill=popink,minimum height=17pt,minimum width=17pt,inner xsep=5pt,inner sep=0]{\popxbold\fontsize{8}{8}\selectfont\color{popcream}\thepage\,/\,\pageref*{LastPage}};}

% ============================================================
% Titres
% ============================================================
\newcommand{\chaptitre}[2]{%
  \begin{tcolorbox}[enhanced,colback=poporange,colframe=popink,boxrule=2.6pt,arc=11pt,
      drop shadow=popink,left=15pt,right=15pt,top=12pt,bottom=11pt,before skip=3pt,after skip=18pt]
    \poppill{#2}{popink}{popcream}\par\vspace{9pt}
    {\raggedright\hyphenpenalty=10000\exhyphenpenalty=10000\popdisplay\fontsize{22}{24}\selectfont\color{popcream}\MakeUppercase{#1}\par}
  \end{tcolorbox}
}
% Section numérotée : pastille orange avec le numéro + titre Archivo
\newcounter{popsec}
\newcommand{\coursec}[1]{%
  \stepcounter{popsec}\setcounter{popsub}{0}%
  \par\vspace{16pt}\noindent
  \begin{minipage}{\linewidth}
  \tikz[baseline=-0.7ex]\node[draw=popink,line width=1.6pt,fill=poporange,rounded corners=4pt,minimum size=22pt,inner sep=0]{\popdisplay\fontsize{12}{12}\selectfont\color{popcream}\thepopsec};%
  \hspace{8pt}{\popdisplay\fontsize{15}{17}\selectfont\color{popink}\MakeUppercase{#1}}\par
  \vspace{4pt}\noindent{\color{poporange}\rule{36pt}{3.6pt}}
  \end{minipage}\par\nopagebreak\vspace{8pt}
}
\newcounter{popsub}
\newcommand{\courssub}[1]{%
  \stepcounter{popsub}%
  \par\vspace{9pt}\noindent{\popxbold\fontsize{11.5}{13}\selectfont\color{popink}{\color{poporange}\thepopsec.\thepopsub}\hspace{6pt}#1}\par\nopagebreak\vspace{4pt}
}

% ============================================================
% Boîtes pédagogiques — même recette : fond blanc, bordure épaisse colorée,
% ombre dure encre, badge plein. Titre optionnel [#1].
% ============================================================
\tcbset{popbase/.style={breakable,enhanced,colback=white,boxrule=2.3pt,arc=9pt,
  shadow={3pt}{-3pt}{0pt}{popink},left=14pt,right=14pt,top=11pt,bottom=11pt,before skip=11pt,after skip=15pt,
  fontupper=\popsemi\fontsize{10.6}{14.5}\selectfont\color{popink},
  fontlower=\popsemi\fontsize{10.6}{14.5}\selectfont\color{popink}}}
% Coche dessinée (le glyphe ✓ n'existe pas dans les polices du document)
\newcommand{\popcheck}[1]{\tikz[baseline=-0.5ex]\draw[#1,line width=1.6pt,line cap=round,line join=round](-0.13,0.0)--(-0.04,-0.09)--(0.13,0.1);}
\providecommand{\coche}{\popcheck{popgreen}}
\providecommand{\resultat}[1]{\par\smallskip\noindent\fcolorbox{popgreen}{popgreenL}{\parbox{\dimexpr\linewidth-2\fboxsep-2\fboxrule\relax}{\textbf{Résultat.} #1}}\par\smallskip}
\newcommand{\popboxhead}[3]{\poppill{#1}{#2}{white}\ifx\relax#3\relax\else\hspace{6pt}{\popxbold\fontsize{10.6}{12}\selectfont\color{popink}#3}\fi\par\vspace{7pt}}

\newtcolorbox{aretenir}[1][]{popbase,colframe=popgreen,before upper={\popboxhead{À retenir}{popgreen}{#1}}}
\newtcolorbox{exemplebox}[1][]{popbase,colframe=poporange,before upper={\popboxhead{Exemple}{poporange}{#1}}}
\newtcolorbox{definition}[1][]{popbase,colframe=popblue,before upper={\popboxhead{Définition}{popblue}{#1}}}
\newtcolorbox{propriete}[1][]{popbase,colframe=poppurple,before upper={\popboxhead{Propriété}{poppurple}{#1}}}
\newtcolorbox{methode}[1][]{popbase,colframe=popgold,before upper={\popboxhead{Méthode}{popgold}{#1}}}
\newtcolorbox{attention}[1][]{popbase,colframe=poppink,before upper={\popboxhead{Attention}{poppink}{#1}}}
\newtcolorbox{experience}[1][]{popbase,colframe=popblue,colback=popblueL,before upper={\popboxhead{Expérience}{popblue}{#1}}}
% « Le savais-tu ? » : carte violette pleine (contexte, culture, Cameroun)
\newtcolorbox{savaistu}[1][]{popbase,colback=poppurple,colframe=popink,
  fontupper=\popsemi\fontsize{10.4}{14.2}\selectfont\color{white},
  before upper={\poppill{Le savais-tu\,?}{white}{poppurple}\ifx\relax#1\relax\else\hspace{6pt}{\popxbold\color{white}#1}\fi\par\vspace{7pt}}}
% Exercice résolu : énoncé en haut, \tcblower puis la solution sur fond crème
\newcounter{popexo}\newcounter{popexr}
\newtcolorbox{exoresolu}[1][]{popbase,colframe=poporange,colbacklower=popcream,
  segmentation style={popink,line width=1.2pt,dashed},
  before upper={\stepcounter{popexr}\popboxhead{Exercice résolu \thepopexr}{poporange}{#1}},
  before lower={\poppill{Solution}{popink}{popcream}\par\vspace{6pt}}}
% Exercice (non résolu en place) — la correction est dans la partie Corrigés
\newtcolorbox{exercice}[1][]{popbase,colframe=popink,
  before upper={\stepcounter{popexo}\popboxhead{Exercice \thepopexo}{popink}{#1}}}
% Corrigé
\newtcolorbox{corrige}[1][]{popbase,colframe=popgreen,colback=popgreenL,
  before upper={\popboxhead{Corrigé}{popgreen}{#1}}}
% Objectifs / prérequis / bilan
\newtcolorbox{objectifs}{popbase,colframe=popink,colback=poporangeL,
  before upper={\popboxhead{Objectifs de la leçon}{popink}{}}}
\newtcolorbox{prerequis}{popbase,colframe=popink,colback=popblueL,
  before upper={\popboxhead{Prérequis}{popblue}{}}}
\newtcolorbox{bilan}{popbase,colframe=popink,colback=white,boxrule=2.8pt,
  before upper={\popboxhead{Fiche bilan}{poporange}{L'essentiel à connaître pour l'examen}}}
% Figure encadrée avec légende
\newtcolorbox{popfigure}[1][]{enhanced,breakable=false,colback=white,colframe=popline,boxrule=1.4pt,arc=8pt,
  left=10pt,right=10pt,top=10pt,bottom=8pt,before skip=10pt,after skip=12pt,halign=center,
  lower separated=false,halign lower=center,
  fontlower=\popsemi\fontsize{9}{11.5}\selectfont\color{popmuted},
  #1}
\newcommand{\legende}[1]{\par\vspace{6pt}{\popsemi\fontsize{9}{11.5}\selectfont\color{popmuted}#1}}

% ============================================================
% Signature OnBuch+
% ============================================================
\newcommand{\poplogoinline}[1]{{\poplogo\fontsize{#1}{#1}\selectfont\color{popink}OnBuch\color{poporange}+}}
\newcommand{\signatureonbuch}{%
  \begin{tcolorbox}[enhanced,colback=popink,colframe=popink,boxrule=2.2pt,arc=10pt,
      drop shadow=poporange,left=14pt,right=14pt,top=10pt,bottom=10pt,before skip=0pt,after skip=0pt]
    \tikz[baseline=-0.7ex]\node[circle,fill=popgreen,draw=popcream,line width=1.3pt,minimum size=22pt,inner sep=0]{\popcheck{white}};%
    \hspace{9pt}\begin{minipage}[c]{0.62\linewidth}
      {\popxbold\fontsize{12}{13}\selectfont\color{popcream}Signé\ \textbullet\ {\poplogo OnBuch\color{poporange}+}}\par\vspace{2pt}
      {\raggedright\popsemi\fontsize{8}{10.5}\selectfont\color{popline}\MakeUppercase{Équipe pédagogique OnBuch+}\\\MakeUppercase{Conforme au programme officiel MINESEC}\par}
    \end{minipage}\hfill
    {\poplogo\fontsize{22}{22}\selectfont\color{popcream}OnBuch\color{poporange}+}
  \end{tcolorbox}%
}

% ============================================================
% Page de garde
% #1 titre · #2 matière · #3 niveau · #4 module · #5 n° leçon · #6 durée ·
% #7 description / objectifs (texte ou itemize)
% ============================================================
\newcommand{\pagegarde}[7]{%
  \thispagestyle{empty}
  \newgeometry{a4paper,margin=1.7cm,top=1.6cm,bottom=1.4cm}%
  \begin{tikzpicture}[remember picture,overlay]
    \fill[poporange!18] ([xshift=-3.2cm,yshift=-3.6cm]current page.north east) circle (4.6cm);
    \fill[poppurple!14] ([xshift=2.4cm,yshift=7.6cm]current page.south west) circle (3.8cm);
  \end{tikzpicture}%
  {\poplogo\fontsize{30}{30}\selectfont\color{popink}OnBuch\color{poporange}+}\hfill
  \raisebox{0.6ex}{\poppill{Cours signé \textbullet\ Premium}{popink}{popcream}}\par
  \vspace{1pt}\popsquiggle{poppurple}\par
  \vspace{8pt}
  \begin{tcolorbox}[enhanced,colback=poporange,colframe=popink,boxrule=2.8pt,arc=13pt,
      drop shadow=popink,left=18pt,right=18pt,top=12pt,bottom=13pt,after skip=10pt]
    \poppill{#2 \textbullet\ #3}{popink}{popcream}\hspace{5pt}\poppill{#4}{white}{popink}\par\vspace{8pt}
    {\popdisplay\fontsize{14}{14}\selectfont\color{popink}LEÇON #5}\par\vspace{4pt}
    {\raggedright\hyphenpenalty=10000\exhyphenpenalty=10000\popdisplay\fontsize{25}{27}\selectfont\color{popcream}\MakeUppercase{#1}\par}
    \vspace{8pt}
    {\popxbold\fontsize{9.5}{9.5}\selectfont\color{popink}\MakeUppercase{Durée conseillée : #6}}
  \end{tcolorbox}
  \begin{tcolorbox}[enhanced,colback=white,colframe=popink,boxrule=2.2pt,arc=10pt,
      drop shadow=popink,left=16pt,right=16pt,top=10pt,bottom=10pt,after skip=8pt]
    \poppill{Dans cette leçon}{poppurple}{white}\par\vspace{5pt}
    \popsemi\fontsize{9.3}{12.6}\selectfont\color{popink}#7
  \end{tcolorbox}
  \noindent\begin{minipage}[t]{0.487\linewidth}
    \begin{tcolorbox}[enhanced,colback=popgreen,colframe=popink,boxrule=2.2pt,arc=10pt,
        drop shadow=popink,left=11pt,right=11pt,top=10pt,bottom=10pt,height=3.1cm,valign=center]
      \tcbox[colback=white,colframe=popink,boxrule=1.3pt,arc=5pt,boxsep=0pt,nobeforeafter,
        left=3pt,right=3pt,top=3pt,bottom=3pt,valign=center]{\qrcode[height=1.9cm]{https://play.google.com/store/apps/details?id=com.luvvix.onbuch&hl=fr}}%
      \hspace{8pt}\begin{minipage}[c]{0.5\linewidth}
        {\popdisplay\fontsize{11}{12.5}\selectfont\color{white}\MakeUppercase{Télécharge OnBuch}}\par\vspace{4pt}
        {\raggedright\popsemi\fontsize{8.4}{11}\selectfont\color{white}Gratuit \textbullet\ Google Play\\Tuteur IA, annales, concours, résultats\par}
      \end{minipage}
    \end{tcolorbox}
  \end{minipage}\hfill
  \begin{minipage}[t]{0.487\linewidth}
    \begin{tcolorbox}[enhanced,colback=poppurple,colframe=popink,boxrule=2.2pt,arc=10pt,
        drop shadow=popink,left=11pt,right=11pt,top=10pt,bottom=10pt,height=3.1cm,valign=center]
      \tikz[baseline=-0.7ex]\node[circle,fill=white,draw=popink,line width=1.3pt,minimum size=1.6cm,inner sep=0]{\popdisplay\fontsize{24}{24}\selectfont\color{poppurple}+};%
      \hspace{8pt}\begin{minipage}[c]{0.6\linewidth}
        {\popdisplay\fontsize{11}{12.5}\selectfont\color{white}\MakeUppercase{Passe à OnBuch+}}\par\vspace{4pt}
        {\raggedright\popsemi\fontsize{8.4}{11}\selectfont\color{white}Tous les cours signés, Tuteur IA illimité, fiches PDF — onglet OnBuch+ de l'app\par}
      \end{minipage}
    \end{tcolorbox}
  \end{minipage}
  \vspace{8pt}
  \popcontenu
  \vfill
  \signatureonbuch
  \restoregeometry
  \newpage
}

% Rangée « Ce cours contient » (page de garde)
\newcommand{\poptile}[3]{% #1 couleur #2 gros texte #3 légende
  \begin{tcolorbox}[enhanced,colback=#1,colframe=popink,boxrule=2pt,arc=9pt,shadow={2.5pt}{-2.5pt}{0pt}{popink},
      left=6pt,right=6pt,top=9pt,bottom=9pt,halign=center,valign=center,height=1.75cm,width=\linewidth,nobeforeafter]
    {\popdisplay\fontsize{12.5}{13}\selectfont\color{white}\MakeUppercase{#2}}\par\vspace{3pt}
    {\centering\popsemi\fontsize{7.8}{9.5}\selectfont\color{white}#3\par}
  \end{tcolorbox}}
\newcommand{\popcontenu}{%
  \par\noindent{\popxboldls\fontsize{7.5}{7.5}\selectfont\color{popmuted}CE COURS SIGNÉ CONTIENT}\par\vspace{7pt}
  \noindent\begin{minipage}[t]{0.235\linewidth}\poptile{popblue}{Cours}{complet, illustré, pas à pas}\end{minipage}\hfill
  \begin{minipage}[t]{0.235\linewidth}\poptile{poporange}{Méthodes}{et exercices résolus}\end{minipage}\hfill
  \begin{minipage}[t]{0.235\linewidth}\poptile{poppink}{Exercices}{type examen + corrigés}\end{minipage}\hfill
  \begin{minipage}[t]{0.235\linewidth}\poptile{popgreen}{Fiche bilan}{l'essentiel à retenir}\end{minipage}\par}

% Sommaire
\newtcolorbox{sommaire}{enhanced,colback=white,colframe=popink,boxrule=2.2pt,arc=10pt,
  drop shadow=popink,left=18pt,right=18pt,top=13pt,bottom=13pt,before skip=6pt,after skip=20pt,
  before upper={\poppill{Sommaire}{poppurple}{white}\par\vspace{9pt}\popsemi\fontsize{10.6}{14.5}\selectfont\color{popink}}}

% Signature de fin de cours — poussée tout en bas de la dernière page
\newcommand{\findecours}{%
  \par\vfill\nopagebreak
  \begin{center}\popsquiggle{poporange}\end{center}\vspace{4pt}
  \signatureonbuch
}
\AtBeginDocument{\def\llbracket{\mathopen{[\mkern-3.5mu[}}\def\rrbracket{\mathclose{]\mkern-3.5mu]}}} % intervalles d'entiers (glyphe ⟦ absent de la police)
% Glyphes absents de Fira Math : redéfinis à partir de glyphes disponibles
\AtBeginDocument{%
  \def\setminus{\mathbin{\backslash}}\let\smallsetminus\setminus
  \def\vdots{\mathord{\vbox{\baselineskip3.2pt\lineskiplimit0pt\kern4pt\hbox{.}\hbox{.}\hbox{.}}}}%
  \def\ddots{\mathinner{\mkern1mu\raise6.5pt\hbox{.}\mkern2mu\raise3.6pt\hbox{.}\mkern2mu\raise0.7pt\hbox{.}\mkern1mu}}%
  \def\triangle{\mathord{\scalebox{0.9}{$\symup{\Delta}$}}}%
  \def\nabla{\mathord{\raisebox{\depth}{\scalebox{1}[-1]{$\symup{\Delta}$}}}}%
  \def\checkmark{\popcheck{popdark}}%
  \def\star{\mathbin{\ast}}\def\bigstar{\mathord{\ast}}%
  \def\diamond{\mathbin{\scalebox{0.6}{$\Diamond$}}}%
  \def\bigcup{\mathop{\mathchoice{\raisebox{-0.3em}{\scalebox{1.7}{$\cup$}}}{\scalebox{1.25}{$\cup$}}{\cup}{\cup}}\displaylimits}%
  \def\bigcap{\mathop{\mathchoice{\raisebox{-0.3em}{\scalebox{1.7}{$\cap$}}}{\scalebox{1.25}{$\cap$}}{\cap}{\cap}}\displaylimits}%
  \def\lesseqqgtr{\mathrel{\lessgtr}}\def\bigoplus{\mathop{\oplus}\displaylimits}%
}
\providecommand{\ssi}{si et seulement si} % abréviation fréquente
\AtBeginDocument{\providecommand{\wideparen}{\overparen}} % arc de cercle

%% FIGFIX-BEGIN v4 — correctifs globaux des figures (docs/onbuch-generation/tools/figfix)
\usepackage{adjustbox}\usepackage{etoolbox}
% toute figure TikZ/pgfplots plus large que la ligne est réduite à la largeur disponible
\BeforeBeginEnvironment{tikzpicture}{\begin{adjustbox}{max width=\linewidth}}
\AfterEndEnvironment{tikzpicture}{\end{adjustbox}}
% légendes pgfplots lisibles par-dessus les courbes
\pgfplotsset{every axis/.append style={legend style={fill=white,fill opacity=0.92,text opacity=1,draw=black!15,inner sep=3pt}}}
% étiquettes d'arêtes (syntaxe edge["..."]) avec fond blanc
\tikzset{every edge quotes/.append style={fill=white,inner sep=1.2pt}}
%% FIGFIX-END
\begin{document}

\pagegarde{Microorganismes dans notre environnement}{SVTEEHB}{3ème}{SVTEEHB – classe de 3ème}{N04}{5 séances (fiche de progression harmonisée)}{Cette leçon fait découvrir le monde invisible des microorganismes (bactéries, virus, champignons microscopiques, organismes unicellulaires), leurs modes de vie (nutrition, respiration, reproduction) et les voies par lesquelles ils contaminent l'organisme humain. Elle s'appuie sur l'observation au microscope optique et sur l'interprétation de résultats expérimentaux pour construire des démarches scientifiques applicables à la vie courante.
\vspace{4pt}
\begin{multicols}{2}\small
\begin{itemize}
\item À la fin de cette leçon, je sais définir ce qu'est un microorganisme et citer ses différents groupes.
\item À la fin de cette leçon, je sais décrire les caractéristiques des bactéries, des virus, des champignons microscopiques et des organismes unicellulaires.
\item À la fin de cette leçon, je sais observer des microbes au microscope optique et réaliser un schéma d'observation légendé.
\item À la fin de cette leçon, je sais expliquer les modes de nutrition et de respiration des microbes à partir de résultats expérimentaux.
\item À la fin de cette leçon, je sais décrire les modes de reproduction des microorganismes et interpréter une courbe de croissance bactérienne.
\item À la fin de cette leçon, je sais identifier les différentes voies de pénétration des microbes dans l'organisme.
\end{itemize}\end{multicols}}

\begin{sommaire}
\begin{multicols}{2}
\begin{enumerate}
\item Découverte des microorganismes
\item Les différents groupes de microorganismes
\item Observation des microbes au microscope optique (TP N°5)
\item Nutrition et respiration des microbes
\item Reproduction des microbes
\item La contamination par les microorganismes
\item Activité d'intégration
\item Méthodes et exercices résolus
\item Exercices
\item Corrigés des exercices
\item Fiche bilan
\end{enumerate}
\end{multicols}
\end{sommaire}

\begin{objectifs}
\begin{itemize}
\item À la fin de cette leçon, je sais définir ce qu'est un microorganisme et citer ses différents groupes.
\item À la fin de cette leçon, je sais décrire les caractéristiques des bactéries, des virus, des champignons microscopiques et des organismes unicellulaires.
\item À la fin de cette leçon, je sais observer des microbes au microscope optique et réaliser un schéma d'observation légendé.
\item À la fin de cette leçon, je sais expliquer les modes de nutrition et de respiration des microbes à partir de résultats expérimentaux.
\item À la fin de cette leçon, je sais décrire les modes de reproduction des microorganismes et interpréter une courbe de croissance bactérienne.
\item À la fin de cette leçon, je sais identifier les différentes voies de pénétration des microbes dans l'organisme.
\item À la fin de cette leçon, je sais appliquer ces notions à des situations concrètes d'hygiène et de prévention de la vie courante.
\item À la fin de cette leçon, je sais rédiger et justifier une démarche, une réponse ou une conclusion scientifique.
\end{itemize}
\end{objectifs}

\begin{prerequis}
\begin{itemize}
\item La cellule : structure générale (membrane, cytoplasme, noyau) vue en 4ème
\item Notion d'être vivant unicellulaire et pluricellulaire (classification des êtres vivants)
\item Utilisation de base du microscope optique (grossissement, préparation d'une lame)
\item Notions d'hygiène du corps et des aliments vues en éducation à la santé
\item Notion de maladie et d'agent pathogène (début d'année de 3ème)
\end{itemize}
\end{prerequis}

\newpage

\coursec{Découverte des microorganismes}

\courssub{Situation d'accroche et problématique}

Imagine un instant : nous sommes à \textbf{Yaoundé}, un midi de semaine. Un groupe d'élèves achète du \emph{soya} (brochettes de viande grillée) chez un vendeur ambulant près du lycée. Quelques heures plus tard, plusieurs d'entre eux sont pris de violents vomissements et de diarrhées. Au même moment, à \textbf{Ngaoundéré}, une épidémie de grippe se propage rapidement dans les dortoirs d'un lycée : toux, fièvre, fatigue. Et au marché central, des sacs d'arachides moisies, recouverts d'un duvet verdâtre, sont saisis et détruits par les services d'hygiène.

Trois situations, un point commun : des \textbf{êtres invisibles à l'œil nu} sont à l'œuvre. Comment de si petits organismes peuvent-ils rendre malade, gâter nos aliments et se transmettre d'une personne à l'autre ? Quels sont ces \textbf{microorganismes} ? Comment vivent-ils, se reproduisent-ils et pénètrent-ils dans notre organisme ? C'est à ces questions que cette leçon va répondre, en commençant par découvrir ce que sont les microbes et comment on a prouvé leur existence.

\courssub{Qu'est-ce qu'un microorganisme ?}

\begin{definition}[Microorganisme (ou microbe)]
Un \cle{microorganisme} (ou \cle{microbe}) est un \textbf{être vivant} de très petite taille, \textbf{invisible à l'œil nu} et \textbf{observable uniquement au microscope} (optique ou électronique selon sa taille).
\end{definition}

Contrairement à ce que l'on pourrait croire au premier abord, un microbe n'est pas une « saleté » ou une « impureté » : c'est un \textbf{organisme vivant} à part entière. Il naît, se nourrit, respire, se reproduit et meurt. Il possède un matériel génétique (ADN ou ARN), un métabolisme et une structure cellulaire (pour les bactéries, champignons, protistes) ou acellulaire (pour les virus).

\begin{exemplebox}[Ordres de grandeur : la taille des microbes]
Pour se représenter leur petitesse, comparons :
\begin{itemize}
    \item Un \textbf{virus} (ex. grippe, VIH) : environ \SI{100}{nm} (\num{0,1} µm).
    \item Une \textbf{bactérie} (ex. \emph{Escherichia coli}, \emph{Salmonella}) : environ \SI{1}{à 5}{µm}.
    \item Une \textbf{cellule humaine} (ex. globule rouge) : environ \SI{7}{à 8}{µm} de diamètre.
    \item Un \textbf{grain de sable fin} : environ \SI{100}{µm} (\num{0,1} mm).
\end{itemize}
Il faut aligner \textbf{environ 1000 bactéries} bout à bout pour faire \SI{1}{mm} !
\end{exemplebox}

\begin{popfigure}
\begin{tikzpicture}[scale=0.9, every node/.style={font=\small}]
    % Échelle logarithmique horizontale : x = log10(taille en nm) * 2cm environ
    % 1 nm -> 0, 10 nm -> 2, 100 nm -> 4, 1 µm (1000 nm) -> 6, 10 µm -> 8, 100 µm -> 10, 1 mm -> 12
    \draw[->, thick, popdark] (0,0) -- (13,0) node[right] {Taille (échelle logarithmique)};
    % Graduations principales
    \foreach \val/\label in {0/1\,nm, 2/10\,nm, 4/100\,nm, 6/1\,µm, 8/10\,µm, 10/100\,µm, 12/1\,mm} {
        \draw (\val,-0.15) -- (\val,0.15);
        \node[below, align=center] at (\val,-0.5) {\label};
    }
    % Barres de taille représentatives (hauteur arbitraire pour visibilité)
    % Virus ~100 nm (x=4)
    \draw[popblue, very thick] (4, 0.5) -- (4, 1.5) node[above, popblue] {Virus (\approx 100\,nm)};
    % Bactérie ~1-5 µm (x=6 à 6.7)
    \draw[popgreen, very thick] (6, 0.5) -- (6, 2.0) node[above, popgreen] {Bactérie (\approx 1-5\,µm)};
    % Cellule humaine ~10 µm (x=8)
    \draw[poporange, very thick] (8, 0.5) -- (8, 2.5) node[above, poporange] {Cellule humaine (\approx 10\,µm)};
    % Grain de sable ~100 µm (x=10)
    \draw[poppurple, very thick] (10, 0.5) -- (10, 3.0) node[above, poppurple] {Grain de sable (\approx 100\,µm)};
    % Représentation schématique des objets (cercles de taille relative non proportionnelle à l'échelle log, mais ordre de grandeur)
    % On place des icônes au-dessus des barres
    \node[circle, draw=popblue, thick, minimum size=0.6cm, fill=popblue!20] at (4, 2.2) {V};
    \node[ellipse, draw=popgreen, thick, minimum width=0.8cm, minimum height=0.4cm, fill=popgreen!20] at (6, 2.7) {B};
    \node[circle, draw=poporange, thick, minimum size=1cm, fill=poporange!20] at (8, 3.3) {C};
    \node[circle, draw=poppurple, thick, minimum size=1.5cm, fill=poppurple!20] at (10, 4.0) {S};
    \node[align=center, font=\tiny] at (4, 1.8) {ADN/ARN +\\
capsïde};
    \node[align=center, font=\tiny] at (6, 2.3) {Cellule\\ procaryote};
    \node[align=center, font=\tiny] at (8, 2.9) {Cellule\\ eucaryote};
    \node[align=center, font=\tiny] at (10, 3.6) {Minéral};
\end{tikzpicture}
\legende{Comparaison des ordres de grandeur (échelle logarithmique) : virus, bactérie, cellule humaine et grain de sable. Notez l'écart immense entre le virus et le grain de sable (facteur 10 millions).}
\end{popfigure}

\begin{aretenir}[Exemple chiffré]
Pour égaler \SI{1}{mm}, il faut aligner environ \textbf{1000 bactéries} (de \SI{1}{µm}) bout à bout. Un virus est \textbf{10 fois plus petit} qu'une bactérie. Une cellule humaine est \textbf{10 à 100 fois plus grande} qu'une bactérie.
\end{aretenir}

\courssub{La fin de la génération spontanée : l'expérience décisive de Pasteur}

Jusqu'au milieu du \textsc{xixe} siècle, beaucoup de savants croyaient à la \textbf{génération spontanée} : l'idée que la vie pouvait apparaître « toute seule » à partir de matière inerte (viande pourrie $\to$ asticots, bouillon nutritif $\to$ microbes). \textbf{Louis Pasteur}, chimiste et biologiste français, a mis fin à ce débat par une expérience d'une élégance redoutable en \textbf{1862}.

\begin{experience}[Expérience de Pasteur : les flacons à col de cygne (1862)]
\textbf{Matériel :} Flacons à long col recourbé (forme de col de cygne), bouillon nutritif (eau + levure + sucre), source de chaleur.
\textbf{Protocole :}
\begin{enumerate}
    \item Remplir deux flacons identiques de bouillon nutritif.
    \item Faire bouillir le bouillon dans \textbf{les deux} flasks assez longtemps pour \textbf{stériliser} le liquide et l'air contenu dans le col (tuer tous les microbes éventuels).
    \item Laisser refroidir à température ambiante.
    \item \textbf{Flacon Témoin (A) :} Laisser le col intact (en col de cygne).
    \item \textbf{Flacon Expérimental (B) :} Casser le col du flacon (ou l'incliner pour que le bouillon touche le col).
\end{enumerate}
\textbf{Observations (après quelques jours) :}
\begin{itemize}
    \item \textbf{Flacon A (col intact) :} Le bouillon reste \textbf{limpide, stérile}. Aucun microbe ne s'est développé.
    \item \textbf{Flacon B (col cassé) :} Le bouillon devient \textbf{trouble, odorant}. Des microbes (bactéries, levures) s'y sont multipliés.
\end{itemize}
\textbf{Interprétation :}
\begin{itemize}
    \item L'air entre dans les deux flacons (le col n'est pas bouché).
    \item Dans le flacon A, les \textbf{poussières et microbes présents dans l'air} sont \textbf{piégés dans la courbe du col} (par gravité) et n'atteignent pas le bouillon. Le bouillon reste stérile $\to$ \textbf{pas de génération spontanée}.
    \item Dans le flacon B, les microbes de l'air tombent directement dans le bouillon $\to$ \textbf{contamination}.
    \item \textbf{Conclusion :} Les microbes ne naissent pas de la matière inerte ; ils \textbf{proviennent de l'air} (et de l'environnement). \og \textbf{Omne vivum ex vivo} \fg (Tout vivant vient du vivant).
\end{itemize}
\end{experience}

\begin{popfigure}
\begin{tikzpicture}[scale=0.85, every node/.style={font=\small}]
    % Flacon A (Témoin - col de cygne intact)
    \begin{scope}[xshift=-5cm]
        % Col de cygne
        \draw[popdark, very thick] (0,0) -- (0,1) .. controls (0,2) and (-1,2.5) .. (-1,3) .. controls (-1,3.5) and (0,3.5) .. (0,4) -- (0,5);
        % Corps du flacon
        \draw[popdark, very thick] (-1.5,0) -- (-1.5,-3) .. controls (-1.5,-3.5) and (1.5,-3.5) .. (1.5,-3) -- (1.5,0);
        % Bouillon limpide
        \fill[popblue!20] (-1.5,-3) .. controls (-1.5,-3.5) and (1.5,-3.5) .. (1.5,-3) -- (1.5,0) -- (-1.5,0) -- cycle;
        % Poussières/microbes piégés dans le col
        \foreach \i in {1,...,12} {
            \pgfmathsetmacro{\x}{-0.5 + rand*0.8}
            \pgfmathsetmacro{\y}{1.5 + rand*1.5}
            \fill[popmuted] (\x,\y) circle (0.04);
        }
        \draw[->, poporange, thick] (0.5, 2.5) -- (1.5, 2.5) node[right, poporange] {Poussières/microbes piégés};
        \node[align=center, font=\bfseries] at (0, -4) {Flacon A (Témoin)\\Col intact};
        \node[align=center, popgreen] at (0, -4.7) {Bouillon \textbf{limpide}\\\textbf{Stérile}};
    \end{scope}
    % Flacon B (Expérimental - col cassé)
    \begin{scope}[xshift=5cm]
        % Col cassé (droit)
        \draw[popdark, very thick] (0,0) -- (0,4);
        % Corps du flacon
        \draw[popdark, very thick] (-1.5,0) -- (-1.5,-3) .. controls (-1.5,-3.5) and (1.5,-3.5) .. (1.5,-3) -- (1.5,0);
        % Bouillon trouble
        \fill[poporange!30] (-1.5,-3) .. controls (-1.5,-3.5) and (1.5,-3.5) .. (1.5,-3) -- (1.5,0) -- (-1.5,0) -- cycle;
        % Microbes dans le bouillon
        \foreach \i in {1,...,30} {
            \pgfmathsetmacro{\x}{rand*1.2}
            \pgfmathsetmacro{\y}{-2.8 + rand*2.5}
            \fill[poporange] (\x,\y) circle (0.05);
        }
        % Flèches entrée air/microbes
        \draw[->, popred, thick] (-2.5, 2) -- (-1.5, 2) node[left, popred] {Air + microbes};
        \draw[->, popred, thick] (2.5, 2) -- (1.5, 2);
        \node[align=center, font=\bfseries] at (0, -4) {Flacon B (Expérimental)\\Col cassé};
        \node[align=center, popred] at (0, -4.7) {Bouillon \textbf{trouble}\\\textbf{Contaminé}};
    \end{scope}
    % Légende commune
    \node[align=center] at (0, 5.5) {\textbf{Expérience de Pasteur (1862)} : Réfutation de la génération spontanée};
\end{tikzpicture}
\legende{Les flacons à col de cygne de Pasteur. À gauche (témoin) : le col intact piège les microbes de l'air, le bouillon reste stérile. À droite (expérimental) : le col cassé laisse entrer les microbes, le bouillon se trouble. Preuve que les microbes viennent de l'air, pas de la matière inerte.}
\end{popfigure}

\begin{savaistu}[Louis Pasteur et la naissance de la microbiologie]
En 1862, \textbf{Louis Pasteur} (1822--1895) présente ses résultats à l'Académie des sciences. Il démontre que \textbf{l'air est le véhicule des microbes} (poussières, gouttelettes). Cette expérience fonde la \textbf{microbiologie} et valide la \textbf{stérilisation par la chaleur} (pasteurisation). Elle ouvre la voie à l'asepsie en chirurgie (Lister) et à la compréhension des maladies infectieuses.
\end{savaistu}

\begin{methode}[Démarche scientifique : Témoin vs Expérimental]
Pour conclure qu'un microbe est responsable d'une contamination ou d'une transformation, on utilise une \textbf{démarche comparative} rigoureuse :
\begin{enumerate}
    \item Préparer \textbf{deux milieux identiques} (même bouillon, même volume, même température).
    \item \textbf{Stériliser les deux} (chauffer pour tuer tout microbe préexistant).
    \item \textbf{Exposer l'un} au facteur testé (air ambiant, main sale, eau de rivière...) $\to$ \textbf{Milieu expérimental}.
    \item \textbf{Protéger l'autre} de ce facteur (col de cygne, couvercle stérile, autoclave...) $\to$ \textbf{Milieu témoin}.
    \item \textbf{Incuber} dans les mêmes conditions (température, durée).
    \item \textbf{Comparer} : si seul le milieu expérimental se trouble (développement microbien), le facteur testé est la source de contamination.
\end{enumerate}
\emph{C'est la méthode de référence en microbiologie (recherche, industrie agroalimentaire, hôpital).}
\end{methode}

\courssub{Ubiquité des microorganismes : ils sont partout !}

L'expérience de Pasteur le montre : les microbes sont dans l'air. Mais pas seulement. Les microorganismes ont conquis \textbf{tous les milieux} de la Terre où l'eau liquide existe, même extrêmes. On parle d'\textbf{ubiquité}.

\begin{center}
\begin{tabularx}{\linewidth}{|l|X|X|}
\hline
\rowcolor{popblueL}
\textbf{Milieu} & \textbf{Exemples de microbes présents} & \textbf{Rôle / Conséquence} \\
\hline
\textbf{Air} (poussières, aérosols) & Spores de champignons (moisissures), bactéries, virus & Transmission maladies respiratoires (grippe, tuberculose), contamination aliments \\
\hline
\textbf{Eau} (rivières, puits, robinet, mer) & Bactéries (\emph{E. coli}, \emph{Vibrio cholerae}), virus, algues, protozoaires & Maladies hydriques (choléra, typhoïde), indicateurs de pollution fécale \\
\hline
\textbf{Sol} (humus, rhizosphère) & Bactéries décompositrices, \emph{Rhizobium} (fixation azote), champignons, actinomycètes & Fertilité des sols, décomposition matière organique, antibiotiques naturels \\
\hline
\textbf{Aliments} (lait, viande, fruits, farine) & Levures, bactéries lactiques, moisissures, \emph{Salmonella}, \emph{Clostridium} & Fermentations utiles (yaourt, fromage, vin, \emph{kwem}, \emph{foléré}) ou altérations / toxi-infections \\
\hline
\textbf{Corps humain} (peau, tube digestif, bouche, nez) & \textbf{Flore commensale} : \emph{Staphylococcus epidermidis}, \emph{E. coli}, \emph{Lactobacillus}, \emph{Candida} & Protection barrière (concurrence), digestion (vitamines K, B12), éducation immunité \\
\hline
\textbf{Milieux extrêmes} (sources chaudes, lacs salés, glace) & Archées (thermophiles, halophiles) & Modèles origine de la vie, enzymes industrielles (Taq polymérase) \\
\hline
\end{tabularx}
\end{center}

\begin{attention}[Piège fréquent : « Microbe = Maladie »]
\textbf{FAUX !} La grande majorité des microorganismes sont \textbf{inoffensifs}, voire \textbf{indispensables}.
\begin{itemize}
    \item Votre tube digestif héberge \textbf{\num{e14} bactéries} (100 000 milliards) : votre \textbf{microbiote intestinal}. Il pèse \SI{1}{à 2}{kg} ! Il digère les fibres, synthétise des vitamines, bloque les pathogènes.
    \item Seule une \textbf{minorité} (quelques centaines d'espèces sur des millions) sont \textbf{pathogènes} (agents de maladies).
    \item Ne confondez pas \og microbe \fg et \og germe pathogène \fg.
\end{itemize}
\end{attention}

\courssub{Microbes utiles et microbes nuisibles : une distinction vitale}

\begin{propriete}[Classification selon l'impact sur l'homme]
\begin{itemize}
    \item \textbf{Microorganismes \textbf{utiles} (ou bénéfiques) :} Interviennent dans les cycles biogéochimiques (C, N, P), la fabrication d'aliments (fermentations lactiques, alcooliques, acétiques), la production d'antibiotiques, d'enzymes, de vaccins, le traitement des eaux usées, la digestion (microbiote).
    \item \textbf{Microorganismes \textbf{inoffensifs} (commensaux, saprophytes) :} Vivent dans l'environnement ou sur l'homme sans causer de pathologie (ex. la plupart des bactéries du sol, \emph{Staphylococcus epidermidis} sur la peau).
    \item \textbf{Microorganismes \textbf{pathogènes} (nuisibles) :} Capables de provoquer une \textbf{maladie infectieuse} chez un hôte sensible (bactéries : \emph{Salmonella}, \emph{Mycobacterium tuberculosis} ; virus : VIH, grippe, rougeole ; champignons : \emph{Candida albicans}, \emph{Aspergillus} ; protozoaires : \emph{Plasmodium} paludisme).
\end{itemize}
\end{propriete}

\begin{exemplebox}[Exemples concrets au Cameroun]
\begin{itemize}
    \item \textbf{Utiles :} Fermentation du \textbf{manioc} pour le \emph{bobolo}, \emph{miondo}, \emph{kwem} (bactéries lactiques, levures) ; fermentation du \emph{foléré} (bissap) ; fabrication du fromage local (wara) ; compostage des déchets agricoles ; \emph{Rhizobium} dans les champs de soja/arachide (engrais naturel).
    \item \textbf{Nuisibles :} \emph{Salmonella} / \emph{Shigella} / \emph{E. coli} pathogènes dans le soya mal cuit ou l'eau de boisson non traitée (diarrhées) ; \emph{Vibrio cholerae} (choléra) ; virus de la grippe, rougeole, COVID-19 ; \emph{Aspergillus flavus} sur arachides/maïs mal séchés (aflatoxines cancérigènes) ; \emph{Plasmodium} (paludisme) transmis par moustique.
\end{itemize}
\end{exemplebox}

\begin{exoresolu}[Analyse d'une situation : contamination d'un jus de fruit]
\textbf{Énoncé :} Un élève prépare un jus d'ananas frais. Il le laisse dans un verre ouvert sur la table à \SI{28}{\celsius}. Au bout de 24 h, le jus a un goût alcoolisé, des bulles remontent à la surface et une fine pellicule blanchâtre flotte. L'élève pense que « le jus a tourné tout seul » (génération spontanée).
\begin{enumerate}
    \item Que s'est-il passé réellement ? Nommez les microorganismes probables.
    \item Proposez une expérience (type Pasteur) pour prouver que ces microbes viennent de l'air et non du jus lui-même.
\end{enumerate}
\tcblower
\textbf{Solution détaillée :}
\begin{enumerate}
    \item \textbf{Interprétation :} Ce n'est \textbf{pas} de la génération spontanée. Les levures (champignons microscopiques, ex. \emph{Saccharomyces}) et bactéries acétiques présentes dans l'air (poussières, surface du fruit, ustensiles) sont tombées dans le jus. Elles ont trouvé un milieu riche en sucres (glucose, fructose), à température idéale. Elles ont effectué une \textbf{fermentation alcoolique} (levures : sucre $\to$ alcool + CO$_2$ + énergie) puis une \textbf{fermentation acétique} (bactéries acétiques : alcool $\to$ acide acétique = vinaigre). La pellicule est un biofilm de bactéries acétiques (mère de vinaigre).
    \item \textbf{Expérience témoin :}
    \begin{itemize}
        \item Préparer deux volumes égaux de jus d'ananas stérilisé (ébullition \SI{10}{min}).
        \item Verser dans deux flacons stériles.
        \item \textbf{Flacon A (Témoin) :} Fermer hermétiquement (bouchon stérile / col de cygne).
        \item \textbf{Flacon B (Expérimental) :} Laisser ouvert à l'air ambiant.
        \item Incuber les deux à \SI{28}{\celsius} pendant 48 h.
    \end{itemize}
    \textbf{Résultat attendu :} Flacon A : jus limpide, pas de bulles, goût sucré inchangé. Flacon B : trouble, bulles, goût alcoolisé/vinaigré.
    \textbf{Conclusion :} Seule l'exposition à l'air a provoqué la transformation. Les microbes viennent de l'air.
\end{enumerate}
\end{exoresolu}

\courssub{Savoirs admis : ce que la science a solidement établi}

Au terme de cette découverte, retenons les \textbf{savoirs admis} (consensus scientifique mondial, enseignés comme acquis) qui fondent la microbiologie moderne :

\begin{aretenir}[Savoirs admis sur les microorganismes (niveau 3ème)]
\begin{enumerate}
    \item \textbf{Existence et nature :} Les microorganismes sont des \textbf{êtres vivants réels}, observables au microscope (optique pour bactéries/champignons/protistes ; électronique pour virus). Leur existence n'est pas une hypothèse, c'est un \textbf{fait expérimental} (microscopie, cultures sur gélose, PCR, séquençage).
    \item \textbf{Origine :} Ils ne naissent \textbf{pas par génération spontanée}. \og Omne vivum ex vivo \fg : tout vivant provient d'un vivant préexistant (Pasteur, 1862). Ils se dispersent dans l'environnement (air, eau, sol, vecteurs).
    \item \textbf{Rôle dans les maladies :} Le lien de causalité entre un microbe spécifique et une maladie précise est établi par les \textbf{postulats de Koch} (1884) : (1) Le microbe est présent chez tout malade ; (2) On peut l'isoler et le cultiver en pur ; (3) La culture pure inoculée à un hôte sain reproduit la maladie ; (4) On ré-isole le même microbe chez le nouvel hôte. \emph{Travaux de Pasteur (anthrax, rage) et Koch (tuberculose, choléra).}
    \item \textbf{Diversité :} Ils appartiennent à des groupes très différents : \textbf{Bactéries} (procaryotes), \textbf{Archées} (procaryotes extrêmophiles), \textbf{Champignons microscopiques} (levures, moisissures -- eucaryotes), \textbf{Protistes} (amibes, paramécies, plasmodies -- eucaryotes), \textbf{Virus} (acellulaires, parasites intracellulaires obligatoires).
    \item \textbf{Ubiquité :} Ils colonisent tous les milieux aqueux de la biosphère, y compris le corps humain (microbiotes).
\end{enumerate}
\end{aretenir}

\begin{exercice}[Entraînement : Vocabulaire et raisonnement]
\begin{enumerate}
    \item Définissez le terme « microorganisme » en utilisant les mots : vivant, microscope, œil nu.
    \item Dans l'expérience de Pasteur, quel est le rôle du col de cygne ? Pourquoi le bouillon du flacon témoin reste-t-il stérile ?
    \item Donnez trois exemples de milieux où l'on trouve des microbes au Cameroun, en précisant pour chacun un microbe utile et un microbe nuisible possible.
    \item Un élève affirme : « Les virus ne sont pas des êtres vivants car on ne les voit pas au microscope optique. » Qu'en pensez-vous ? Justifiez.
    \item Complétez : « La démarche scientifique pour prouver qu'une contamination vient de l'air consiste à comparer un milieu \_\_\_\_\_\_\_\_ et un milieu \_\_\_\_\_\_\_\_. »
\end{enumerate}
\end{exercice}

\begin{corrige}[Exercice n°1]
\begin{enumerate}
    \item Un microorganisme est un \textbf{être vivant} de taille microscopique, \textbf{invisible à l'œil nu} et \textbf{observable au microscope}.
    \item Le col de cygne \textbf{piège les poussières et microbes de l'air} par gravité dans sa courbe. L'air circule mais les particules contaminantes n'atteignent pas le bouillon. Le bouillon reste stérile car \textbf{aucun microbe n'a pu entrer} (pas de génération spontanée).
    \item Exemples : \textbf{Sol} : utile = \emph{Rhizobium} (fixation azote pour légumineuses) ; nuisible = \emph{Clostridium tetani} (tétanos). \textbf{Eau de boisson} : utile = bactéries de station d'épuration ; nuisible = \emph{Vibrio cholerae} (choléra). \textbf{Intestin} : utile = \emph{E. coli} commensal (vitamine K) ; nuisible = \emph{Salmonella typhi} (fièvre typhoïde).
    \item \textbf{Faux.} Les virus \textbf{sont des êtres vivants} (au sens large : ils ont un génome, évoluent, se reproduisent). Ils sont \textbf{acellulaires} et \textbf{parasites obligatoires} : ils ne se reproduisent qu'à l'intérieur d'une cellule hôte. Leur taille (\approx \SI{100}{nm}) impose le \textbf{microscope électronique} pour les voir, mais cela ne remet pas en cause leur statut d'entité biologique.
    \item « ... un milieu \textbf{témoin (stérile, protégé)} et un milieu \textbf{expérimental (exposé à l'air)}. »
\end{enumerate}
\end{corrige}

\coursec{Les différents groupes de microorganismes}

Dans la section précédente, nous avons découvert que notre environnement — l'air, l'eau, le sol, la nourriture, notre propre peau — est peuplé d'êtres vivants invisibles à l'œil nu : les \cle{microorganismes}. Mais tous ces microbes ne se ressemblent pas ! Certains sont de véritables cellules, d'autres n'en sont pas ; certains possèdent un noyau, d'autres non ; certains nous rendent malades, d'autres nous sont utiles. Comment s'y retrouver dans cette diversité ? Les scientifiques les ont classés en \cle{quatre grands groupes} : les \cle{bactéries}, les \cle{virus}, les \cle{champignons microscopiques} et les \cle{organismes unicellulaires} (protistes). Étudions-les un par un.

\courssub{Les bactéries : des cellules sans noyau}

\textbf{Une observation pour commencer.} Un élève laisse un verre de lait à l'air libre pendant deux jours : le lait tourne et prend une odeur aigre. Au microscope, on observe des milliards de petits bâtonnets en mouvement. Ce sont des \cle{bactéries}, les microorganismes les plus abondants de notre planète.

\begin{definition}[Bactérie]
Une \cle{bactérie} est un microorganisme \cle{unicellulaire} (formé d'une seule cellule) \cle{sans noyau véritable} : son matériel génétique (un chromosome circulaire d'\cle{ADN}) flotte librement dans le cytoplasme. On dit que c'est un organisme \cle{procaryote}.
\end{definition}

Une bactérie mesure en général de \SI{1}{\micro\meter} à \SI{5}{\micro\meter} (soit environ mille fois plus petite qu'un grain de sable). Sa cellule comprend :
\begin{itemize}
\item une \cle{paroi} rigide, composée de peptidoglycane, qui protège la cellule et lui donne sa forme ;
\item une \cle{membrane} plasmique qui contrôle les échanges avec le milieu extérieur ;
\item du \cle{cytoplasme}, liquide intérieur où se déroulent les réactions de la vie (métabolisme) ;
\item un filament d'\cle{ADN circulaire}, qui contient l'information génétique, non enfermé dans un noyau ;
\item parfois un ou plusieurs \cle{flagelles}, sortes de longs filaments qui battent et permettent à la bactérie de se déplacer.
\end{itemize}

\begin{popfigure}
\begin{center}
\begin{tikzpicture}[scale=1.0]
% Corps de la bactérie (bacille)
\draw[fill=popgreenL, draw=popgreen, line width=1.2pt, rounded corners=18pt] (0,0) ellipse (2.6 and 1.1);
% Membrane interne
\draw[draw=popgreen!70, rounded corners=14pt] (0,0) ellipse (2.35 and 0.9);
% ADN circulaire (chromosome)
\draw[draw=poppurple, line width=1.1pt] (-0.4,0.1) .. controls (-1.0,0.5) and (0.2,0.7) .. (0.5,0.2)
   .. controls (0.8,-0.3) and (-0.6,-0.5) .. (-1.0,-0.1)
   .. controls (-1.3,0.2) and (-0.1,0.4) .. (-0.4,0.1);
% Ribosomes (points)
\foreach \p in {(1.5,0.4),(1.8,-0.3),(-1.8,0.3),(-1.6,-0.4),(0.9,-0.5),(1.2,0.6)}
  \fill[poporange] \p circle (0.07);
% Flagelle
\draw[draw=popblue, line width=1.2pt, decorate, decoration={coil, segment length=6pt, amplitude=3pt}] (2.6,0) -- (4.6,0.5);
% Légendes fléchées
\draw[->, popdark] (-0.3,0.35) -- (-0.3,1.9) node[above, font=\small]{3. Chromosome (ADN circulaire)};
\draw[->, popdark] (1.5,0.47) -- (1.5,1.9) node[above, font=\small]{4. Cytoplasme (ribosomes)};
\draw[->, popdark] (-2.2,-0.9) -- (-2.2,-1.8) node[below, font=\small]{1. Paroi};
\draw[->, popdark] (0.6,-1.0) -- (0.6,-1.8) node[below, font=\small]{2. Membrane};
\draw[->, popdark] (4.2,0.45) -- (4.2,1.6) node[above, font=\small]{5. Flagelle};
\end{tikzpicture}
\end{center}
\legende{Schéma d'une bactérie de forme bacille : 1. paroi protectrice (peptidoglycane) ; 2. membrane plasmique ; 3. chromosome circulaire libre dans le cytoplasme (pas de noyau) ; 4. cytoplasme avec ribosomes ; 5. flagelle assurant le déplacement.}
\end{popfigure}

\textbf{Des formes variées.} On distingue les bactéries selon leur forme :
\begin{itemize}
\item les \cle{coques} : forme sphérique (exemple : le \emph{Staphylocoque}, responsable d'infections cutanées) ;
\item les \cle{bacilles} : forme de bâtonnet (exemples : \emph{Escherichia coli}, habitant naturel de notre intestin, et le \cle{bacille de Koch}, agent de la tuberculose) ;
\item les \cle{spirilles} : forme de spirale (exemple : le vibrion cholérique, \emph{Vibrio cholerae}, agent du choléra, qui sévit encore lors d'épidémies dans certaines régions du Cameroun).
\end{itemize}

\begin{attention}[Bactéries : amies ou ennemies ?]
Toutes les bactéries ne sont pas dangereuses ! La plupart sont inoffensives, et beaucoup nous sont utiles : \emph{Escherichia coli} participe à la digestion et à la synthèse de vitamines dans notre intestin, et certaines bactéries lactiques transforment le lait en yaourt. Seule une petite partie des bactéries est \cle{pathogène} (provoque des maladies : tuberculose, choléra, méningite, etc.).
\end{attention}

\courssub{Les virus : des microorganismes non cellulaires}

\textbf{Une énigme scientifique.} À la fin du XIX\textsuperscript{e} siècle, les savants filtrent le jus de plantes malades à travers des filtres en porcelaine qui retiennent toutes les bactéries... et pourtant le liquide filtré rend encore les plantes malades ! Ils découvrent ainsi des « microbes » encore plus petits que les bactéries : les \cle{virus}.

\begin{definition}[Virus]
Un \cle{virus} est un microorganisme \cle{non cellulaire}, \cle{parasite obligatoire}, formé d'un \cle{acide nucléique} (ADN \textbf{ou} ARN, son information génétique) enfermé dans une enveloppe protéique appelée \cle{capside}.
\end{definition}

Les virus sont infiniment plus petits que les bactéries (de \SI{20}{\nano\meter} à \SI{300}{\nano\meter} environ) : on ne peut les voir qu'au microscope électronique. Un virus n'a ni membrane, ni cytoplasme, ni organites, ni ribosomes : \textbf{ce n'est pas une cellule}. Il ne possède pas de métabolisme propre.

\begin{popfigure}
\begin{center}
\begin{tikzpicture}[scale=1.0]
% Capside (polygone régulier à 5 côtés = icosaèdre vu de face)
\draw[fill=popblueL, draw=popblue, line width=1.2pt]
  (90:1.5) -- (162:1.5) -- (234:1.5) -- (306:1.5) -- (18:1.5) -- cycle;
% Acide nucléique en spirale (représentation schématique)
\draw[draw=poppurple, line width=1.1pt, decorate, decoration={coil, segment length=5pt, amplitude=4pt}] (-0.7,-0.4) -- (0.7,0.5);
% Protéines de surface (spicules / glycoprotéines)
\foreach \a in {90,162,234,306,18}
  \draw[draw=popblue, line width=1pt] (\a:1.5) -- (\a:1.9);
\foreach \a in {90,162,234,306,18}
  \fill[popblue] (\a:1.9) circle (0.08);
% Légendes
\draw[->, popdark] (1.05,1.05) -- (2.6,1.7) node[right, font=\small]{1. Capside (protéines)};
\draw[->, popdark] (0.1,0.15) -- (2.6,-0.6) node[right, font=\small]{2. Acide nucléique (ADN ou ARN)};
\draw[->, popdark] (90:1.9) -- (-1.2,2.6) node[left, font=\small]{3. Protéines de surface (fixation)};
\end{tikzpicture}
\end{center}
\legende{Schéma d'un virus (structure icosaédrique simplifiée) : 1. capside, enveloppe protéique protectrice ; 2. acide nucléique (ADN ou ARN) contenant l'information génétique ; 3. protéines de surface (spicules) permettant la reconnaissance et la fixation sur la cellule hôte.}
\end{popfigure}

\textbf{Un parasite strict.} Un virus est incapable de se nourrir, de respirer ou de se multiplier tout seul. Pour se reproduire, il doit obligatoirement pénétrer dans une \cle{cellule hôte} et détourner le matériel de cette cellule (ribosomes, énergie, nucléotides) pour fabriquer de nouveaux virus. C'est pourquoi on l'appelle \cle{parasite obligatoire} (ou parasite strict).

Exemples de virus et de maladies associées :
\begin{itemize}
\item le \cle{virus de la grippe} (Influenza), qui provoque chaque année des épidémies saisonnières ;
\item le \cle{VIH} (virus de l'immunodéficience humaine), responsable du Sida ;
\item le \cle{virus Ebola}, qui a causé de graves épidémies en Afrique de l'Ouest et centrale ;
\item le virus de la poliomyélite, contre lequel les enfants camerounais sont vaccinés dans le cadre du PEV (Programme Élargi de Vaccination).
\end{itemize}

\begin{attention}[Erreur fréquente : « un virus est une petite cellule »]
FAUX ! Un virus \textbf{n'est pas une cellule} : il n'a ni membrane, ni cytoplasme, ni ribosomes. Il ne peut \textbf{ni se nourrir ni se reproduire sans une cellule hôte}. Conséquence médicale majeure : les antibiotiques, qui ciblent des structures bactériennes (paroi, ribosomes, ADN), sont \textbf{sans effet sur les virus}. C'est pourquoi on ne soigne pas la grippe, le rhume ou le COVID-19 avec des antibiotiques !
\end{attention}

\courssub{Les champignons microscopiques : levures et moisissures}

Qui n'a jamais vu apparaître un duvet blanc, vert ou noir sur un pain oublié, ou sur des arachides mal conservées ? Ce sont des \cle{moisissures}, des champignons microscopiques.

\begin{definition}[Champignon microscopique]
Un \cle{champignon microscopique} est un microorganisme dont les cellules possèdent un \cle{noyau véritable} délimité par une membrane nucléaire : ce sont des organismes \cle{eucaryotes}. On distingue les \cle{levures}, unicellulaires, et les \cle{moisissures}, formées de filaments microscopiques (hyphes).
\end{definition}

\begin{itemize}
\item Les \cle{levures} sont des champignons \cle{unicellulaires} de forme ovale ou sphérique. Elles se reproduisent principalement par \cle{bourgeonnement} : un petit bourgeon apparaît sur la cellule mère, grossit, reçoit une copie du noyau, puis se détache pour devenir une cellule fille. Exemple : la levure de boulanger (\emph{Saccharomyces cerevisiae}).
\item Les \cle{moisissures} forment un réseau de \cle{filaments} (hyphes) ramifiés qui s'étendent sur la nourriture (substrat) et forment un feutrage visible à l'œil nu (le mycélium). Exemples : la moisissure du pain (\emph{Rhizopus}), la moisissure de l'arachide (\emph{Aspergillus}), le \emph{Penicillium} (producteur de pénicilline, premier antibiotique découvert).
\end{itemize}

\begin{popfigure}
\begin{center}
\begin{tikzpicture}[scale=1.0]
% Cellule mère
\draw[fill=popgoldL, draw=popgold, line width=1.2pt] (0,0) ellipse (1.3 and 1.0);
% Noyau mère
\draw[fill=poppurpleL, draw=poppurple] (-0.2,0) circle (0.35);
% Mitochondries (symboles)
\foreach \p in {(-0.8,0.5),(-0.5,-0.6),(0.5,0.6)}
  \draw[poporange, line width=0.8pt] \p ++(-0.1,-0.05) -- ++(0.2,0.1) \p ++(0.1,-0.05) -- ++(-0.2,0.1);
% Bourgeon
\draw[fill=popgoldL, draw=popgold, line width=1.2pt] (1.7,0.9) circle (0.55);
% Noyau du bourgeon (en division/migration)
\draw[fill=poppurpleL, draw=poppurple] (1.6,0.9) circle (0.18);
% Flèche de croissance
\draw[->, popdark, thick] (1.0,0.9) -- (1.3,0.9);
% Légendes
\draw[->, popdark] (-0.2,0.35) -- (-0.2,1.7) node[above, font=\small]{2. Noyau (eucaryote)};
\draw[->, popdark] (-1.3,-0.6) -- (-1.3,-1.5) node[below, font=\small]{1. Cellule mère};
\draw[->, popdark] (1.7,1.45) -- (1.7,2.1) node[above, font=\small]{3. Bourgeon (cellule fille en formation)};
\end{tikzpicture}
\end{center}
\legende{Schéma d'une levure (\emph{Saccharomyces}) en bourgeonnement : 1. cellule mère ; 2. noyau véritable (organisme eucaryote) ; 3. bourgeon recevant une copie du noyau avant de se séparer.}
\end{popfigure}

\begin{savaistu}[Les levures de nos beignets]
Les levures utilisées pour faire lever la pâte de beignets (les fameux « puff-puff » vendus dans nos villes) sont des champignons microscopiques \textbf{vivants} ! En se nourrissant du sucre de la pâte (fermentation alcoolique), elles dégagent du gaz carbonique (\ce{CO2}) qui forme des bulles : la pâte gonfle et les beignets deviennent moelleux. Sans microbes, pas de bons beignets ! De même, le vin de palme et le bilbil (bière de mil) doivent leur fabrication à l'action de levures.
\end{savaistu}

\courssub{Les organismes unicellulaires : les protistes}

Il existe enfin des microorganismes formés d'\textbf{une seule cellule à noyau} (eucaryotes), qui ne sont ni des bactéries, ni des champignons, ni des animaux, ni des plantes : ce sont les \cle{protistes} (ou organismes unicellulaires eucaryotes). Ils vivent surtout dans les milieux humides (eaux stagnantes, sols humides, ou à l'intérieur d'autres êtres vivants en tant que parasites).

Exemples :
\begin{itemize}
\item l'\cle{amibe} (\emph{Amoeba proteus}) : cellule sans forme fixe, qui se déplace et capture sa nourriture (bactéries, petites particules) grâce à des prolongements changeants de son cytoplasme, les \cle{pseudopodes} (pieds faux) ; elle possède une vacuole contractile pour l'osmorégulation ;
\item la \cle{paramécie} (\emph{Paramecium}) : cellule en forme de sabot, recouverte de milliers de \cle{cils} vibratiles qui lui permettent de nager rapidement dans l'eau et de créer un courant pour amener la nourriture vers sa bouche (cytostome) ;
\item le \cle{Plasmodium} : protiste parasite intracellulaire, transmis à l'Homme par la piqûre d'un moustique vecteur, l'anophèle femelle.
\end{itemize}

\begin{popfigure}
\begin{center}
\begin{tikzpicture}[scale=1.0]
% Amibe : contour irrégulier (forme amiboïde)
\draw[fill=popgreenL, draw=popgreen, line width=1.2pt]
  plot[smooth cycle, tension=0.9] coordinates {(-2.2,0.3) (-1.4,1.1) (-0.3,0.9) (0.6,1.4) (1.5,0.7) (2.3,0.5) (1.8,-0.4) (2.5,-1.0) (1.0,-0.9) (0.2,-1.3) (-0.9,-0.8) (-1.9,-1.1) (-2.5,-0.4)};
% Noyau
\draw[fill=poppurpleL, draw=poppurple] (-0.2,0.1) circle (0.4);
% Vacuole contractile
\draw[fill=white, draw=popblue, line width=1pt] (1.1,0.3) circle (0.28);
\draw[popblue] (1.1,0.3) circle (0.35); % membrane
% Vacuole digestive
\draw[fill=poporangeL, draw=poporange] (-1.2,-0.3) circle (0.22);
% Légendes
\draw[->, popdark] (-0.2,0.5) -- (-0.2,1.9) node[above, font=\small]{1. Noyau};
\draw[->, popdark] (2.3,0.5) -- (3.2,1.4) node[right, font=\small]{2. Pseudopode (déplacement, phagocytose)};
\draw[->, popdark] (1.1,0.3) -- (3.2,-0.5) node[right, font=\small]{3. Vacuole contractile (osmorégulation)};
\draw[->, popdark] (-1.2,-0.3) -- (-3.0,-0.3) node[left, font=\small]{4. Vacuole digestive};
\draw[->, popdark] (-1.9,-1.0) -- (-1.9,-1.9) node[below, font=\small]{5. Cytoplasme (ectoplasme / endoplasme)};
\end{tikzpicture}
\end{center}
\legende{Schéma d'une amibe : 1. noyau (eucaryote) ; 2. pseudopodes, prolongements cytoplasmiques temporaires servant au déplacement et à la capture des proies (phagocytose) ; 3. vacuole contractile (régulation de l'eau) ; 4. vacuole digestive ; 5. cytoplasme.}
\end{popfigure}

\begin{exemplebox}[Le Plasmodium et le paludisme]
Le \emph{Plasmodium} (espèces \emph{falciparum}, \emph{malariae}, \emph{ovale}, \emph{vivax}) est un protiste parasite intracellulaire obligatoire, transmis par la piqûre de l'anophèle femelle infectée. Il est l'agent du \cle{paludisme} (malaria), qui reste la première cause de consultation et de mortalité infantile au Cameroun. Le parasite se multiplie d'abord dans le foie (phase hépatique silencieuse), puis dans les globules rouges (phase érythrocytaire), qu'il détruit cycliquement, provoquant les accès de fièvre caractéristiques, une anémie et des troubles neurologiques graves (paludisme cérébral). La prévention repose sur l'utilisation des moustiquaires imprégnées d'insecticide à longue durée d'action (MILD), la pulvérisation intra-domiciliaire, la chimioprévention du paludisme saisonnier (CPS) et la prise en charge rapide par les ACT (artémisinine-based combination therapies).
\end{exemplebox}

\courssub{Comment distinguer les quatre groupes ?}

Pour classer un microorganisme observé ou décrit, on utilise des \cle{critères de distinction} simples et logiques :
\begin{enumerate}
\item \textbf{Possède-t-il une structure cellulaire (membrane, cytoplasme) ?}
      \begin{itemize}
      \item Si \textbf{non} : c'est un \textbf{virus} (acide nucléique + capside).
      \item Si \textbf{oui} : c'est une cellule ; on passe au critère 2.
      \end{itemize}
\item \textbf{Sa cellule possède-t-elle un noyau véritable (membrane nucléaire) ?}
      \begin{itemize}
      \item Si \textbf{non} : c'est une \textbf{bactérie} (procaryote).
      \item Si \textbf{oui} : c'est un eucaryote ; on passe au critère 3.
      \end{itemize}
\item \textbf{Parmi les eucaryotes microscopiques}, on distingue :
      \begin{itemize}
      \item les \textbf{champignons microscopiques} : paroi en chitine (souvent), nutrition par absorption (saprotrophes ou parasites), reproduction par spores ou bourgeonnement (levures) ;
      \item les \textbf{organismes unicellulaires (protistes)} : pas de paroi en chitine (souvent péllicule), nutrition variée (phagocytose, photosynthèse, absorption), locomotion souvent par pseudopodes, cils ou flagelles.
      \end{itemize}
\end{enumerate}

\begin{popfigure}
\begin{center}
\begin{tikzpicture}[
  grow=down,
  level 1/.style={sibling distance=55mm, level distance=18mm},
  level 2/.style={sibling distance=28mm, level distance=16mm},
  every node/.style={draw, rounded corners=3pt, fill=popcream, font=\small, align=center, inner sep=3pt},
  edge from parent/.style={draw=popdark, thick, -},
  edge from parent path={(\tikzparentnode.south) -- ++(0,-4mm) -| (\tikzchildnode.north)}
]
\node[fill=popblueL, font=\small\bfseries, text width=3.5cm, align=center] {Microorganismes}
  child { node[fill=poporangeL, text width=2.5cm, align=center] {Virus\\ (non cellulaires\\ parasites obligatoires)} }
  child { node[fill=popgreenL, text width=2.5cm, align=center] {Bactéries\\ (procaryotes\\ sans noyau)} }
  child { node[fill=poppurpleL, font=\small\bfseries, text width=3.5cm, align=center] {Cellules à noyau\\ (eucaryotes)}
    child { node[fill=popgoldL, text width=2.5cm, align=center] {Champignons microscopiques\\ (levures, moisissures)\\ nutrition absorbante} }
    child { node[fill=poppinkL, text width=2.5cm, align=center] {Organismes unicellulaires\\ (protistes : amibes,\\ paramécies, Plasmodium)} }
  };
\end{tikzpicture}
\end{center}
\legende{Arbre de classification dichotomique des microorganismes en quatre groupes, basé sur la structure cellulaire et la présence d'un noyau.}
\end{popfigure}

\begin{center}
\rowcolors{1}{popcream}{white}
\begin{tabularx}{\linewidth}{|l|X|X|X|X|}
\hline
\rowcolor{popblueL} \textbf{Critère} & \textbf{Bactéries} & \textbf{Virus} & \textbf{Champignons microscopiques} & \textbf{Organismes unicellulaires (Protistes)} \\
\hline
Structure cellulaire & Cellule procaryote (sans noyau) & Non cellulaire (capside + acide nucléique) & Cellule(s) eucaryote(s) (noyau vrai) & Cellule eucaryote (noyau vrai) \\
\hline
Taille typique & \SI{1}{\micro\meter} à \SI{5}{\micro\meter} & \SI{20}{\nano\meter} à \SI{300}{\nano\meter} & \SI{5}{\micro\meter} à \SI{10}{\micro\meter} (levures) & \SI{10}{\micro\meter} à \SI{500}{\micro\meter} \\
\hline
Paroi & Oui (peptidoglycane) & Non (capside protéique) & Oui (souvent chitine) & Non (péllicule souple) ou test (silice) \\
\hline
Locomotion & Flagelles (parfois) & Aucune (transport passif) & Aucune (sporulation, croissance apicale) & Pseudopodes, cils, flagelles \\
\hline
Nutrition & Absorption (saprotrophe ou parasite) & Parasite obligatoire (détournement machinerie hôte) & Absorption (saprotrophe ou parasite) & Phagocytose, absorption, photosynthèse \\
\hline
Reproduction & Division binaire (scissiparité) & Réplication dans cellule hôte & Bourgeonnement (levures), spores (moisissures) & Division mitotique (scissiparité), sexuée parfois \\
\hline
Exemples courants & \emph{E. coli}, Bacille de Koch, Vibrion cholérique & Grippe, VIH, Ebola, Polio, Rougeole & Levure de boulanger, \emph{Penicillium}, \emph{Aspergillus} & Amibe, Paramécie, \emph{Plasmodium}, Euglène \\
\hline
Maladies humaines & Tuberculose, choléra, méningite, pneumonie & Grippe, Sida, Ebola, poliomyélite, rougeole, COVID-19 & Mycoses cutanées (teignes), candidoses, aspergillose & Paludisme, amibiase (dysenterie amibienne), toxoplasmose \\
\hline
Visibilité microscope optique & Oui (coloration souvent nécessaire) & Non (microscope électronique requis) & Oui & Oui \\
\hline
\end{tabularx}
\end{center}

\begin{exoresolu}[Identification d'un microorganisme observé]
Énoncé. Au laboratoire du centre de santé, on observe au microscope optique (objectif \SI{40}{\times}) un microorganisme formé d'une seule cellule possédant un noyau bien visible. Il se déplace activement grâce à des prolongements changeants de son corps (pseudopodes) et engloutit des bactéries.
\begin{enumerate}
\item À quel groupe de microorganismes appartient-il ? Justifiez votre réponse en citant \textbf{deux} critères d'observation.
\item Peut-il s'agir d'un virus ? Expliquez pourquoi.
\item Quel est le genre probable de cet organisme ? Citez une maladie qu'il peut provoquer s'il est pathogène.
\end{enumerate}
\tcblower
Solution détaillée.
\begin{enumerate}
\item Ce microorganisme appartient au groupe des \textbf{organismes unicellulaires (protistes)}.
      \textbf{Justification :}
      \begin{itemize}
      \item Critère 1 : Il possède une \textbf{structure cellulaire} complète ( cytoplasme, noyau, organites visibles), ce qui exclut les virus.
      \item Critère 2 : Sa cellule possède un \textbf{noyau véritable} (eucaryote), ce qui exclut les bactéries (procaryotes).
      \item Critère 3 (discriminant) : Il se déplace et se nourrit par \textbf{pseudopodes} (phagocytose), mode de vie caractéristique des amibes (protistes), distinct de la nutrition absorbante des champignons.
      \end{itemize}
\item \textbf{Non}, il ne peut pas s'agir d'un virus, pour deux raisons fondamentales :
      \begin{itemize}
      \item Un virus n'est \textbf{pas une cellule} : il ne possède ni cytoplasme, ni noyau, ni organites, ni membrane plasmique.
      \item Un virus est \textbf{invisible au microscope optique} (taille < \SI{300}{\nano\meter} < pouvoir de résolution \approx \SI{200}{\nano\meter}) ; seul le microscope électronique permet de l'observer.
      \end{itemize}
\item Le genre probable est \textbf{\emph{Amoeba}} (amibe). S'il s'agit de l'espèce \emph{Entamoeba histolytica} (amibe dysentérique), elle est responsable de l'\textbf{amibiase} (dysenterie amibienne), une diarrhée sanglante et muqueuse fréquente en zone tropicale.
\end{enumerate}
\end{exoresolu}

\begin{aretenir}[L'essentiel de la section]
\begin{itemize}
\item Les microorganismes se classent en \cle{quatre groupes} : bactéries, virus, champignons microscopiques et organismes unicellulaires (protistes).
\item Les \cle{bactéries} sont des cellules \cle{procaryotes} (sans noyau, ADN circulaire libre) : formes en coques, bacilles ou spirilles ; paroi en peptidoglycane.
\item Les \cle{virus} sont des entités \cle{non cellulaires} (acide nucléique ADN ou ARN + capside protéique) ; ce sont des \cle{parasites obligatoires} stricts, invisibles au microscope optique.
\item Les \cle{champignons microscopiques} (levures unicellulaires, moisissures filamenteuses) et les \cle{organismes unicellulaires} (amibes, paramécies, Plasmodium) sont des \cle{eucaryotes} (noyau vrai, organites).
\item La classification dichotomique repose sur : 1) structure cellulaire (oui/non) $\rightarrow$ 2) présence d'un noyau (procaryote/eucaryote) $\rightarrow$ 3) mode de nutrition/locomotion pour les eucaryotes.
\end{itemize}
\end{aretenir}

Maintenant que nous savons reconnaître les quatre groupes de microorganismes grâce à leurs caractéristiques cellulaires et morphologiques, il est temps de passer à la pratique : au laboratoire, nous allons apprendre à \textbf{observer les microbes au microscope optique} (TP n°5), afin de vérifier par nous-mêmes ce que nous venons d'étudier et d'identifier des bactéries, des levures et éventuellement des protistes sur lame préparée.

\coursec{Observation des microbes au microscope optique (TP N°5)}

Après avoir identifié les principaux groupes de microorganismes (bactéries, virus, champignons microscopiques, protistes) et leurs caractéristiques générales, nous devons maintenant les \textbf{observer réellement}. Les microorganismes étant invisibles à l'œil nu, le microscope optique est l'outil indispensable du biologiste pour explorer ce monde microscopique. Ce TP N°5 te permettra d'acquérir les gestes techniques rigoureux de la microscopie et de réaliser des schémas d'observation scientifiques conformes aux attendus du BEPC.

\courssub{Rappel des règles de sécurité et d'utilisation du microscope}

\begin{attention}[Sécurité au laboratoire — Règles impératives]
La manipulation de microorganismes, même non pathogènes en classe (levures, moisissures, bactéries de yaourt), impose une hygiène stricte pour éviter toute contamination croisée ou accidentelle :
\begin{itemize}
    \item \textbf{Port de la blouse} obligatoire (coton, manches longues) pour protéger tes vêtements et ta peau.
    \item \textbf{Interdiction absolue de porter les mains à la bouche, au nez ou aux yeux} pendant la manipulation (pas de grignotage, pas de stylo machouillé).
    \item \textbf{Pas de pipetage à la bouche} : utilise systématiquement une poire à pipeter ou une pipette automatique.
    \item \textbf{Lavage des mains} minutieux au savon (ou solution hydroalcoolique) \textbf{avant} et \textbf{après} la séance.
    \item Désinfection du plan de travail avant et après le TP (eau de Javel diluée ou alcool \SI{70}{\%}).
    \item Élimination des lames et lamelles usagées dans un bac à déchets contaminés (ou trempage dans un désinfectant) ; ne les jette pas à la poubelle classique.
\end{itemize}
\end{attention}

\begin{methode}[Prise en main du microscope optique — Étapes ordonnées]
Pour observer sans casser le matériel et obtenir une image nette, suis cette procédure dans l'ordre :
\begin{enumerate}
    \item \textbf{Transport} : porte le microscope d'une main sous la base (pied), l'autre tenant le bras (colonne). Pose-le délicatement sur une table stable, loin du bord.
    \item \textbf{Éclairage} : branche la lampe (ou oriente le miroir vers une source lumineuse diffuse, jamais le soleil direct). Ouvre le diaphragme (iris) au maximum pour commencer.
    \item \textbf{Choix de l'objectif} : place l'\textbf{objectif de plus faible grossissement} (souvent \textbf{$\times 4$} ou \textbf{$\times 10$}) en position de travail (clic audible). \textbf{Jamais l'objectif $\times 40$ ou $\times 100$ au départ}.
    \item \textbf{Préparation} : dépose la lame (échantillon côté couverture) sur la platine, cale-la sous les pinces. Centre la zone d'intérêt au-dessus du trou de la platine (diaphragme).
    \item \textbf{Mise au point grossière} : \textbf{regarde par le côté} (pas dans l'oculaire), abaisse l'objectif au maximum vers la lame avec la \textbf{vis macrométrique} (grosse vis) jusqu'à \SI{1}{\milli\meter}--\SI{2}{\milli\meter} de la lamelle.
    \item \textbf{Mise au point fine (observation)} : regarde dans l'oculaire. Tourne la \textbf{vis macrométrique} dans le sens de l'\textbf{éloignement} (objectif qui monte) jusqu'à voir une image. Affine avec la \textbf{vis micrométrique} (petite vis) pour la netteté parfaite.
    \item \textbf{Augmentation du grossissement} : centre la zone d'intérêt. Passe à l'objectif supérieur ($\times 10$ puis $\times 40$). \textbf{N'utilise plus la vis macrométrique} : seulement la vis micrométrique pour la netteté. Referme légèrement le diaphragme pour augmenter le contraste.
    \item \textbf{Rangement} : remet l'objectif $\times 4$ (ou $\times 10$) en place, baisse la platine au maximum, éteins la lumière, débranche, replace la housse.
\end{enumerate}
\end{methode}

\begin{popfigure}
\begin{tikzpicture}[scale=0.85, every node/.style={font=\small}]
    % Base / Pied
    \draw[popdark, thick, fill=popblueL!30] (-2.5,-3) rectangle (2.5,-2.5);
    % Colonne / Bras
    \draw[popdark, thick, fill=popblueL!20] (-0.8,-2.5) -- (-0.8,4) -- (0.8,4) -- (0.8,-2.5) -- cycle;
    % Platine
    \draw[popdark, thick, fill=popmuted!20] (-1.8,-0.5) rectangle (1.8,0);
    % Trou platine
    \draw[popdark, thick, fill=white] (-0.3,-0.4) rectangle (0.3,-0.1);
    % Pinces
    \draw[popdark, thick] (-1.8,-0.5) -- (-2.2,-0.8) (-1.8,0) -- (-2.2,0.3);
    \draw[popdark, thick] (1.8,-0.5) -- (2.2,-0.8) (1.8,0) -- (2.2,0.3);
    % Vis macrométrique (gauche)
    \draw[popdark, thick] (-1.5,-2.5) -- (-1.5,-3.5) node[below] {Vis macrométrique};
    \draw[popdark, thick] (-1.5,-3.5) circle (0.4);
    % Vis micrométrique (droite)
    \draw[popdark, thick] (1.5,-2.5) -- (1.5,-3.5) node[below] {Vis micrométrique};
    \draw[popdark, thick] (1.5,-3.5) circle (0.2);
    % Révolver / Porte-objectifs
    \draw[popdark, thick, fill=poporangeL!40] (-0.5,3.2) rectangle (0.5,3.8);
    % Objectifs (3 rectangles verticaux)
    \draw[popdark, thick, fill=popmuted!30] (-0.3,3.5) rectangle (0.3,2.5) node[midway, right=3pt] {Objectif $\times 4$};
    \draw[popdark, thick, fill=popmuted!30] (-0.1,3.5) rectangle (0.1,2.2) node[midway, right=3pt] {Objectif $\times 10$};
    \draw[popdark, thick, fill=popmuted!30] (0.1,3.5) rectangle (-0.1,1.9) node[midway, left=3pt] {Objectif $\times 40$};
    % Tube optique
    \draw[popdark, thick, fill=popblueL!20] (-0.6,4) -- (-0.6,5.5) -- (1.5,5.5) -- (1.5,4) -- cycle;
    % Oculaire
    \draw[popdark, thick, fill=poppurpleL!40] (-0.6,5.5) rectangle (0.2,6.2) node[midway, above=2pt] {Oculaire ($\times 10$)};
    % Miroir / Lampe
    \draw[popdark, thick, fill=popgoldL!50] (-1.2,-1.5) rectangle (1.2,-1.2) node[midway] {Lampe / Miroir};
    % Diaphragme
    \draw[popdark, thick, fill=popgreenL!40] (-0.8,-1.2) rectangle (0.8,-0.8) node[midway] {Diaphragme (Iris)};
    % Flèches lumière
    \draw[poporange, thick, ->] (0,-1.5) -- (0,-0.8);
    \draw[poporange, thick, ->] (0,2.5) -- (0,3.5);
    \draw[poporange, thick, ->] (0,3.5) -- (0,5.5);
    % Légendes principales
    \node[popdark, align=center] at (-4.5, 5.5) {\textbf{1. Oculaire}\\(grossit l'image\\réelle)};
    \node[popdark, align=center] at (3.5, 3.0) {\textbf{2. Objectifs}\\(\textbf{$\times 4$, $\times 10$, $\times 40$})};
    \node[popdark, align=center] at (3.5, 0.5) {\textbf{3. Platine}\\(support de la lame)};
    \node[popdark, align=center] at (-4.5, -1.0) {\textbf{4. Diaphragme}\\(contrôle contraste)};
    \node[popdark, align=center] at (-4.5, -2.8) {\textbf{5. Vis macrométrique}\\(mise au point grossière)};
    \node[popdark, align=center] at (3.5, -2.8) {\textbf{6. Vis micrométrique}\\(mise au point fine)};
    \node[popdark, align=center] at (-4.5, 4.2) {\textbf{7. Révolver}\\(change objectif)};
    \node[popdark, align=center] at (3.5, 5.8) {\textbf{8. Bras / Colonne}\\(portage)};
    \node[popdark, align=center] at (0, -3.8) {\textbf{9. Pied / Base}\\(stabilité)};
\end{tikzpicture}
\legende{Schéma du microscope optique de laboratoire (type monoculaire) avec ses organes principaux. Le trajet de la lumière (flèches orange) va de la source (lampe/miroir) vers l'oculaire en traversant le diaphragme, la préparation, l'objectif et le tube.}
\end{popfigure}

\courssub{Préparation des échantillons : lame, lamelle et coloration}

Pour être observés, les microorganismes doivent être déposés sur une \textbf{lame} (lame porte-objet), recouverts d'une \textbf{lamelle} (couvre-objet) et souvent colorés pour augmenter le contraste.

\begin{experience}[Préparation d'un frottis et coloration simple (Bleu de méthylène)]
\textbf{Matériel :} Lame et lamelle propres, boucle de platine (ou aiguille), source de chaleur (bec Bunsen ou allume-gaz), bleu de méthylène (solution aqueuse \SI{1}{\%}), eau distillée, papier buvard, microscope.
\textbf{Échantillons types :} Levures de boulanger (\emph{Saccharomyces cerevisiae}), Moisissure (\emph{Penicillium} ou \emph{Aspergillus} sur pain/fruit), Bactéries lactiques (yaourt nature) ou préparation prête (ex. \emph{Bacillus subtilis}).

\textbf{Protocole :}
\begin{enumerate}
    \item \textbf{Nettoyage} : passe la lame et la lamelle à la flamme (côté frottis vers le haut) ou essuie-les à l'alcool \SI{70}{\%} / papier lentille. \emph{But : éliminer poussières et graisses.}
    \item \textbf{Dépôt (Frottis)} :
    \begin{itemize}
        \item \emph{Levures / Bactéries (liquide) :} dépose une \textbf{goutte d'eau distillée} au centre de la lame. Prélève une infime quantité d'échantillon (pointe de boucle) et émulsionne dans la goutte pour obtenir une couche fine. Étale en cercle d'environ \SI{1}{\centi\meter}.
        \item \emph{Moisissures (solide) :} dépose une goutte d'eau, prélève avec une aiguille stérilisée quelques filaments (mycélium) et spores à la périphérie de la colonie (zone jeune), émulsionne délicatement.
    \end{itemize}
    \item \textbf{Séchage et Fixation} : laisse sécher à l'air (ou passe rapidement à la flamme, face préparation vers le haut, sans surchauffer). \emph{But : coller les cellules à la lame (fixation thermique).}
    \item \textbf{Coloration} : recouvre le frottis sec de \textbf{bleu de méthylène} (1 à 2 min). \emph{But : colorer les structures acides (noyau, cytoplasme) en bleu pour les rendre visibles.}
    \item \textbf{Rinçage et Séchage} : rince délicatement à l'eau distillée (flacon laveur ou bécher) pour enlever l'excès de colorant. Tamponne le pourtour avec du papier buvard (ne frotte pas la préparation). Laisse sécher.
    \item \textbf{Montage} : dépose une \textbf{goutte d'eau} (ou de glycérol pour conservation) sur la préparation. Pose la \textbf{lamelle} en l'inclinant (\SI{45}{\degree}) pour chasser les bulles d'air. Tamponne l'excès d'eau.
\end{enumerate}

\textbf{Observations attendues (Grossissement $\times 400$ = Oculaire $\times 10$ $\times$ Objectif $\times 40$) :}
\begin{itemize}
    \item \textbf{Levures :} Cellules isolées, ovales ou sphériques (\SIrange{3}{8}{\micro\meter}), cytoplasme bleu homogène, noyau parfois visible (bleu foncé), \textbf{bourgeonnement} fréquent (petite cellule fille attachée à la mère).
    \item \textbf{Bactéries (bacilles) :} Petits bâtonnets (\SIrange{1}{5}{\micro\meter}), alignés en chaînes ou isolés, bleu uniforme (pas de noyau visible).
    \item \textbf{Moisissures :} \textbf{Hyphes} (filaments longs, transparents, septés ou non), \textbf{conidiophores} (pieds porteurs) terminés par des \textbf{conidies} (spores arrondies en chaînes ou bouquets, ex. \emph{Penicillium} = pinceau).
\end{itemize}
\end{experience}

\begin{methode}[Calcul du grossissement total]
Le grossissement total ($G_t$) est le produit du grossissement de l'oculaire ($G_{oc}$) par celui de l'objectif ($G_{ob}$) utilisé.
\[
G_t = G_{oc} \times G_{ob}
\]
\textbf{Exemple standard :} Oculaire $\times 10$ + Objectif $\times 40$ $\Rightarrow$ $G_t = 10 \times 40 = \textbf{$\times 400$}$.
\begin{itemize}
    \item Objectif $\times 4$ $\Rightarrow$ $\times 40$ (recherche, vue d'ensemble).
    \item Objectif $\times 10$ $\Rightarrow$ $\times 100$ (détails colonies, filaments).
    \item Objectif $\times 40$ $\Rightarrow$ $\times 400$ (détails cellulaires, morphologie).
    \item Objectif $\times 100$ (immersion huile) $\Rightarrow$ $\times 1000$ (bactéries, détails fins) — \emph{rare en 3ème, nécessite huile d'immersion}.
\end{itemize}
\end{methode}

\begin{exemplebox}[Calcul de la taille réelle à partir du dessin]
Au grossissement $\times 400$, tu dessines une levure qui mesure \textbf{20 mm} (2 cm) de long sur ton papier. Quelle est sa taille réelle ?
\[
\text{Taille réelle} = \frac{\text{Taille sur le dessin}}{\text{Grossissement total}} = \frac{20 \text{ mm}}{400} = 0,05 \text{ mm} = \mathbf{50\ \mu\text{m}}.
\]
\textbf{Attention :} Une levure typique fait \SIrange{5}{10}{\micro\meter}. Si tu trouves 50 $\mu$m, ton dessin est trop grand ou le grossissement noté est faux. Vérifie toujours la cohérence avec les tailles connues (voir l'encadré « À retenir » ci-dessous).
\end{exemplebox}

\courssub{Conduite de l'observation et réalisation des schémas}

\begin{attention}[Erreur fréquente — Ordre des grossissements]
\textbf{Ne commence JAMAIS par l'objectif $\times 40$ (ou $\times 100$).}
\begin{itemize}
    \item Risque 1 : l'objectif, très long, écrase la lamelle et la lame (casse) en descendant (vis macrométrique).
    \item Risque 2 : le champ est trop petit, tu ne trouves pas l'échantillon.
    \item \textbf{Règle d'or :} $\times 4$ (ou $\times 10$) $\rightarrow$ centrer $\rightarrow$ $\times 10$ $\rightarrow$ centrer $\rightarrow$ $\times 40$ (uniquement vis micrométrique).
\end{itemize}
\end{attention}

\begin{methode}[Règles du schéma d'observation scientifique (Exigible au BEPC)]
Un schéma d'observation au microscope n'est \textbf{pas un dessin artistique}, c'est un \textbf{document scientifique codifié}. Il doit comporter \textbf{obligatoirement} :
\begin{enumerate}
    \item \textbf{Titre} : Précis, en haut, souligné. Ex: \emph{« Levures de boulanger (Saccharomyces cerevisiae) observées au microscope optique, grossissement $\times 400$ »}.
    \item \textbf{Grossissement total} : Noté explicitement (ex: $\times 400$).
    \item \textbf{Dessin au crayon à papier (HB ou 2H)} : Propre, traits nets, sans gommages visibles, sans coloriage (pas de feutre, pas de crayon de couleur). Représente \textbf{ce que tu vois} (quelques cellules représentatives, pas des centaines), en respectant les proportions.
    \item \textbf{Légendes} : Au crayon, traits de rappel (trait droit, règle) ne se croisant pas, terminés par le nom de la structure. Ex: \emph{Paroi cellulaire}, \emph{Cytoplasme}, \emph{Noyau}, \emph{Bourgeon}.
    \item \textbf{Échelle (optionnel mais recommandé)} : Barre d'échelle (ex: 10 $\mu$m) calculée.
\end{enumerate}
\end{methode}

\begin{popfigure}
\begin{tikzpicture}[scale=1.2, every node/.style={font=\small}]
    % Cadre du champ microscopique
    \draw[popdark, thick] (0,0) circle (4cm);
    % Fond léger
    \fill[popcream!50] (0,0) circle (4cm);
    % Levures (dessin manuel de 4 cellules représentatives)
    % Cellule 1 (centre-haut)
    \begin{scope}[shift={(0,1.0)}]
        \draw[popdark, thick, fill=popblueL!20] (0,0) ellipse (0.6cm and 0.48cm);
        \fill[popblue!60] (0.06,0.06) circle (0.07cm);
        \draw[popdark, thick, fill=popblueL!40] (0.48,-0.12) ellipse (0.21cm and 0.15cm);
        \fill[popblue!60] (0.51,-0.09) circle (0.03cm);
    \end{scope}
    % Cellule 2 (gauche)
    \begin{scope}[shift={(-1.5,0.5)}, rotate=20]
        \draw[popdark, thick, fill=popblueL!20] (0,0) ellipse (0.5cm and 0.4cm);
        \fill[popblue!60] (0.05,0.05) circle (0.06cm);
    \end{scope}
    % Cellule 3 (droite)
    \begin{scope}[shift={(1.2,1.3)}, rotate=-15]
        \draw[popdark, thick, fill=popblueL!20] (0,0) ellipse (0.55cm and 0.44cm);
        \fill[popblue!60] (0.055,0.055) circle (0.066cm);
        \draw[popdark, thick, fill=popblueL!40] (0.44,-0.11) ellipse (0.19cm and 0.14cm);
        \fill[popblue!60] (0.47,-0.08) circle (0.028cm);
    \end{scope}
    % Cellule 4 (bas)
    \begin{scope}[shift={(0.3,-1.5)}, rotate=30]
        \draw[popdark, thick, fill=popblueL!20] (0,0) ellipse (0.45cm and 0.36cm);
        \fill[popblue!60] (0.045,0.045) circle (0.054cm);
        \draw[popdark, thick, fill=popblueL!40] (0.36,-0.09) ellipse (0.16cm and 0.11cm);
        \fill[popblue!60] (0.38,-0.07) circle (0.023cm);
    \end{scope}
    % Légendes (flèches)
    \draw[popdark, ->, thick] (1.5, 2.5) -- (0.8, 1.3) node[above right, align=left] {Paroi cellulaire};
    \draw[popdark, ->, thick] (1.5, 2.0) -- (0.5, 1.0) node[above right, align=left] {Cytoplasme};
    \draw[popdark, ->, thick] (1.5, 1.5) -- (0.1, 0.6) node[above right, align=left] {Noyau};
    \draw[popdark, ->, thick] (-2.5, 2.5) -- (-1.0, 1.2) node[above left, align=right] {Bourgeon (cellule fille)};
    \draw[popdark, ->, thick] (-2.5, 2.0) -- (-1.3, 0.6) node[above left, align=right] {Cicatrice de bourgeonnement};
    % Barre d'échelle
    \draw[popdark, thick] (-3.5, -3.5) -- (-2.5, -3.5) node[midway, below] {10 $\mu$m};
    % Titre simulé (hors dessin)
    \node[popdark, align=center, font=\bfseries\small] at (0, 4.5) {Levures de boulanger -- Grossissement $\times 400$};
\end{tikzpicture}
\legende{Schéma type d'observation de levures bourgeonnantes (\emph{Saccharomyces cerevisiae}) au grossissement $\times 400$. Respect des conventions : titre, grossissement, dessin au crayon, légendes à la règle, barre d'échelle. On note les cellules mères ovales, les noyaux (petits points sombres), les bourgeons (cellules filles) et les cicatrices de bourgeonnement (anneaux sur la paroi mère).}
\end{popfigure}

\courssub{Interprétation : identification des groupes microbiens}

L'observation au microscope optique (limite de résolution $\approx \SI{0,2}{\micro\meter}$) permet de distinguer les \textbf{champignons microscopiques} (levures, moisissures) et les \textbf{bactéries} (grandes bacilles, coccis). Les \textbf{virus} (\SIrange{20}{300}{\nano\meter}) restent \textbf{invisibles} au microscope optique (nécessitent le microscope électronique).

\begin{aretenir}[Critères d'identification microscopique ($\times 400$)]
\begin{center}
\begin{tabularx}{\linewidth}{|l|X|X|X|}\hline
\textbf{Groupe} & \textbf{Taille typique} & \textbf{Morphologie clé} & \textbf{Structures visibles} \\ \hline
\textbf{Levures} (Champignons) & \SIrange{3}{10}{\micro\meter} & Cellules isolées, ovales/sphériques, \textbf{bourgeonnement} & Paroi épaisse, cytoplasme dense, noyau unique, vacuole \\ \hline
\textbf{Moisissures} (Champignons) & Hyphes : \SIrange{2}{10}{\micro\meter} larg. & \textbf{Filaments (hyphes)} ramifiés (mycélium), \textbf{spores (conidies)} en chaînes/bouquets & Septa (cloisons) parfois, conidiophores spécialisés (\emph{Penicillium}, \emph{Aspergillus}) \\ \hline
\textbf{Bactéries} (Procaryotes) & \SIrange{0,5}{5}{\micro\meter} & \textbf{Bâtonnets (bacilles)}, \textbf{Sphères (coccis)}, Spirilles & \textbf{Pas de noyau}, cytoplasme granulaire uniforme, parfois flagelles (coloration spé.) \\ \hline
\textbf{Protistes} (ex. Paramecium) & \SIrange{50}{300}{\micro\meter} & Cellule unique complexe, cils/vacuoles contractiles & Noyau (macro/micronoyau), vacuoles digestives, cils vibratiles \\ \hline
\end{tabularx}
\end{center}
\end{aretenir}

\begin{savaistu}[Le virus, le grand absent du microscope optique]
Au Cameroun, comme partout, le diagnostic du VIH/Sida, de la fièvre jaune, de la rougeole ou du COVID-19 ne se fait \textbf{jamais} en observant le virus au microscope optique. On utilise :
\begin{itemize}
    \item La \textbf{sérologie} (recherche d'anticorps ou d'antigènes par test rapide ELISA / TDR).
    \item La \textbf{biologie moléculaire} (PCR : amplification de l'ARN/DN viral).
    \item Le \textbf{microscope électronique} (transmission ou balayage) pour la morphologie virale (capside, enveloppe) — uniquement en centre de recherche spécialisé (ex. Centre Pasteur du Cameroun, IMPM).
\end{itemize}
Rappelle-toi : \textbf{« Pas de noyau visible $\neq$ virus »}. Une bactérie n'a pas de noyau non plus, mais elle \textbf{se voit} au microscope optique.
\end{savaistu}

\courssub{Compte rendu de TP : rédaction et justification de la conclusion}

Le compte rendu (CR) de TP suit la \textbf{démarche d'investigation scientifique}. C'est l'exercice type du BEPC (épreuve pratique ou écrite).

\begin{methode}[Structure type du Compte Rendu de TP (BEPC)]
\begin{enumerate}
    \item \textbf{Titre du TP} : Ex: \emph{Observation de microorganismes au microscope optique}.
    \item \textbf{Problématique / Question} : Ex: \emph{Comment distinguer au microscope optique une levure, une moisissure et une bactérie ?}
    \item \textbf{Hypothèses} : Ex: \emph{Les levures montrent des bourgeons ; les moisissures montrent des filaments et spores ; les bactéries sont de petits bâtonnets sans noyau.}
    \item \textbf{Matériel et Méthode} : Liste + Protocole numéroté (préparation, coloration, observation $\times 40 \rightarrow \times 100 \rightarrow \times 400$).
    \item \textbf{Résultats} :
    \begin{itemize}
        \item Tableau récapitulatif (Groupe, Grossissement, Description morphologique).
        \item \textbf{Schémas légendés} (au moins 2 groupes différents, dont levure).
    \end{itemize}
    \item \textbf{Interprétation / Discussion} : Compare tes dessins aux schémas types (voir l'encadré ci-dessus). Justifie l'appartenance à tel groupe par les critères observés (ex: « Présence de bourgeons $\Rightarrow$ Levure » ; « Filaments septés + conidies en pinceau $\Rightarrow$ \emph{Penicillium} (Moisissure) » ; « Petits bâtonnets alignés, pas de noyau $\Rightarrow$ Bacilles (Bactéries) »).
    \item \textbf{Conclusion} : Réponse directe à la problématique, validée par les résultats. Ex: \emph{« Le microscope optique à $\times 400$ permet de différencier les trois groupes : levures (cellules bourgeonnantes), moisissures (hyphes et spores), bactéries (bacilles/coccis procaryotes). Les virus ne sont pas observables. »}
    \item \textbf{Critique / Amélioration} : Limites (résolution, coloration unique), propositions (coloration de Gram, microscope électronique).
\end{enumerate}
\end{methode}

\begin{exoresolu}[Rédaction d'une conclusion argumentée]
\textbf{Énoncé :} Lors d'un TP, un élève observe au grossissement $\times 400$ une préparation colorée au bleu de méthylène. Il voit des structures filamenteuses ramifiées, transparentes, portant à leur extrémité des chaînes de petites cellules sphériques. Il dessine ce qu'il voit. Rédige la conclusion argumentée de cet élève.

\textbf{Solution détaillée :}
\begin{enumerate}
    \item \textbf{Analyse des observations :} Structures filamenteuses = \textbf{hyphes}. Ramifiées = \textbf{mycélium}. Extrémités porteuses de chaînes de cellules sphériques = \textbf{conidiophores} portant des \textbf{conidies} (spores asexuées) en chaînes.
    \item \textbf{Comparaison aux modèles (Savoirs) :} Cette organisation (mycélium + conidiophores + conidies en chaînes) est caractéristique des \textbf{moisissures} (champignons filamenteux), genre type \emph{Aspergillus} ou \emph{Penicillium}.
    \item \textbf{Exclusion :} Ce ne sont pas des levures (pas de cellules isolées bourgeonnantes), ni des bactéries (taille trop grande, structures eucaryotes complexes, spores), ni des protistes (pas de cils, organisation fongique).
    \item \textbf{Rédaction finale :}
\end{enumerate}
\tcblower
\textbf{Conclusion type :}
\emph{« La préparation observée au grossissement $\times 400$ montre un mycélium composé d'hyphes ramifiés et septés, portant des conidiophores terminés par des chaînes de conidies sphériques. Ces caractères morphologiques correspondent au mode de reproduction asexuée des \textbf{moisissures} (champignons microscopiques filamenteux), probablement du genre \emph{Aspergillus} ou \emph{Penicillium}. Il ne s'agit pas de levures (absence de bourgeonnement), ni de bactéries (taille $> \SI{2}{\micro\meter}$, structures eucaryotes visibles, spores), ni de virus (invisibles à ce grossissement). »}
\end{exoresolu}

\begin{exercice}[Entraînement BEPC — Schéma et Interprétation]
Tu observes au microscope (Oculaire $\times 10$, Objectif $\times 40$) une préparation de yaourt nature colorée au bleu de méthylène. Tu vois de nombreux petits bâtonnets (longueur $\approx \SI{3}{\micro\meter}$), souvent alignés par paires ou en chaînes courtes, au cytoplasme bleu uniforme, sans noyau distinct.
\begin{enumerate}
    \item Calcule le grossissement total utilisé.
    \item Dessine (schéma au crayon, 3--4 cellules) ce que tu observes en respectant les conventions (titre, grossissement, légendes : paroi, cytoplasme). Indique une barre d'échelle de \SI{5}{\micro\meter}.
    \item Identifie le groupe microbien et l'espèce probable (genre). Justifie par deux arguments morphologiques.
    \item Pourquoi n'utilise-t-on pas l'objectif $\times 100$ (immersion) en classe de 3ème pour cette observation ?
\end{enumerate}
\end{exercice}

\begin{corrige}[Exercice n°1]
\begin{enumerate}
    \item $G_t = 10 \times 40 = \textbf{$\times 400$}$.
    \item \textbf{Schéma attendu :} Titre : « Bactéries lactiques (yaourt) $\times 400$ ». Dessin : 3--4 bâtonnets alignés (paires/chaînes), traits nets crayon. Légendes : Paroi cellulaire, Cytoplasme. Barre \SI{5}{\micro\meter} = \SI{2}{\milli\meter} sur le papier ($5 \times 400 = 2000 \mu$m = 2 mm).
    \item \textbf{Groupe :} \textbf{Bactéries} (Procaryotes). \textbf{Genre probable :} \emph{Lactobacillus} (bacilles Gram +, chaînes).
    \item \textbf{Arguments :} (1) Taille micrométrique (\SI{3}{\micro\meter}) compatible bactérie ; (2) Absence de noyau délimité (procaryote) ; (3) Forme bacillaire en chaînes (morphologie \emph{Lactobacillus}).
    \item L'objectif $\times 100$ nécessite de l'\textbf{huile d'immersion} (manipulation délicate, risque de salir les autres objectifs, nettoyage impératif) et n'est pas indispensable pour voir la morphologie bacillaire en chaînes à $\times 400$. De plus, les microscopes de collège n'en sont pas toujours équipés.
\end{enumerate}
\end{corrige}

\begin{attention}[Pièges à éviter dans le CR de TP (BEPC)]
\begin{itemize}
    \item \textbf{Oublier le grossissement} sur le schéma ou dans le titre ($\rightarrow$ 0 pt pour le schéma).
    \item \textbf{Dessiner au stylo / feutre / couleur} ($\rightarrow$ Pénalité, non scientifique).
    \item \textbf{Légendes sans trait de rappel} (flèches courbes, traits qui se croisent) ou \textbf{écrites sur le dessin}.
    \item \textbf{Confondre « Noyau » et « Nucléoïde »} : écrire « Noyau visible » pour une bactérie est une \textbf{erreur scientifique grave} (procaryote = pas de noyau). Écrire « cytoplasme granulaire, pas de noyau » est correct.
    \item \textbf{Conclusion sans justification} : « C'est une bactérie » sans citer les critères observés (taille, forme, absence de noyau) = conclusion incomplète.
\end{itemize}
\end{attention}

\coursec{Nutrition et respiration des microbes}

\courssub{Transition depuis l'observation microscopique}

Dans la section précédente, nous avons mis en évidence, grâce au microscope optique (TP N°5), la diversité morphologique des microorganismes : bactéries sous forme de bacilles, coccis ou spirilles ; levures bourgeonnantes ; moisissures à mycélium et spores. Nous les avons \emph{vus}, mais sont-ils \emph{vivants} au sens biologique strict ? Tout organisme vivant se caractérise par deux fonctions fondamentales : la \textbf{nutrition} (apport de matière et d'énergie pour se construire et fonctionner) et la \textbf{respiration} (libération d'énergie utilisable par la cellule). Cette section vise à démontrer expérimentalement que les microbes se nourrissent et respirent, et à comprendre comment ces processus conditionnent leur développement dans notre environnement — y compris dans nos assiettes.

\courssub{Les besoins nutritifs des microorganismes : une expérience fondatrice}

\emph{Observation initiale.} Si l'on laisse un morceau de pain humide à l'air libre, un duvet vert-bleu (moisissure) apparaît en quelques jours. Si l'on prépare un bouillon de viande stérilisé et qu'on l'expose à l'air, il se trouble au bout de 24 à 48 h. Dans les deux cas, des microbes invisibles au départ se sont multipliés. D'où leur viennent la matière et l'énergie nécessaires à cette prolifération ?

\begin{experience}[Mise en évidence des besoins nutritifs : culture sur gélose]
\textbf{Matériel :} Boîtes de Pétri, gélose nutritive (milieu solide gélifié contenant de l'eau, des sels minéraux, des peptones — hydrolysats de protéines — et du glucose), brûleur Bunsen, anses de platine, échantillons divers (air, eau de mare, doigts non lavés, surface de table), étuve à \SI{30}{\degreeCelsius}.

\textbf{Protocole :}
\begin{enumerate}
    \item Préparer des boîtes de Pétri contenant de la gélose nutritive (milieu complet) et, pour le contrôle, des boîtes avec de l'eau gélifiée seule (agar + eau, sans nutriments organiques ni sels minéraux).
    \item Stériliser l'anse à la flamme, la laisser refroidir.
    \item Ensemencer chaque boîte par étalement (surface) ou par piquage (profondeur) avec l'échantillon testé.
    \item Refermer les boîtes, les sceller avec du parafilm (pour éviter les contaminations aériennes ultérieures et la déshydratation).
    \item Incuber à \SI{30}{\degreeCelsius} pendant 24 à 48 h.
\end{enumerate}

\textbf{Observations :}
\begin{itemize}
    \item Sur gélose nutritive : apparition de nombreuses \textbf{colonies} visibles à l'œil nu, de formes, couleurs et textures variées (blanches, jaunes, pigmentées, muqueuses, sèches). Chaque colonie provient de la multiplication d'une seule cellule mère (ou d'un petit groupe) déposée lors de l'ensemencement.
    \item Sur eau gélifiée (témoin négatif) : \textbf{aucune colonie} visible, même après plusieurs jours.
\end{itemize}

\textbf{Interprétation :}
La croissance microbienne (multiplication cellulaire $\rightarrow$ colonie visible) \textbf{nécessite la présence de matière organique (source de carbone et d'énergie) et de sels minéraux (sources d'azote, phosphore, soufre, magnésium, potassium, oligo-éléments)}. L'eau est le solvant indispensable à tous les échanges métaboliques. Le témoin (eau + agar) prouve que l'agar lui-même n'est pas assimilable par la majorité des bactéries courantes (sauf rares agarolytiques) et que l'eau seule ne suffit pas.

\textbf{Conclusion :} Les microorganismes sont des êtres \textbf{hétérotrophes} (pour la grande majorité) : ils prélèvent dans le milieu extérieur des molécules organiques toutes faites (glucides, protéines, lipides) et des sels minéraux dissous. L'eau est le milieu réactionnel obligatoire.
\end{experience}

\begin{propriete}[Besoins nutritifs fondamentaux des microbes -- Établie expérimentalement]
Pour se développer et se multiplier, les microorganismes ont besoin de :
\begin{enumerate}
    \item \textbf{Eau} : milieu de vie, solvant des nutriments et des déchets, participant aux réactions d'hydrolyse.
    \item \textbf{Sels minéraux} : azote (N), phosphore (P), potassium (K), magnésium (Mg), soufre (S), fer (Fe), oligo-éléments — pour la synthèse des protéines, acides nucléiques, ATP, coenzymes, paroi.
    \item \textbf{Matière organique} : source de carbone et d'énergie chimique. Principalement des \textbf{glucides} (glucose, saccharose, amidon hydrolysé) et des \textbf{protéines} (acides aminés, peptides). Certaines bactéries utilisent des lipides ou des hydrocarbures.
\end{enumerate}
L'absence d'un seul de ces facteurs arrête la croissance.
\end{propriete}

\begin{attention}[Piège fréquent : confusion entre « se nourrir » et « manger »]
Les microbes \textbf{n'ont pas de bouche, ni de tube digestif}. Ils n'ingèrent pas des morceaux solides. Ils \textbf{absorbent uniquement des molécules \emph{dissoutes} de faible masse molaire} (glucose, acides aminés, ions minéraux) à travers leur membrane cytoplasmique. Les grosses molécules (amidon, protéines natives, cellulose) doivent être \textbf{hydrolysées extracellulairement} par des \textbf{enzymes sécrétées} (exoenzymes : amylases, protéases, cellulases, lipases) avant d'être absorbées. \textbf{Ne pas écrire : « la bactérie mange le sucre » mais « la bactérie absorbe le glucose dissous ».}
\end{attention}

\courssub{Modes de nutrition selon les groupes microbiens}

\courssub{Nutrition des bactéries : absorption et décomposition}

Les bactéries sont des procaryotes unicellulaires. Leur paroi (peptidoglycane) est poreuse aux petites molécules ; la membrane cytoplasmique (phospholipides + protéines transporteurs) assure la \textbf{perméabilité sélective} et le \textbf{transport actif} (consommateur d'ATP) des nutriments contre gradient de concentration.

\begin{itemize}
    \item \textbf{Saprophytes (décomposeurs) :} La grande majorité des bactéries du sol, de l'eau, des déchets. Elles sécrètent des exoenzymes qui décomposent la matière organique morte (protéines $\rightarrow$ acides aminés ; polysaccharides $\rightarrow$ sucres simples). Elles \textbf{recyclent} les éléments chimiques (C, N, P, S) dans les écosystèmes. Exemple : \emph{Bacillus subtilis} dans le sol, bactéries putréfiantes dans les fosses septiques.
    \item \textbf{Parasites / Pathogènes :} Elles absorbent les nutriments directement dans les tissus de l'hôte vivant (sang, lymphe, mucus), lui causant des lésions. Exemple : \emph{Salmonella typhi} (fièvre typhoïde), \emph{Mycobacterium tuberculosis}.
    \item \textbf{Symbiotes (mutualistes) :} Échange d'avantages. Exemple : bactéries de la flore intestinale (colibacilles) qui synthétisent la vitamine K et digèrent la cellulose résiduelle ; bactéries \emph{Rhizobium} dans les nodules des légumineuses (fixation de l'azote atmosphérique).
    \item \textbf{Autotrophes (rares) :} Bactéries chimiosynthétiques (nitrifiantes : \emph{Nitrosomonas}, \emph{Nitrobacter}) qui oxydent l'ammonium en nitrites puis nitrates pour obtenir de l'énergie, fixant le CO$_2$ comme source de carbone. Bactéries photosynthétiques (cyanobactéries, bactéries pourpres soufrées).
\end{itemize}

\courssub{Nutrition des champignons microscopiques : levures et moisissures}

Les champignons (eucaryotes) partagent le mode de nutrition \textbf{par absorption} (osmotrophie), mais leur morphologie diffère.

\begin{exemplebox}[La levure de boulanger (\emph{Saccharomyces cerevisiae})]
Unicellulaire, ovoïde (5--10 µm). Elle se nourrit de \textbf{sucres simples} (glucose, fructose, saccharose, maltose) qu'elle absorbe par transporteurs membranaires. Elle possède des enzymes intracellulaires (invertase pour hydrolyser le saccharose, maltase pour le maltose). Sur un milieu sucré (jus de fruit, pâte à pain, moût de bière), elle se multiplie rapidement par bourgeonnement.
\end{exemplebox}

\begin{exemplebox}[Les moisissures (ex. \emph{Aspergillus}, \emph{Penicillium}, \emph{Rhizopus})]
Filamenteuses : le corps végétatif est un \textbf{mycélium} ramifié composé d'hyphes (filaments de 5--10 µm de large). L'extrémité des hyphes (apex) sécrète puissamment des exoenzymes (amylases, pectinases, cellulases, protéases) qui \textbf{digèrent le substrat solide} (pain, fruit, fromage, bois, cuir) en molécules diffusibles. Les nutriments sont absorbés sur toute la surface du mycélium et transportés vers les zones de croissance. C'est pourquoi elles colonisent les aliments solides en surface et en profondeur.
\end{exemplebox}

\begin{savaistu}[Le « bâton de manioc » et le yaourt : biotechnologie traditionnelle camerounaise]
Au Cameroun, la transformation du manioc en \textbf{bâton de manioc} (bobolo, miondo) et du lait en \textbf{yaourt} (lait caillé) repose entièrement sur la nutrition et la respiration de microbes utiles.
\begin{itemize}
    \item \textbf{Bâton de manioc :} Le manioc trempé (rettissage) est envahi par des bactéries lactiques (\emph{Lactobacillus}, \emph{Leuconostoc}) et des levures. Elles fermentent les glucides (amidon hydrolysé) en acide lactique, acide acétique, éthanol, CO$_2$. L'acidification (pH $\approx$ 4) empêche le développement de microbes pathogènes et dégrade les cyanogénosides toxiques (linamarine) du manioc amer. La pâte devient digeste, savoureuse et se conserve des semaines.
    \item \textbf{Yaourt :} Inoculation de ferments lactiques (\emph{Streptococcus thermophilus}, \emph{Lactobacillus bulgaricus}) dans du lait pasteurisé. Fermentation du lactose en acide lactique $\rightarrow$ coagulation des caséines (caillage), acidification, arômes. Conservation naturelle par acidité.
\end{itemize}
Ces pratiques ancestrales sont des applications maîtrisées de la physiologie microbienne !
\end{savaistu}

\courssub{Respiration des microorganismes : expérience clé et diversité métabolique}

\emph{Question déclenchante :} Comment les microbes libèrent-ils l'énergie stockée dans les molécules organiques qu'ils absorbent ?

\begin{experience}[Mise en évidence du dégagement de CO$_2$ par la levure -- Fermentation alcoolique]
\textbf{Matériel :} Deux tubes à essai, levure de boulanger fraîche (ou sèche réhydratée), eau tiède (\SI{35}{\degreeCelsius}), sucre en poudre (saccharose), eau de chaux (solution saturée d'hydroxyde de calcium \ce{Ca(OH)2}, limpide), bouchons à deux trous, tubes de liaison, étuve ou bain-marie à \SI{35}{\degreeCelsius}.

\textbf{Dispositif expérimental :}
\begin{popfigure}
\begin{tikzpicture}[scale=0.9, every node/.style={font=\small}]
    % Couleurs
    \colorlet{liquid}{popblue!30}
    \colorlet{gas}{poporange!50}
    \colorlet{tube}{popgray!50}
    \colorlet{yeast}{poporange!70!popdark!60}
    % Tube 1 : Milieu de culture
    \draw[tube, line width=2pt] (0,0) rectangle (3,6);
    \fill[liquid] (0.1,0.1) rectangle (2.9,4);
    % Levure au fond
    \foreach \i in {1,...,30} {
        \pgfmathsetmacro{\x}{0.2 + rand*2.5}
        \pgfmathsetmacro{\y}{0.2 + rand*1.5}
        \fill[yeast] (\x,\y) circle (0.06);
    }
    % Bulles CO2
    \foreach \i in {1,...,15} {
        \pgfmathsetmacro{\x}{0.5 + rand*2}
        \pgfmathsetmacro{\y}{4.2 + rand*1.5}
        \fill[gas, opacity=0.7] (\x,\y) circle (0.08);
    }
    \draw[->, poporange, thick] (1.5,5.5) -- (1.5,6.5) node[above] {Gaz dégagé};
    % Tube de liaison
    \draw[tube, line width=2pt] (1.3,6) -- (1.3,7) -- (5.7,7) -- (5.7,5.5);
    \draw[->, poporange, thick] (3.5,7.2) -- (3.5,7.5) node[above] {\ce{CO2}};
    % Tube 2 : Eau de chaux
    \draw[tube, line width=2pt] (5,0) rectangle (7,6);
    \fill[popblue!10] (5.1,0.1) rectangle (6.9,4);
    % Bulles entrant
    \foreach \i in {1,...,10} {
        \pgfmathsetmacro{\x}{5.3 + rand*1.4}
        \pgfmathsetmacro{\y}{0.5 + rand*2.5}
        \fill[gas, opacity=0.7] (\x,\y) circle (0.06);
    }
    % Trouble (précipité CaCO3)
    \fill[popgray!30] (5.1,0.1) rectangle (6.9,4);
    \draw[popgray, densely dashed] (5.1,2) -- (6.9,2);
    \node[popgray] at (6,2.3) {Trouble : \ce{CaCO3}};
    % Légendes
    \node[align=center] at (1.5,-0.5) {Tube 1 : Levure +\\Eau sucrée (\SI{35}{\degreeCelsius})};
    \node[align=center] at (6,-0.5) {Tube 2 : Eau de chaux\\\ce{Ca(OH)2}};
    \node[align=left, popdark] at (0,6.5) {1. Milieu nutritif};
    \node[align=left, popdark] at (0,6.0) {2. Levures};
    \node[align=left, popdark] at (0,5.5) {3. Gaz (\ce{CO2})};
    \node[align=left, popdark] at (0,5.0) {4. Tube de liaison};
    \node[align=left, popdark] at (0,4.5) {5. Eau de chaux};
    \node[align=left, popdark] at (0,4.0) {6. Trouble (\ce{CaCO3})};
\end{tikzpicture}
\legende{Dispositif classique de mise en évidence de la fermentation alcoolique par \emph{Saccharomyces cerevisiae}. Le \ce{CO2} produit barboté dans l'eau de chaux la trouble par formation de carbonate de calcium insoluble.}
\end{popfigure}

\textbf{Protocole :}
\begin{enumerate}
    \item Dans le tube 1 : mélanger 2 g de levure + 20 mL d'eau tiède + 5 g de sucre. Boucher hermétiquement avec le tube de liaison plongeant dans le tube 2.
    \item Dans le tube 2 : verser 15 mL d'eau de chaux limpide.
    \item Placer l'ensemble à \SI{35}{\degreeCelsius} (température optimale pour la levure).
    \item Observer pendant 30--60 min.
\end{enumerate}

\textbf{Observations :}
\begin{itemize}
    \item Dans le tube 1 : le mélange mousse, des bulles remontent en continu.
    \item Dans le tube de liaison : passage de bulles de gaz.
    \item Dans le tube 2 : l'eau de chaux devient \textbf{blanchâtre, trouble (laitier)} au passage des bulles. Un dépôt blanc (précipité) se forme au fond.
\end{itemize}

\textbf{Interprétation :}
\begin{itemize}
    \item L'eau de chaux (\ce{Ca(OH)2} aqueuse) est un \textbf{réactif spécifique du dioxyde de carbone \ce{CO2}} : \ce{Ca(OH)2 + CO2 -> CaCO3 v + H2O}. Le trouble prouve la présence de \ce{CO2} dans le gaz émis.
    \item La levure a consommé le sucre (glucose + fructose issu de l'hydrolyse du saccharose par l'invertase) et a dégagé du \ce{CO2}.
    \item L'expérience a lieu \textbf{sans apport d'air extérieur} (dispositif clos) : la levure n'a pas utilisé le dioxygène \ce{O2} de l'air. C'est une \textbf{fermentation} (respiration anaérobie).
\end{itemize}

\textbf{Bilan global de la fermentation alcoolique :}
\[
\ce{C6H12O6 -> 2 C2H5OH + 2 CO2 + Energie (2 ATP)}
\]
(Glucose $\rightarrow$ Éthanol + Dioxyde de carbone + Énergie utilisable par la cellule).

\textbf{Conclusion :} Les levures (et de nombreuses bactéries) sont capables de \textbf{respérer sans dioxygène} : elles dégradent le glucose de façon incomplète en éthanol et \ce{CO2}, libérant peu d'énergie (2 ATP par glucose), mais leur permettant de survivre et de se multiplier en milieu anaérobie.
\end{experience}

\begin{methode}[Test de l'eau de chaux -- Preuve de dégagement de \ce{CO2}]
Pour mettre en évidence la production de dioxyde de carbone par un organisme (microbe, plante au noir, animal) ou une réaction chimique :
\begin{enumerate}
    \item Préparer de l'\textbf{eau de chaux} : ajouter de l'eau à de la chaux vive (\ce{CaO}) ou éteinte (\ce{Ca(OH)2}), agiter, laisser décanter, filtrer. La solution limpide obtenue est saturée en \ce{Ca(OH)2} ($\approx$ \SI{1,5}{g/L} à \SI{20}{\degreeCelsius}).
    \item Faire \textbf{barbotter} le gaz à tester dans cette solution (tube plongeant au fond).
    \item \textbf{Résultat positif :} L'eau de chaux devient \textbf{trouble (blanchâtre)} par formation de \textbf{carbonate de calcium \ce{CaCO3}} (insoluble).
    \item \textbf{Contrôle :} Faire barbotter de l'air ambiant (\SI{0,04}{\%} \ce{CO2}) : trouble très lent ou absent. Faire barbotter de l'air expiré (\SI{4}{\%} \ce{CO2}) : trouble rapide.
\end{enumerate}
\textbf{Attention :} Si l'on continue à faire passer du \ce{CO2} en excès, le trouble disparaît : \ce{CaCO3 + CO2 + H2O -> Ca(HCO3)2} (bicarbonate de calcium soluble). Le test n'est donc valide que pour la \textbf{première apparition du trouble}.
\end{methode}

\courssub{Aérobies, anaérobies : la diversité des respirations microbiennes}

\begin{definition}[Respiration cellulaire (sens large)]
Ensemble des réactions d'oxydation de substrats organiques (glucose, acides gras, acides aminés) aboutissant à la production d'ATP. Selon l'accepteur final d'électrons, on distingue :
\begin{itemize}
    \item \textbf{Respiration aérobie :} Accepteur final = \textbf{dioxygène \ce{O2}}. Dégradation complète du glucose en \ce{CO2} + \ce{H2O}. Rendement élevé ($\approx$ 30--38 ATP/glucose).
    \item \textbf{Respiration anaérobie (fermentation) :} Accepteur final = \textbf{molécule organique} (pyruvate, acétaldéhyde, etc.) issue du substrat initial. Pas d'\ce{O2} requis. Dégradation incomplète. Rendement faible (2 ATP/glucose par phosphorylation au niveau du substrat).
\end{itemize}
\end{definition}

\begin{propriete}[Classification des microbes selon leurs exigences en \ce{O2} -- Établie expérimentalement (milieux de culture, tubes de Thioglycollate)]
\begin{center}
\begin{tabularx}{\linewidth}{|l|X|X|}\hline
\textbf{Groupe} & \textbf{Exigence en \ce{O2}} & \textbf{Exemples \& Conséquences} \\ \hline
\textbf{Aérobies stricts} (obligatoires) & \ce{O2} \textbf{indispensable} (accepteur final). Meurent sans \ce{O2}. & \emph{Mycobacterium tuberculosis}, \emph{Pseudomonas aeruginosa}, moisissures (\emph{Aspergillus}), levures (en présence d'\ce{O2}). Croissent en surface des milieux liquides. \\ \hline
\textbf{Anaérobies stricts} (obligatoires) & \ce{O2} \textbf{toxique} (absence de catalase, superoxyde dismutase). Meurent en présence d'\ce{O2}. & \emph{Clostridium tetani} (tétanos), \emph{Clostridium botulinum} (botulisme), \emph{Bacteroides} (flore intestinale). Croissent au fond des milieux profonds / anaérobiose. \\ \hline
\textbf{Anaérobies facultatifs} (ou aérobies facultatifs) & \ce{O2} \textbf{utilisé si présent} (respiration aérobie, haut rendement) ; \textbf{fermentation si absent} (bas rendement). Grande adaptabilité. & \emph{Escherichia coli}, \emph{Salmonella}, \emph{Staphylococcus}, \emph{Saccharomyces cerevisiae} (levure), bactéries lactiques. Colonisent des milieux variés. \\ \hline
\textbf{Microaérophiles} & Faible taux d'\ce{O2} requis (2--10 \%), \ce{O2} atmosphérique (\SI{21}{\%}) toxique. & \emph{Helicobacter pylori} (ulcère gastrique), \emph{Campylobacter}. \\ \hline
\end{tabularx}
\end{center}
\end{propriete}

\begin{exemplebox}[Fermentation lactique : le yaourt et le lait caillé]
Les \textbf{bactéries lactiques} (\emph{Lactobacillus}, \emph{Streptococcus}, \emph{Leuconostoc}) sont des anaérobies facultatifs (souvent aéro-intolérants). Elles fermentent le glucose (ou lactose, galactose) principalement en \textbf{acide lactique} :
\[
\ce{C6H12O6 -> 2 CH3-CHOH-COOH + 2 ATP} \quad \text{(Fermentation homolactique)}
\]
ou en acide lactique + éthanol + \ce{CO2} (hétérolactique).
L'acidification (pH 4--4,5) coagule les protéines du lait (caséines) $\rightarrow$ texture gélifiée, conservation par inhibition des putréfactifs. C'est la base du yaourt, du lait caillé (pendé), du fromage, de la choucroute, des cornichons, du silo d'ensilage.
\end{exemplebox}

\courssub{Conditions favorables au développement microbien : lien avec la conservation des aliments}

La vitesse de croissance microbienne (temps de génération, biomasse finale) dépend de paramètres physico-chimiques. Les comprendre permet de \textbf{maîtriser la conservation des aliments} (lutte contre les altérations et les toxi-infections).

\begin{popfigure}
\begin{tikzpicture}
\begin{axis}[
    width=\linewidth,
    height=8cm,
    axis lines=middle,
    xlabel={Température (\si{\degreeCelsius})},
    ylabel={Vitesse de croissance / Nombre de colonies (échelle arbitraire)},
    xmin=0, xmax=50,
    ymin=0, ymax=1.1,
    xtick={0,10,20,30,37,45,50},
    ytick=\empty,
    domain=0:50,
    samples=200,
    clip=false,
    tick label style={font=\small},
    label style={font=\small},
    grid=major,
    grid style={popline!50},
]
% Courbe gaussienne asymétrique type mésophile (optimum 37°C)
\addplot[popcurve, thick, popblue, domain=0:50] 
    {exp(-(x-37)^2/(2*8^2)) * (1 + 0.02*x)};
% Annotations
\node[popblue, anchor=south, font=\small\bfseries] at (axis cs:37,1.02) {Optimum \textasciitilde 37 °C};
\draw[popblue, dashed, thick] (axis cs:37,0) -- (axis cs:37,1);
\node[popmuted, anchor=north east, font=\footnotesize, align=right] at (axis cs:48,0.9) {Zone de\\ destruction\\ thermique};
\node[popmuted, anchor=north west, font=\footnotesize, align=left] at (axis cs:2,0.9) {Zone de\\ survie /\\ croissance\\ très lente};
\node[popgreen, anchor=south, font=\footnotesize] at (axis cs:4,0.15) {Réfrigération (0--4 °C)};
\node[poporange, anchor=south, font=\footnotesize] at (axis cs:20,0.4) {Température ambiante};
\node[popred, anchor=south, font=\footnotesize] at (axis cs:37,0.6) {Température corporelle};
\end{axis}
\end{tikzpicture}
\legende{Influence de la température sur la vitesse de croissance des bactéries mésophiles (majorité des pathogènes humains et des altérants alimentaires). L'optimum est proche de 37 °C. En dessous de 10 °C, la croissance est fortement ralentie (bactériostase) ; au-dessus de 45--50 °C, les protéines se dénaturent (bactéricidie).}
\end{popfigure}

\begin{aretenir}[Facteurs clés du développement microbien -- Règle des « 3 + 1 »]
Pour qu'un microbe se développe (multiplication), il lui faut simultanément :
\begin{enumerate}
    \item \textbf{Eau disponible (Aw = activité de l'eau) :} Aw > 0,90 pour la plupart des bactéries ; > 0,70 pour les moisissures ; > 0,60 pour les levures xérophiles. \textbf{Séchage, salage, sucrage, lyophilisation} abaissent Aw $\rightarrow$ conservation.
    \item \textbf{Température adaptée :} Chaque espèce a un minimum, un optimum, un maximum.
        \begin{itemize}
            \item \textbf{Psychrophiles :} Optimum < 15 °C (eaux froides, réfrigérateur — altération lente).
            \item \textbf{Mésophiles :} Optimum 25--40 °C (majorité pathogènes humains, levures, moisissures communes). \textbf{Optimum 37 °C} = température corps humain.
            \item \textbf{Thermophiles :} Optimum 45--70 °C (sources chaudes, compost, silos).
        \end{itemize}
        \textbf{Règle pratique :} \textbf{Froid (0--4 °C) = ralentissement (bactériostase)} ; \textbf{Chauffer > 60--65 °C (pasteurisation) ou > 100 °C (stérilisation) = destruction (bactéricidie)}.
    \item \textbf{Nutriments :} Source de C, N, sels minéraux (voir Propriété 1). Milieux riches (viande, lait, œufs, plats cuisinés) = haut risque. Milieux acides (pH < 4,5), salés, sucrés = sélectifs (moisissures, levures, lactiques).
    \item \textbf{(+1) pH :} Optimum généralement neutre (6,5--7,5) pour les bactéries ; acide (4--6) pour les levures et moisissures. Acidification (vinaigre, fermentation lactique) = conservation.
\end{enumerate}
\textbf{Autres facteurs :} Pression osmotique, potentiel d'oxydo-réduction (Eh), présence d'inhibiteurs (antibiotiques naturels, épices, fumage), atmosphère (\ce{O2}, \ce{CO2}, \ce{N2}).
\end{aretenir}

\begin{savaistu}[La chaîne du froid au Cameroun : un défi de santé publique]
Au Cameroun, la rupture de la chaîne du froid (transport de la viande, du poisson, du lait, des vaccins dans des véhicules non réfrigérés, marchés sans électricité, coupures fréquentes) expose les aliments à des températures mésophiles (25--35 °C) pendant des heures. Une \emph{E. coli} ou \emph{Salmonella} divise toutes les 20--30 min à 37 °C : en 6 h, 1 bactérie $\rightarrow$ $2^{12} = 4096$ ; en 12 h $\rightarrow$ 16 millions. C'est la cause majeure des \textbf{maladies diarrhéiques} (deuxième cause de mortalité < 5 ans) et des \textbf{toxi-infections alimentaires collectives (TIAC)}. La maîtrise de la température est le levier n°1 de la sécurité alimentaire.
\end{savaistu}

\courssub{Exercice résolu : Analyse d'une expérience de fermentation}

\begin{exoresolu}[Identification du métabolisme d'une souche bactérienne]
\textbf{Énoncé :} On ensemence trois tubes de milieu de culture liquide (bouillon nutritif + indicateur de pH rouge de phénol, rouge à pH neutre, jaune en milieu acide) avec la même souche bactérienne. Les tubes sont incubés à \SI{37}{\degreeCelsius} pendant 24 h dans des conditions différentes :
\begin{itemize}
    \item Tube A : Bouchon coton (air libre).
    \item Tube B : Bouchon caoutchouc hermétique + aiguille pour prélever le gaz (anaérobiose).
    \item Tube C : Bouchon hermétique + générateur d'anaérobiose (sachet absorbant l'\ce{O2} + \ce{CO2}).
\end{itemize}
Résultats après 24 h :
\begin{itemize}
    \item Tube A : Milieu trouble, \textbf{rouge} (pH neutre/basique). Gaz prélevé : ne trouble pas l'eau de chaux.
    \item Tube B : Milieu trouble, \textbf{jaune} (pH acide). Gaz prélevé : \textbf{trouble l'eau de chaux}.
    \item Tube C : Milieu trouble, \textbf{jaune} (pH acide). Pas de gaz récupérable.
\end{itemize}
\textbf{Questions :}
\begin{enumerate}
    \item Interpréter la couleur du milieu dans chaque tube.
    \item Identifier le type de respiration dans les tubes B et C.
    \item Classer cette bactérie parmi les groupes aérobies/anaérobies. Justifier.
    \item Proposer une identification probable (genre) en lien avec les résultats.
\end{enumerate}

\tcblower
\textbf{Solution détaillée :}

\begin{enumerate}
    \item \textbf{Interprétation des couleurs :}
        \begin{itemize}
            \item \textbf{Tube A (rouge) :} pH neutre/basique. En présence d'\ce{O2}, la bactérie a effectué une \textbf{respiration aérobie complète} : oxydation totale du glucose en \ce{CO2} + \ce{H2O}. Le \ce{CO2} dissous forme de l'acide carbonique, mais l'ammoniac (\ce{NH3}) issu de la déamination d'acides aminés (peptones du bouillon) alcalinise le milieu $\rightarrow$ rouge. Pas de trouble de l'eau de chaux car le \ce{CO2} est faible ou réabsorbé, ou le gaz prélevé est de l'air résiduel.
            \item \textbf{Tubes B et C (jaunes) :} pH acide. En l'absence d'\ce{O2}, la bactérie a effectué une \textbf{fermentation} : dégradation incomplète du glucose en acides organiques (acide lactique, acétique, formique, succinique) $\rightarrow$ chute du pH $\rightarrow$ jaune.
        \end{itemize}
    \item \textbf{Type de respiration :}
        \begin{itemize}
            \item Tube B : \textbf{Fermentation mixte (gaz + acides)}. Le gaz trouble l'eau de chaux $\rightarrow$ production de \ce{CO2} (et peut-être \ce{H2}). Typique de la fermentation \emph{entérobactérienne} (forme mixte : acides + \ce{CO2} + \ce{H2} + éthanol).
            \item Tube C : \textbf{Fermentation sans production gazeuse nette} (ou gaz réabsorbé par le générateur). Production d'acides uniquement.
        \end{itemize}
    \item \textbf{Classification :} \textbf{Anaérobie facultatif (aérobie facultatif)}.
        \textbf{Justification :} Elle croît \textbf{bien en présence d'\ce{O2}} (Tube A : trouble = croissance) \textbf{ET bien en l'absence d'\ce{O2}} (Tubes B et C : trouble = croissance). Elle adapte son métabolisme : respiration aérobie si \ce{O2}, fermentation si pas d'\ce{O2}.
    \item \textbf{Identification probable :} \textbf{\emph{Escherichia coli}} (ou \emph{Klebsiella}, \emph{Enterobacter}).
        \emph{Indices :} Bâtonnet Gram négatif (sous-entendu par croissance sur bouillon standard), mésophile (\SI{37}{\degreeCelsius}), anaérobie facultatif, fermentation du glucose avec production d'acides \textbf{et} de gaz (\ce{CO2}, \ce{H2}) $\rightarrow$ test de l'eau de chaux positif (Tube B). C'est le colibacille, indicateur de contamination fécale, anaérobie facultatif type.
\end{enumerate}
\end{exoresolu}

\begin{attention}[Erreurs fréquentes à ne pas commettre au BEPC]
\begin{itemize}
    \item \textbf{Confondre « respiration » et « ventilation » :} Les microbes n'ont pas de poumons. La respiration est \textbf{cellulaire} (mitochondries ou cytoplasme). Ne pas écrire « la bactérie respire par sa membrane » mais « la bactérie réalise sa respiration cellulaire au niveau de sa membrane cytoplasmique (procaryotes) ou de ses mitochondries (levures, moisissures) ».
    \item \textbf{Oublier le rôle de l'eau :} « Les microbes ont besoin de nourriture » $\rightarrow$ \textbf{Incomplet}. Il faut citer \textbf{l'eau, les sels minéraux, la matière organique}. L'eau est souvent oubliée car « évidente », mais c'est le facteur limitant n°1 (Aw).
    \item \textbf{Dire « la fermentation produit de l'énergie » sans préciser :} C'est \textbf{peu d'énergie (2 ATP)} vs \textbf{beaucoup (30--38 ATP)} en aérobie. La fermentation est un « pis-aller » métabolique.
    \item \textbf{Mélanger « anaérobie strict » et « anaérobie facultatif » :} \emph{Clostridium} (tétanos) \textbf{meurt} à l'air. \emph{E. coli} (colibacille) \textbf{vit} à l'air. C'est une distinction vitale en clinique (prélèvement, transport, culture).
    \item \textbf{Test de l'eau de chaux :} Ne pas écrire « l'eau de chaux devient blanche ». Elle devient \textbf{trouble (laitière)} par formation d'un \textbf{précipité blanc de \ce{CaCO3}}. « Blanche » = couleur, « trouble » = aspect physique (diffusion de la lumière par les particules en suspension).
\end{itemize}
\end{attention}

\courssub{Synthèse : ce qu'il faut retenir pour le BEPC}

\begin{aretenir}[Nutrition et respiration des microbes -- Points clés BEPC]
\begin{enumerate}
    \item \textbf{Besoins :} Eau (Aw), Sels minéraux (N, P, K, Mg, S…), Matière organique (Glucides, Protéines). $\rightarrow$ Hétérotrophie majoritaire (saprophytes, parasites, symbiotes).
    \item \textbf{Mode de nutrition :} \textbf{Absorption de molécules dissoutes} (transport membranaire). Hydrolyse extracellulaire préalable par exoenzymes si macromolécules (amidon, protéines).
    \item \textbf{Respiration :} Deux types majeurs.
        \begin{itemize}
            \item \textbf{Aérobie :} \ce{O2} requis. Glucose + \ce{O2} $\rightarrow$ \ce{CO2} + \ce{H2O} + \textbf{$\approx$ 38 ATP}. (Bactéries aérobies strictes, levures + \ce{O2}, moisissures).
            \item \textbf{Anaérobie (Fermentation) :} Sans \ce{O2}. Glucose $\rightarrow$ Produits incomplets (Acide lactique \textbf{OU} Éthanol + \ce{CO2}) + \textbf{2 ATP}. (Bactéries lactiques, levures - \ce{O2}, Clostridies, Entérobactéries).
        \end{itemize}
    \item \textbf{Preuve expérimentale du \ce{CO2} :} Barbotage dans l'\textbf{eau de chaux} $\rightarrow$ \textbf{Trouble (\ce{CaCO3})}.
    \item \textbf{Groupes selon \ce{O2} :} Aérobies stricts, Anaérobies stricts, Anaérobies facultatifs (le plus courant : \emph{E. coli}, \emph{Levure}, \emph{Staphylococcus}), Microaérophiles.
    \item \textbf{Facteurs de croissance :} Température (optimum 37 °C pour mésophiles), pH (neutre pour bactéries), Aw ($>$ 0,9), Nutriments.
    \item \textbf{Conservation des aliments :} Agir sur ces facteurs : \textbf{Froid} (ralentir), \textbf{Chaleur} (pasteuriser/stériliser), \textbf{Séchage/Salage/Sucrage} (baisser Aw), \textbf{Acidification} (baisser pH), \textbf{Stérilisation/Asséptie} (éliminer).
\end{enumerate}
\end{aretenir}

\begin{exercice}[Application : Conservation du « Koki » (haricots + huile de palme + épices, cuit à la vapeur dans des feuilles de bananier)]
Le Koki est un plat traditionnel camerounais. Après cuisson à la vapeur (\SI{100}{\degreeCelsius} pendant 1--2 h), il est laissé dans ses feuilles à température ambiante (28--32 °C).
\begin{enumerate}
    \item La cuisson à la vapeur a-t-elle stérilisé le Koki ? Justifier (température, durée, spores).
    \item Quels groupes de microbes (selon \ce{O2}) risquent de se développer dans le Koki refroidi ? Pourquoi ?
    \item Le pH du Koki est proche de 6,5. L'Aw est élevée ($>$ 0,95). Prédire l'évolution de la flore microbienne en 24 h à \SI{30}{\degreeCelsius}.
    \item Proposer deux méthodes traditionnelles ou modernes pour prolonger sa conservation sans réfrigérateur, en expliquant le principe biophysique/microbiologique.
\end{enumerate}
\end{exercice}

\begin{corrige}[Exercice : Conservation du Koki]
\begin{enumerate}
    \item \textbf{Non, pas stérilisation stricte.} \SI{100}{\degreeCelsius} pendant 1--2 h détruit les formes végétatives (bactéries, levures, moisissures) mais \textbf{pas les spores bactériennes} (\emph{Bacillus}, \emph{Clostridium}) qui résistent à \SI{100}{\degreeCelsius} pendant des heures (stérilisation = \SI{121}{\degreeCelsius} / 20 min / 1 bar en autoclave). De plus, la contamination est possible à l'ouverture/manipulation.
    \item \textbf{Anaérobies facultatifs et stricts.} L'intérieur du Koki (épais, dense, riche en lipides) devient rapidement \textbf{anaérobie} (consommation de l'\ce{O2} résiduel par la flore résiduelle ou chimisme). Les anaérobies stricts (\emph{Clostridium} — risque botulisme si pH > 4,6 et Aw haute) et facultatifs (\emph{Bacillus cereus} — émétique, diarrhéique ; \emph{Staphylococcus aureus} si contamination manuelle post-cuisson) sont favorisés.
    \item \textbf{Pullulation rapide.} Conditions idéales : T° optimum mésophile (\SI{30}{\degreeCelsius}), pH neutre favorable aux bactéries, Aw max, nutriments riches (protéines haricots, lipides palme, glucides). En 24 h : charge microbienne $> 10^7$--$10^8$ UFC/g. Risque toxi-infection élevé (\emph{B. cereus}, \emph{Cl. perfringens}, \emph{S. aureus}).
    \item \textbf{Méthodes :}
        \begin{enumerate}
            \item \textbf{Acidification / Fermentation lactique contrôlée :} Ajouter du jus de citron, du vinaigre, ou laisser fermenter naturellement (si souches lactiques dominantes) pour abaisser pH $<$ 4,5 $\rightarrow$ inhibition des pathogènes (sauf moisissures/levures). Principe : \textbf{barrière pH}.
            \item \textbf{Séchage / Fumage :} Étaler en fines couches au soleil ou au-dessus du feu $\rightarrow$ baisse drastique de Aw ($<$ 0,60) $\rightarrow$ arrêt croissance bactérienne (moisissures possibles si Aw $>$ 0,70). Principe : \textbf{barrière Aw}. Le fumage ajoute des composés phénoliques antibactériens.
        \end{enumerate}
        \textit{Autre :} Réchauffage rigoureux à cœur ($>$ \SI{70}{\degreeCelsius}) avant chaque consommation (destruction formes végétatives, pas spores ni thermostables entérotoxines de \emph{S. aureus} déjà produites).
\end{enumerate}
\end{corrige}

\coursec{Reproduction des microbes}

Dans la section précédente, nous avons vu comment les microbes se nourrissent et respirent : ils puisent dans leur milieu la matière et l'énergie nécessaires à leur survie. Mais à quoi sert cette énergie ? À grandir, certes, mais surtout à \cle{se multiplier}. C'est là une des caractéristiques les plus redoutables des microorganismes : leur capacité à se reproduire à une vitesse extraordinaire. Une seule bactérie déposée sur un aliment peut donner, en quelques heures, des millions de descendants. Comprendre comment les microbes se reproduisent, c'est comprendre pourquoi les aliments se gâtent si vite, pourquoi une infection peut devenir grave en peu de temps, et pourquoi l'hygiène est si importante.

\courssub{La reproduction des bactéries : la scissiparité}

\textbf{Observation.} Un morceau de viande laissé à l'air libre, dans une cuisine chaude, se gâte en quelques heures : il change d'odeur, de couleur, de texture. Pourtant, on n'a rien vu se poser dessus ! D'où viennent ces millions de microbes ? De la multiplication, à une vitesse prodigieuse, de quelques bactéries invisibles présentes dès le départ.

\begin{definition}[La scissiparité]
Les bactéries se reproduisent par \cle{scissiparité} (ou division binaire) : une cellule mère se divise en \textbf{deux cellules filles identiques}, qui grandissent puis se divisent à leur tour. C'est un mode de reproduction \textbf{asexué} : un seul individu suffit, il n'y a ni mâle ni femelle.
\end{definition}

Le mécanisme est simple à retenir en trois étapes :
\begin{enumerate}
\item la bactérie \textbf{grandit} et son matériel génétique (son unique chromosome circulaire) se \textbf{double} ;
\item la cellule \textbf{s'allonge} et un cloisonnement apparaît en son milieu ;
\item la cellule mère se \textbf{sépare en deux} cellules filles identiques, chacune recevant un exemplaire complet du chromosome.
\end{enumerate}

\begin{propriete}[Une population qui double à chaque génération]
À chaque division, une bactérie donne deux bactéries. La population \textbf{double donc à chaque génération} : $1 \rightarrow 2 \rightarrow 4 \rightarrow 8 \rightarrow 16 \rightarrow \ldots$ Après $n$ divisions successives, une bactérie initiale a donné :
\[ N = 2^{n} \text{ bactéries} \]
On parle de \cle{croissance exponentielle}.
\end{propriete}

\begin{experience}[Rapidité de la division bactérienne]
\textbf{Observation expérimentale :} on place une bactérie (par exemple \emph{Escherichia coli}, bactérie normalement présente dans l'intestin) dans un bouillon de culture nutritif maintenu à $37~^\circ\text{C}$. On compte régulièrement le nombre de bactéries.

\textbf{Résultat :} dans ces conditions favorables (nourriture abondante, température convenable), la bactérie se divise environ \textbf{toutes les 20 minutes}.

\textbf{Interprétation :} en une heure, il y a 3 générations ($60/20 = 3$), donc $2^3 = 8$ bactéries. En 7 heures, il y a 21 générations, soit $2^{21} = \num{2097152}$ bactéries : \textbf{plus de 2 millions de bactéries} issues d'une seule !
\end{experience}

\begin{exemplebox}[Un calcul qui fait réfléchir]
Une bactérie se divisant toutes les 20 minutes donne théoriquement :
\begin{itemize}
\item après 1 h : $2^3 = 8$ bactéries ;
\item après 2 h : $2^6 = 64$ bactéries ;
\item après 4 h : $2^{12} = \num{4096}$ bactéries ;
\item après 7 h : $2^{21} \approx 2,1$ millions de bactéries ;
\item après 10 h : $2^{30} \approx 1$ milliard de bactéries !
\end{itemize}
En réalité, la nourriture finit par manquer et les déchets s'accumulent : la croissance ralentit puis s'arrête. Mais le début de la multiplication est bien exponentiel.
\end{exemplebox}

\begin{popfigure}
\begin{center}
\begin{tikzpicture}
\begin{axis}[
    width=0.85\linewidth, height=6.5cm,
    xlabel={Temps (en heures)},
    ylabel={Nombre de bactéries},
    xmin=0, xmax=5, ymin=0, ymax=4200,
    xtick={0,1,2,3,4,5},
    grid=major, grid style={popline!40},
    legend pos=north west,
    axis line style={popdark},
]
\addplot[popcurve, domain=0:5, samples=200] {2^(3*x)};
\addlegendentry{Croissance théorique ($N = 2^{n}$, 1 division / 20 min)}
\addplot[only marks, mark=*, poporange, mark size=2pt] coordinates {(0,1) (1,8) (2,64) (3,512) (4,4096)};
\addlegendentry{Effectifs calculés}
\end{axis}
\end{tikzpicture}
\end{center}
\legende{Courbe de croissance d'une population bactérienne en conditions favorables (une division toutes les 20 minutes). L'effectif double à chaque génération : la courbe monte de plus en plus vite, c'est une croissance exponentielle. Points remarquables : 1 bactérie au départ, 8 après 1 h, 64 après 2 h, 512 après 3 h, \num{4096} après 4 h.}
\end{popfigure}

\begin{methode}[Exploiter une courbe de croissance bactérienne]
Pour lire et exploiter une courbe de croissance, procède toujours dans l'ordre :
\begin{enumerate}
\item \textbf{Repérer les axes} : en abscisse, le temps (avec son unité) ; en ordonnée, l'effectif de la population.
\item \textbf{Observer l'allure de la courbe} : si elle monte de plus en plus vite, la croissance est exponentielle.
\item \textbf{Déterminer le nombre de générations} : diviser le temps écoulé par la durée d'une division (ex. : $n = \dfrac{\text{temps écoulé}}{20 \text{ min}}$).
\item \textbf{Calculer l'effectif} : appliquer $N = 2^{n}$ (en partant d'une seule bactérie).
\item \textbf{Conclure} en rédigeant une phrase complète avec le résultat et son interprétation.
\end{enumerate}
\end{methode}

\begin{exoresolu}[Calcul d'un effectif bactérien]
\textbf{Énoncé :} Une bactérie tombe sur une tranche de \emph{kwanga} (pain de manioc) laissée à l'air. Dans ces conditions, elle se divise toutes les 30 minutes. Combien de bactéries y aura-t-il, en théorie, au bout de 4 heures ?
\tcblower
\textbf{Solution détaillée :}

Étape 1 : calcul du nombre de générations.
\[ n = \frac{4 \text{ h}}{30 \text{ min}} = \frac{240 \text{ min}}{30 \text{ min}} = 8 \text{ générations} \]

Étape 2 : application de la formule $N = 2^{n}$.
\[ N = 2^{8} = 256 \text{ bactéries} \]

\textbf{Conclusion :} au bout de 4 heures, une seule bactérie a théoriquement donné 256 bactéries. Cela explique pourquoi il faut conserver les aliments au frais et ne pas les laisser longtemps à l'air libre : le froid ralentit fortement la division des bactéries.
\end{exoresolu}

\courssub{La reproduction des levures : le bourgeonnement}

Les levures sont des champignons microscopiques unicellulaires (comme la levure de bière ou la levure de boulanger, \emph{Saccharomyces}). Elles ne se divisent pas en deux parties égales comme les bactéries : elles utilisent un autre mode de reproduction asexuée.

\begin{definition}[Le bourgeonnement]
Les levures se reproduisent par \cle{bourgeonnement} : une petite excroissance, le \textbf{bourgeon}, apparaît à la surface de la cellule mère. Ce bourgeon grandit, reçoit une copie du matériel génétique, puis \textbf{se détache} pour devenir une nouvelle levure indépendante.
\end{definition}

Contrairement à la scissiparité, les deux cellules obtenues ne sont pas de même taille au départ : le bourgeon est d'abord plus petit que la cellule mère, puis il grossit. Une même levure peut produire plusieurs bourgeons successivement, et garde sur sa surface des cicatrices de bourgeonnement.

\courssub{La reproduction des moisissures : les spores}

Les moisissures (comme \emph{Penicillium} sur les oranges ou \emph{Rhizopus} sur le pain) sont des champignons filamenteux. Leur mode de reproduction est différent encore.

\begin{definition}[La reproduction par spores]
Les moisissures se reproduisent grâce à des \cle{spores} : de minuscules cellules reproductrices, très légères, produites en grand nombre au sommet des filaments de la moisissure. Les spores sont \textbf{disséminées par l'air} ; lorsqu'une spore tombe sur un milieu favorable (aliment humide, chaleur), elle \textbf{germe} et donne une nouvelle moisissure.
\end{definition}

\begin{savaistu}[Pourquoi le pain « moisit tout seul » ?]
Les spores des moisissures sont si légères qu'elles \textbf{flottent dans l'air} qui nous entoure, partout et en permanence. Il suffit donc qu'une tranche de pain soit exposée à l'air pour que des spores s'y déposent. Si le pain est un peu humide et la température douce, les spores germent : quelques jours plus tard, des taches vertes ou blanches apparaissent. Le pain n'a pas « créé » les moisissures : elles viennent de l'air ! C'est pourquoi on conserve le pain dans un sac fermé.
\end{savaistu}

\begin{popfigure}
\begin{center}
\begin{tikzpicture}[scale=0.9, every node/.style={font=\small}]
% --- Scissiparité ---
\node[font=\bfseries\small, popblue] at (2.6,3.3) {Scissiparité (bactérie)};
% étape 1
\draw[fill=popblueL, draw=popblue, thick, rounded corners=6pt] (0.6,1.6) ellipse (0.45 and 0.9);
\draw[popdark] (0.6,1.6) node[circle, draw, inner sep=1.5pt, popdark] {};
\node[below] at (0.6,0.6) {\footnotesize 1. cellule mère};
% flèche
\draw[->, very thick, poporange] (1.3,1.6) -- (1.9,1.6);
% étape 2
\draw[fill=popblueL, draw=popblue, thick, rounded corners=8pt] (2.9,1.6) ellipse (0.5 and 1.05);
\draw[popdark] (2.9,2.0) node[circle, draw, inner sep=1.5pt] {};
\draw[popdark] (2.9,1.2) node[circle, draw, inner sep=1.5pt] {};
\node[below] at (2.9,0.35) {\footnotesize 2. allongement,};
\node[below] at (2.9,0.05) {\footnotesize chromosome doublé};
% flèche
\draw[->, very thick, poporange] (3.7,1.6) -- (4.3,1.6);
% étape 3
\draw[fill=popblueL, draw=popblue, thick, rounded corners=6pt] (5.0,2.05) ellipse (0.4 and 0.55);
\draw[fill=popblueL, draw=popblue, thick, rounded corners=6pt] (5.0,1.15) ellipse (0.4 and 0.55);
\draw[popdark] (5.0,2.05) node[circle, draw, inner sep=1.5pt] {};
\draw[popdark] (5.0,1.15) node[circle, draw, inner sep=1.5pt] {};
\node[below] at (5.0,0.35) {\footnotesize 3. deux cellules};
\node[below] at (5.0,0.05) {\footnotesize filles identiques};

% --- Bourgeonnement ---
\node[font=\bfseries\small, popgreen!60!black] at (9.6,3.3) {Bourgeonnement (levure)};
% étape 1
\draw[fill=popgreenL, draw=popgreen!60!black, thick] (7.3,1.6) circle (0.55);
\draw[popdark] (7.3,1.6) node[circle, draw, inner sep=1.5pt] {};
\node[below] at (7.3,0.6) {\footnotesize 1. cellule mère};
% flèche
\draw[->, very thick, poporange] (8.1,1.6) -- (8.7,1.6);
% étape 2
\draw[fill=popgreenL, draw=popgreen!60!black, thick] (9.5,1.55) circle (0.55);
\draw[fill=popgreenL, draw=popgreen!60!black, thick] (10.15,2.05) circle (0.3);
\draw[popdark] (9.5,1.55) node[circle, draw, inner sep=1.5pt] {};
\draw[popdark] (10.15,2.05) node[circle, draw, inner sep=1.2pt] {};
\node[below] at (9.7,0.6) {\footnotesize 2. le bourgeon grandit};
% flèche
\draw[->, very thick, poporange] (10.6,1.6) -- (11.2,1.6);
% étape 3
\draw[fill=popgreenL, draw=popgreen!60!black, thick] (11.9,1.9) circle (0.5);
\draw[fill=popgreenL, draw=popgreen!60!black, thick] (12.75,1.25) circle (0.32);
\draw[popdark] (11.9,1.9) node[circle, draw, inner sep=1.5pt] {};
\draw[popdark] (12.75,1.25) node[circle, draw, inner sep=1.2pt] {};
\node[below] at (12.3,0.6) {\footnotesize 3. le bourgeon se détache};
\end{tikzpicture}
\end{center}
\legende{Comparaison de deux modes de reproduction asexuée. \textbf{À gauche :} scissiparité d'une bactérie — la cellule mère s'allonge, son chromosome (point noir) se double, puis la cellule se divise en deux cellules filles identiques. \textbf{À droite :} bourgeonnement d'une levure — un bourgeon apparaît sur la cellule mère, grandit, puis se détache ; les deux cellules sont d'abord de tailles inégales.}
\end{popfigure}

\courssub{La multiplication des virus : un cas particulier}

Les virus ne sont pas de véritables cellules : ils sont beaucoup plus petits que les bactéries et ne possèdent ni membrane cytoplasmique, ni organites, ni les « machines » nécessaires pour fabriquer leurs propres constituants. Ils sont donc \textbf{incapables de se multiplier seuls}.

\begin{definition}[Multiplication des virus]
Pour se multiplier, un virus doit \textbf{pénétrer dans une cellule vivante} (cellule humaine, animale, végétale ou bactérienne), appelée \cle{cellule hôte}. Le virus détourne alors le fonctionnement de cette cellule, qui \textbf{fabrique de nouveaux virus} à sa place. Les nouveaux virus finissent par \textbf{détruire la cellule hôte} et sont libérés : ils peuvent alors infecter d'autres cellules.
\end{definition}

On retient trois étapes :
\begin{enumerate}
\item \textbf{Pénétration} : le virus entre dans la cellule hôte ;
\item \textbf{Fabrication} : la cellule, trompée, fabrique de nombreux nouveaux virus ;
\item \textbf{Destruction} : la cellule hôte éclate et meurt, libérant les virus qui infectent les cellules voisines.
\end{enumerate}

\begin{attention}[Les virus ne se « reproduisent » pas seuls]
Erreur fréquente : dire qu'un virus « se divise » ou « se reproduit » comme une bactérie. C'est \textbf{faux} ! Un virus isolé, sur une table ou dans l'air, est totalement inactif : il ne se nourrit pas, ne respire pas, ne se multiplie pas. Seule une cellule vivante peut fabriquer de nouveaux virus, obligée par le virus qui l'a infectée. C'est pourquoi on dit que les virus sont des \textbf{parasites obligatoires} des cellules, et que les antibiotiques (qui agissent sur les bactéries) sont \textbf{sans effet sur les virus}.
\end{attention}

\courssub{Conséquences : contamination rapide des aliments et gravité des infections}

La rapidité de multiplication des microbes a des conséquences directes sur notre vie quotidienne et notre santé :

\begin{itemize}
\item \textbf{Aliments :} quelques bactéries invisibles déposées sur un aliment peuvent donner des millions de descendants en quelques heures, surtout par temps chaud. L'aliment se gâte ou devient dangereux (intoxications alimentaires). D'où les règles d'hygiène : se laver les mains, conserver les aliments au frais, bien cuire les aliments, couvrir les plats.
\item \textbf{Infections :} lorsque des microbes pénètrent dans notre organisme, ils s'y multiplient très vite si nos défenses ne les arrêtent pas. C'est pourquoi une infection mal soignée peut devenir grave en peu de temps, et pourquoi il faut consulter rapidement un centre de santé et suivre le traitement prescrit jusqu'au bout.
\end{itemize}

\begin{aretenir}[L'essentiel de la section]
\begin{itemize}
\item Les \textbf{bactéries} se reproduisent par \cle{scissiparité} : 1 cellule mère $\rightarrow$ 2 cellules filles identiques. La population double à chaque génération : après $n$ divisions, $N = 2^{n}$.
\item En conditions favorables, une bactérie peut se diviser \textbf{toutes les 20 minutes} : une seule bactérie donne plus de 2 millions de descendants en 7 heures (croissance exponentielle).
\item Les \textbf{levures} se reproduisent par \cle{bourgeonnement} : un bourgeon se forme, grandit, puis se détache.
\item Les \textbf{moisissures} se reproduisent par des \cle{spores} légères, disséminées dans l'air, qui germent sur un milieu favorable.
\item Les \textbf{virus} ne se multiplient pas seuls : ils pénètrent dans une cellule hôte qui fabrique de nouveaux virus, puis est détruite.
\item Cette multiplication très rapide explique la \textbf{contamination rapide des aliments} et la \textbf{gravité possible des infections}.
\end{itemize}
\end{aretenir}

\begin{exercice}[À toi de jouer !]
Une bactérie pathogène pénètre dans un organisme et s'y divise toutes les 20 minutes.
\begin{enumerate}
\item Combien de générations se succèdent en 3 heures ?
\item Calculer le nombre théorique de bactéries présentes au bout de 3 heures.
\item Expliquer pourquoi ce résultat justifie qu'une infection doit être traitée rapidement.
\end{enumerate}
\end{exercice}

\begin{corrige}[Exercice : multiplication d'une bactérie pathogène]
\begin{enumerate}
\item Nombre de générations en 3 heures : $n = \dfrac{3 \text{ h}}{20 \text{ min}} = \dfrac{180}{20} = 9$ générations.
\item Effectif théorique : $N = 2^{9} = 512$ bactéries.
\item En seulement 3 heures, une seule bactérie donne plus de 500 descendants, et ce nombre double encore toutes les 20 minutes. Plus on attend, plus le nombre de bactéries dans l'organisme est grand et plus l'infection est difficile à combattre. Il faut donc traiter rapidement, avant que la population de microbes ne devienne considérable.
\end{enumerate}
\end{corrige}

\coursec{La contamination par les microorganismes}

Nous avons vu dans la section précédente que les microbes savent se reproduire, souvent très vite : une bactérie peut se diviser en deux toutes les vingt minutes dans de bonnes conditions. Mais pour se multiplier \emph{dans} notre corps, encore faut-il qu'ils y \emph{entrent}. Comment un microbe présent dans l'eau d'un puits, dans l'air d'une salle de classe ou sur la peau d'une lame rouillée parvient-il à l'intérieur de notre organisme ? C'est toute la question de la \cle{contamination}, que nous allons maintenant étudier voie par voie.

\courssub{Qu'est-ce que la contamination ?}

\courssub{Observation et définition}

En période de pluies, dans certains quartiers de Yaoundé ou de Douala, plusieurs membres d'une même famille peuvent se mettre à souffrir de diarrhées violentes après avoir bu l'eau d'un puits mal protégé. Quelques jours plus tard, leurs voisins qui partagent la même source d'eau tombent malades à leur tour. Les microbes présents dans l'eau souillée sont \emph{entrés} dans leur organisme et s'y sont \emph{multipliés}, provoquant la maladie. Ce phénomène porte un nom précis.

\begin{definition}[La contamination]
La \cle{contamination} est la \textbf{pénétration} de microorganismes dans l'organisme, suivie de leur \textbf{multiplication} dans celui-ci. Cette multiplication peut provoquer une \textbf{maladie} (on parle alors de maladie infectieuse) lorsque les microbes perturbent le fonctionnement de l'organisme.
\end{definition}

\begin{attention}[Contamination ne veut pas dire maladie automatique]
Être contaminé ne signifie pas toujours être malade. Notre organisme possède des défenses naturelles qui peuvent détruire les microbes avant qu'ils ne se multiplient. La contamination ne débouche sur la maladie que si les microbes franchissent ces défenses et se développent en grand nombre.
\end{attention}

\courssub{Les barrières naturelles de l'organisme}

Notre corps n'est pas une porte ouverte : il est protégé par des \cle{barrières naturelles} qui empêchent la plupart des microbes d'entrer.

\begin{itemize}
\item La \textbf{peau} : intacte, elle forme une enveloppe imperméable aux microbes. C'est notre première ligne de défense.
\item Les \textbf{muqueuses} : elles tapissent l'intérieur du nez, de la bouche, des bronches et de l'intestin ; elles produisent du mucus qui piège les microbes et les poussières.
\item L'\textbf{acidité gastrique} : le suc gastrique de l'estomac est très acide (pH $\approx$ 1,5--2) et détruit la plupart des microbes avalés avec les aliments.
\end{itemize}

\begin{aretenir}[Savoir admis]
La peau, les muqueuses et l'acidité de l'estomac constituent les barrières naturelles de l'organisme. Les microbes ne peuvent nous contaminer que s'ils franchissent ces barrières, par des \textbf{voies de pénétration} particulières.
\end{aretenir}

\courssub{Les voies de pénétration des microbes dans l'organisme}

\begin{propriete}[Les voies de contamination]
Les microbes pénètrent dans l'organisme par des voies variées : la \cle{voie digestive}, la \cle{voie respiratoire}, la \cle{voie cutanée}, la \cle{voie génitale et sanguine}, et la \cle{voie transplacentaire} (de la mère au fœtus).
\end{propriete}

\courssub{La voie digestive : les microbes entrent par la bouche}

\textbf{Observation.} Après de fortes pluies, les eaux de ruissellement mélangent les excréments aux eaux des puits et des marigots. Les personnes qui boivent cette eau sans la traiter développent des diarrhées, parfois graves. C'est le scénario classique des épidémies de choléra qui ont frappé l'Extrême-Nord du Cameroun à plusieurs reprises.

\textbf{Mécanisme.} Les microbes présents dans les \textbf{eaux et les aliments souillés}, ou transportés par les \textbf{mains sales} portées à la bouche, franchissent la bouche puis l'intestin. Ceux qui résistent à l'acidité de l'estomac se multiplient dans le tube digestif et provoquent la maladie.

\textbf{Maladies transmises par voie digestive :}
\begin{itemize}
\item la \textbf{typhoïde} (bactérie \emph{Salmonella Typhi} présente dans l'eau et les aliments souillés) ;
\item le \textbf{choléra} (bactérie \emph{Vibrio cholerae} transmise par l'eau contaminée, provoquant des diarrhées très abondantes, « en eau de riz ») ;
\item les \textbf{diarrhées} infectieuses diverses, fréquentes chez les jeunes enfants.
\end{itemize}

\textbf{Prévention :} boire de l'eau potable (ou la faire bouillir au moins 1 minute), bien laver les mains au savon avant les repas et après être allé aux toilettes, laver les fruits et légumes, bien cuire les aliments.

\courssub{La voie respiratoire : les microbes entrent par le nez et la bouche}

\textbf{Observation.} En classe, quand un élève enrhumé tousse ou éternue sans se couvrir la bouche, plusieurs de ses camarades tombent malades dans les jours qui suivent. Comment le microbe a-t-il voyagé ?

\textbf{Mécanisme.} En toussant, en éternuant ou simplement en parlant, un malade projette dans l'air des \textbf{gouttelettes de salive} microscopiques (aérosols) chargées de microbes. Ces gouttelettes, ainsi que les \textbf{poussières} contaminées, sont inhalées par les personnes proches : les microbes pénètrent alors par le \textbf{nez} et la \textbf{bouche} et se multiplient dans les voies respiratoires.

\textbf{Maladies transmises par voie respiratoire :}
\begin{itemize}
\item la \textbf{grippe} (virus \emph{Influenza}) ;
\item la \textbf{tuberculose} (bactérie \emph{Mycobacterium tuberculosis} qui atteint surtout les poumons) ;
\item la \textbf{rougeole} (virus très contagieux, évitable par la vaccination).
\end{itemize}

\textbf{Prévention :} se couvrir la bouche en toussant ou en éternuant (dans le coude ou un mouchoir), aérer les salles de classe et les chambres plusieurs fois par jour, éviter de cracher par terre, se faire vacciner quand un vaccin existe (rougeole, grippe, tuberculose BCG).

\courssub{La voie cutanée : les microbes entrent par les plaies et les piqûres}

\textbf{Observation.} Un élève marche pieds nus et se blesse avec un clou rouillé planté dans une planche. Quelques jours plus tard, sa plaie s'infecte. Autre cas : une personne piquée la nuit par un moustique développe une forte fièvre une semaine après.

\textbf{Mécanisme.} La peau intacte est une barrière, mais toute \textbf{plaie} (coupure, éraflure, brûlure) est une porte d'entrée. Les \textbf{piqûres} d'animaux (moustiques, tiques, morsures) enfoncent également la barrière cutanée et peuvent inoculer directement des microbes dans le sang ou les tissus.

\textbf{Maladies transmises par voie cutanée :}
\begin{itemize}
\item le \textbf{tétanos} : la bactérie \emph{Clostridium tetani}, présente dans la terre et la rouille, pénètre par une \textbf{plaie souillée} ; elle sécrète une toxine qui provoque des contractures musculaires très graves. Le vaccin antitétanique (souvent associé DTCoq) protège efficacement ;
\item le \textbf{paludisme} : le microbe responsable (le \emph{Plasmodium}, un parasite unicellulaire) est inoculé dans le sang par la \textbf{piqûre d'un moustique femelle}, l'anophèle.
\end{itemize}

\begin{attention}[Erreur fréquente : le paludisme et l'eau sale]
Beaucoup d'élèves croient que le paludisme s'attrape en buvant de l'eau sale. C'est \textbf{faux} ! Le paludisme est dû à un microbe (le \emph{Plasmodium}) inoculé par la \textbf{piqûre d'anophèle} : c'est une contamination par \textbf{voie cutanée et sanguine}, et non par voie digestive. L'eau sale transmet le choléra ou la typhoïde, jamais le paludisme.
\end{attention}

\textbf{Prévention :} nettoyer et désinfecter toute plaie (asepsie : eau, savon, antiseptique), se faire vacciner contre le tétanos (rappels tous les 10--20 ans), dormir sous \textbf{moustiquaire imprégnée d'insecticide} (MII), éliminer les eaux stagnantes où les moustiques pondent (gîtes larvaires).

\courssub{La voie génitale et sanguine}

\textbf{Mécanisme.} Certains microbes pénètrent dans l'organisme :
\begin{itemize}
\item lors de \textbf{rapports sexuels} non protégés, par les muqueuses génitales ;
\item par du \textbf{sang contaminé} : transfusion de sang non contrôlé, partage d'objets piquants ou tranchants souillés (lames de rasoir, aiguilles, ciseaux de coiffure, instruments de circoncision ou de scarification non stérilisés).
\end{itemize}

\textbf{Maladies transmises par ces voies :}
\begin{itemize}
\item le \textbf{VIH/Sida} : le virus (VIH) se transmet par les rapports sexuels non protégés, par le sang contaminé et de la mère à l'enfant ;
\item les \textbf{hépatites} (notamment l'hépatite B et C), transmises par le sang et les rapports sexuels.
\end{itemize}

\textbf{Prévention :} utiliser des \textbf{préservatifs} (masculins ou féminins) à chaque rapport sexuel, ne jamais partager les objets piquants ou tranchants (rasoirs, aiguilles, brosses à dents), exiger du matériel stérilisé ou à usage unique pour tout soin/injection/piercing/tatouage, contrôler le sang avant toute transfusion.

\courssub{La voie transplacentaire : de la mère au fœtus}

\begin{aretenir}[Complément utile]
Certains microbes peuvent passer de la mère au \textbf{fœtus} à travers le \textbf{placenta} pendant la grossesse : c'est la \cle{voie transplacentaire}. C'est ainsi que le VIH, la syphilis, la toxoplasmose ou la rubéole peuvent être transmis de la mère à l'enfant. Le suivi médical de la grossesse (consultations prénatales, dépistages, traitements préventifs) permet de réduire fortement ce risque.
\end{aretenir}

\courssub{Schéma de synthèse : les voies de pénétration dans le corps humain}

\begin{popfigure}
\begin{center}
\begin{tikzpicture}[scale=1.0, every node/.style={font=\small}]
% Corps humain stylisé
\draw[fill=popcream, draw=popdark, thick] (0,3.4) circle (0.55); % tête
\draw[fill=popcream, draw=popdark, thick, rounded corners=6pt] (-0.75,-0.2) rectangle (0.75,2.9); % tronc
\draw[fill=popcream, draw=popdark, thick, rounded corners=4pt] (-1.6,2.6) rectangle (-0.75,0.6); % bras G
\draw[fill=popcream, draw=popdark, thick, rounded corners=4pt] (0.75,2.6) rectangle (1.6,0.6); % bras D
\draw[fill=popcream, draw=popdark, thick, rounded corners=4pt] (-0.65,-0.2) rectangle (-0.08,-2.6); % jambe G
\draw[fill=popcream, draw=popdark, thick, rounded corners=4pt] (0.08,-0.2) rectangle (0.65,-2.6); % jambe D
% Points d'entrée
\fill[popink] (0.18,3.25) circle (0.09); % bouche
\fill[popink] (0,3.5) circle (0.07); % nez
\fill[popink] (-1.2,0.75) circle (0.09); % plaie main
\fill[popink] (0,-0.1) circle (0.09); % organes génitaux
% Flèches voies
\draw[->, very thick, popblue] (-4.4,4.1) -- (-0.6,3.55) node[midway, above, sloped, text=popblue]{\textbf{Voie respiratoire}};
\node[text=popblue, align=center, font=\footnotesize] at (-4.4,4.6) {nez et bouche\\ (grippe, tuberculose)};
\draw[->, very thick, poporange] (-4.4,2.6) -- (-0.35,3.2) node[midway, below, sloped, text=poporange]{\textbf{Voie digestive}};
\node[text=poporange, align=center, font=\footnotesize] at (-4.4,2.0) {bouche : eau, aliments\\ (choléra, typhoïde)};
\draw[->, very thick, popgreen] (-4.4,0.6) -- (-1.35,0.72) node[midway, above, sloped, text=popgreen]{\textbf{Voie cutanée}};
\node[text=popgreen, align=center, font=\footnotesize] at (-4.4,0.0) {plaies, piqûres\\ (tétanos, paludisme)};
\draw[->, very thick, poppurple] (4.4,0.2) -- (0.15,-0.05) node[midway, above, sloped, text=poppurple]{\textbf{Voie génitale et sanguine}};
\node[text=poppurple, align=center, font=\footnotesize] at (4.4,-0.5) {rapports sexuels, sang\\ (VIH/Sida, hépatites)};
\draw[->, very thick, poppink] (4.4,-2.2) -- (0.6,-1.4) node[midway, below, sloped, text=poppink]{\textbf{Voie transplacentaire}};
\node[text=poppink, align=center, font=\footnotesize] at (4.4,-2.9) {de la mère au fœtus\\ (grossesse)};
\end{tikzpicture}
\end{center}
\legende{Schéma des voies de pénétration des microbes dans l'organisme humain. Les flèches indiquent les portes d'entrée : nez et bouche (voie respiratoire), bouche (voie digestive), plaies et piqûres de la peau (voie cutanée), organes génitaux et sang (voie génitale et sanguine), placenta (voie transplacentaire).}
\end{popfigure}

\courssub{Tableau de synthèse : voies, maladies et prévention}

\begin{center}
\small
\begin{tabularx}{\linewidth}{|l|X|X|}
\hline
\rowcolor{popblueL} \textbf{Voie de contamination} & \textbf{Microbes / maladies} & \textbf{Moyens de prévention} \\
\hline
Digestive (bouche) & Choléra, typhoïde, diarrhées & Eau potable ou bouillie, lavage des mains, cuisson des aliments \\
\hline
Respiratoire (nez, bouche) & Grippe, tuberculose, rougeole & Se couvrir la bouche, aérer les pièces, vaccination \\
\hline
Cutanée (plaies, piqûres) & Tétanos, paludisme & Désinfecter les plaies, vaccin antitétanique, moustiquaires imprégnées \\
\hline
Génitale et sanguine & VIH/Sida, hépatites B, C & Préservatifs, objets piquants personnels, sang contrôlé \\
\hline
Transplacentaire & VIH, syphilis, rubéole, toxoplasmose & Suivi médical de la grossesse (CPN) \\
\hline
\end{tabularx}
\end{center}

\courssub{Application à la vie courante : bloquer chaque voie de contamination}

\begin{methode}[Comment prévenir une contamination]
Pour prévenir une contamination, on raisonne en deux étapes :
\begin{enumerate}
\item \textbf{Identifier la voie de pénétration} du microbe responsable de la maladie (digestive, respiratoire, cutanée, génitale et sanguine, transplacentaire) ;
\item \textbf{Appliquer la règle d'hygiène qui bloque cette voie} : lavage des mains et eau potable pour la voie digestive, aération et couverture de la bouche pour la voie respiratoire, asepsie des plaies et moustiquaires pour la voie cutanée, préservatifs et matériel stérilisé pour la voie génitale et sanguine.
\end{enumerate}
\end{methode}

\begin{exemplebox}[L'épidémie de choléra : un exemple de voie digestive]
Lors d'une épidémie de choléra, la bactérie responsable (\emph{Vibrio cholerae}) se propage par l'\textbf{eau et les aliments souillés} par les excréments des malades. La voie de pénétration est la \textbf{voie digestive}. Pour bloquer la propagation, il suffit de fermer cette porte : \textbf{faire bouillir l'eau} de boisson, \textbf{laver les mains au savon} après être allé aux toilettes et avant de préparer les repas, et bien laver les aliments consommés crus. Ces gestes simples, appliqués par toute la communauté, peuvent stopper l'épidémie.
\end{exemplebox}

\begin{savaistu}[Le pouvoir d'une poignée de main]
Une simple poignée de main peut transmettre des \textbf{milliers de microbes} ! Nos mains touchent chaque jour des centaines de surfaces : poignées de porte, argent, téléphones, rampes d'escalier... C'est pourquoi le \textbf{lavage régulier des mains au savon} est considéré par les médecins comme \textbf{le geste barrière le plus efficace et le moins coûteux} pour lutter contre les contaminations. Au Cameroun, les campagnes de lavage des mains (notamment dans les écoles via le programme « École amie de l'hygiène ») ont permis de limiter la propagation de nombreuses épidémies diarrhéiques.
\end{savaistu}

\begin{exoresolu}[Identifier une voie de contamination et proposer une prévention]
\textbf{Énoncé.} Dans un village, plusieurs enfants qui jouent pieds nus dans la cour se blessent avec des objets rouillés. L'un d'eux est hospitalisé : les médecins diagnostiquent un tétanos. (a) Quelle est la voie de pénétration du microbe responsable ? (b) Cite deux mesures qui auraient pu empêcher cette contamination.
\tcblower
\textbf{Solution détaillée.}

(a) Le tétanos est provoqué par une bactérie (\emph{Clostridium tetani}) présente dans la terre et sur les objets rouillés souillés. Elle pénètre dans l'organisme par une \textbf{plaie de la peau} : il s'agit donc de la \textbf{voie cutanée}.

(b) Deux mesures de prévention :
\begin{itemize}
\item \textbf{nettoyer et désinfecter immédiatement toute plaie} (asepsie) avec de l'eau propre, du savon et un antiseptique (ex. : Bétadine\textsuperscript{\textregistered}, Dakin\textsuperscript{\textregistered}) ;
\item \textbf{se faire vacciner contre le tétanos}, vaccin (souvent DTCoq) qui protège durablement l'organisme par la production d'anticorps anti-toxine.
\end{itemize}
On peut aussi conseiller de porter des chaussures pour éviter les blessures, ce qui revient à maintenir intacte la barrière naturelle qu'est la peau.
\end{exoresolu}

\begin{attention}[Pièges de rédaction au BEPC]
\begin{itemize}
\item Ne confonds pas \textbf{voie de pénétration} et \textbf{maladie} : la question « par quelle voie ? » attend une réponse comme « voie digestive », pas « le choléra ».
\item Le \textbf{paludisme} se classe dans la voie \textbf{cutanée et sanguine} (piqûre d'anophèle), jamais dans la voie digestive.
\item La contamination comporte \textbf{deux étapes} : la pénétration \emph{puis} la multiplication du microbe. N'oublie pas la seconde dans ta définition.
\item « Voie sanguine » seule est incomplet pour le VIH/Sida ou les hépatites : on précise \textbf{voie génitale et sanguine} (ou « voie sexuelle et sanguine »).
\end{itemize}
\end{attention}

\begin{exercice}[Pour t'entraîner]
Recopie et complète le tableau suivant : pour chaque maladie, indique la voie de contamination et une mesure de prévention : grippe ; choléra ; paludisme ; VIH/Sida ; tétanos.
\end{exercice}

\begin{corrige}[Exercice]
Grippe : voie respiratoire ; se couvrir la bouche en toussant, aérer les pièces, vaccination. Choléra : voie digestive ; faire bouillir l'eau, laver les mains, cuire les aliments. Paludisme : voie cutanée et sanguine (piqûre d'anophèle) ; dormir sous moustiquaire imprégnée, éliminer les gîtes larvaires. VIH/Sida : voie génitale et sanguine ; préservatifs, objets piquants personnels, sang contrôlé. Tétanos : voie cutanée (plaie souillée) ; désinfecter les plaies, vaccination (rappels).
\end{corrige}

\begin{aretenir}[L'essentiel de la section]
La contamination est la pénétration puis la multiplication de microbes dans l'organisme. Malgré les barrières naturelles (peau, muqueuses, acidité gastrique), les microbes peuvent entrer par cinq voies : \textbf{digestive} (eau et aliments souillés), \textbf{respiratoire} (gouttelettes et poussières), \textbf{cutanée} (plaies et piqûres), \textbf{génitale et sanguine} (rapports sexuels, sang contaminé) et \textbf{transplacentaire} (mère-fœtus). Prévenir une maladie infectieuse, c'est identifier sa voie de contamination et appliquer la règle d'hygiène qui la bloque.
\end{aretenir}

\newpage

\coursec{Activité d'intégration}

\courssub{Objectifs de l'activité}
\begin{itemize}
    \item Mobiliser les connaissances sur les \cle{groupes de microorganismes} (bactéries, champignons microscopiques), leurs \cle{conditions de développement} (nutrition, température) et leurs \cle{modes de contamination} pour résoudre un problème concret.
    \item Mettre en œuvre une \cle{démarche d'investigation scientifique} : observer des résultats expérimentaux (culture sur gélose, courbe de croissance), formuler des hypothèses, les tester par l'interprétation des données, et conclure de manière argumentée.
    \item Proposer des \cle{mesures de prévention} concrètes et adaptées au contexte d'une cantine scolaire camerounaise pour rompre la chaîne de contamination.
\end{itemize}

\courssub{Situation : Enquête sur une toxi-infection alimentaire collective au Collège Bilingue de Douala-Bonabéri}
\label{sit:enquete}

Lundi 14 octobre 2024, 07 h 30. Au Collège Bilingue de Douala-Bonabéri, une cinquantaine d'élèves de classes de 3ème et de 4ème se présentent à l'infirmerie scolaire dans l'heure qui suit la récréation de 10 h. Ils présentent tous des symptômes identiques : \textbf{nausées violentes, vomissements répétés, douleurs abdominales intenses et diarrhée aqueuse}. Aucun cas de fièvre n'est signalé dans les premiers temps. L'infirmière, M.~\textsc{Ngoua}, note une apparition brutale des symptômes (incubation courte, 2 à 4 h après l'ingestion) et suspecte immédiatement une \cle{toxi-infection alimentaire collective} (TIAC).

L'enquête alimentaire rapide révèle que tous les élèves malades ont consommé des \textbf{beignets de haricots} (spécialité locale très prisée) achetés à la cantine de l'établissement vers 09 h 45. Les élèves ayant mangé d'autres aliments (pain, banane, arachides) ou ayant apporté leur repas ne sont pas touchés.

Le Chef d'établissement mandate le Club Sciences et Vie de la Terre (SVTEEHB), encadré par Mme~\textsc{Fotso}, pour mener l'investigation scientifique en collaboration avec le service d'hygiène du district de santé. Voici les éléments recueillis :

\begin{enumerate}
    \item \textbf{Témoignage de la vendeuse (Mme~\textsc{Mbia}) :} « J'ai préparé la pâte (haricots trempés, mixés, assaisonnés avec oignon, piment, sel, cube) \textbf{hier dimanche vers 16 h}. Je l'ai laissée dans une grande bassine en plastique, \textbf{couverte d'un simple tissu}, sur la table de la cuisine extérieure. Ce matin, vers 06 h 30, j'ai fait chauffer l'huile dans la grande marmite. Je n'ai \textbf{pas eu le temps de me laver les mains} avec du savon après avoir balayé la cour et arrangé le bois de chauffe ; je les ai juste essuyées sur mon pagne. J'ai façonné les beignets à la main et les ai frits par fournées jusqu'à 09 h 30. Les derniers beignets sont restés \textbf{à l'air libre dans une corbeille en osier} jusqu'à la vente à 09 h 45. Il faisait très chaud hier et ce matin (environ \SI{32}{\celsius} d'après la météo locale). »
    
    \item \textbf{Analyse microbiologique (réalisée au Centre Pasteur du Cameroun, antenne de Douala, sur prélèvement effectué à 11 h 00) :}
    Un échantillon de beignet restant (prélevé dans la corbeille) et un beignet témoin (fraîchement refrit à 11 h 15 à partir de la même pâte restante mais cuit à cœur) ont été mis en culture sur gélose nutritive standard (gélose Chapman pour \textit{Staphylococcus} et gélose standard pour flore totale) à \SI{37}{\celsius} pendant \SI{24}{\hour}.
    
    \begin{center}
    \begin{tabularx}{\linewidth}{|l|X|X|}\hline
    \rowcolor{popblueL}
    \textbf{Échantillon} & \textbf{Résultat sur Gélose Chapman (spécifique \textit{Staphylococcus aureus})} & \textbf{Résultat sur Gélose Nutritive (Flore totale aérobie)} \\ \hline
    Beignet \textbf{incriminé} (resté à l'air libre 1h30) & \textbf{Nombreuses colonies} jaunes dorées (\textgreater \num{10^5} UFC/g) & \textbf{Nombreuses colonies} de morphologies variées (blanches, crème, pigmentées) (\textgreater \num{10^6} UFC/g) \\ \hline
    Beignet \textbf{témoin} (fraîchement cuit à cœur) & \textbf{Aucune colonie} & \textbf{Aucune colonie} \\ \hline
    \end{tabularx}
    \end{center}
    \textit{UFC/g : Unités Formant Colonies par gramme. Seuil d'alerte pour \textit{S. aureus} dans les aliments : \num{10^3} UFC/g.}
    
    \item \textbf{Données bibliographiques : Courbe de croissance bactérienne type (ex. \textit{Staphylococcus aureus}) en fonction de la température de conservation.}
    
    \begin{popfigure}
    \begin{tikzpicture}
        \begin{axis}[
            width=\linewidth, height=8cm,
            xlabel={Température (\si{\celsius})}, ylabel={Vitesse de croissance relative},
            xmin=0, xmax=45, ymin=0, ymax=1.2,
            domain=0:45, samples=200,
            axis lines=middle, tick style={popdark},
            xlabel style={anchor=north}, ylabel style={anchor=east},
            grid=major, grid style={popline},
        ]
        % Courbe type : croissance nulle < 10, optimale ~37, nulle > 45
        \addplot[popcurve, very thick, poporange] { 
            (x<10 || x>45) * 0 + 
            (x>=10 && x<=45) * ( 
                (x<=37) * ( (x-10)/(37-10) )^2 * 1.0 + 
                (x>37) * ( (45-x)/(45-37) )^2 * 1.0 
            )
        };
        % Zones de fond (rectangles)
        \fill[popgreen!30] (axis cs:0,0) rectangle (axis cs:10,1.2);
        \fill[poporange!30] (axis cs:10,0) rectangle (axis cs:45,1.2);
        \fill[popred!30] (axis cs:45,0) rectangle (axis cs:45,1.2); % ligne fine à 45
        % Annotations
        \node[popgreen!80!black, align=center] at (axis cs:5,0.2) {Zone de\\ non croissance};
        \node[poporange!80!black, align=center] at (axis cs:27,1.1) {Zone de croissance\\ optimale (mésophiles)};
        \node[popred!80!black, align=center] at (axis cs:42,0.2) {Zone de\\ destruction / arrêt};
        % Point température ambiante Douala
        \draw[popblue, dashed, thick] (axis cs:32,0) -- (axis cs:32,0.95) node[above, popblue] {$T_{\text{ambiante}} \approx \SI{32}{\celsius}$};
        \draw[popblue, dashed, thick] (axis cs:32,0.95) -- (axis cs:0,0.95) node[left, popblue] {$v_{\text{rel}} \approx 0.95$};
        \end{axis}
    \end{tikzpicture}
    \legende{Courbe théorique de vitesse de croissance relative des bactéries mésophiles (ex. \textit{Staphylococcus aureus}) en fonction de la température. La température ambiante relevée (\SI{32}{\celsius}) se situe en zone de croissance rapide.}
    \end{popfigure}
\end{enumerate}

\courssub{Consignes}
En vous appuyant \textbf{uniquement} sur les documents fournis et vos connaissances du cours (Leçon 04), répondez aux questions suivantes en justifiant \textbf{systématiquement} chaque affirmation par une référence au document (Doc 1, Doc 2, Doc 3) ou à une notion du cours.

\begin{enumerate}
    \item \textbf{Identification de l'agent pathogène :} D'après les résultats de culture (Doc 2), quel est le \cle{genre} bactérien majoritaire identifié sur le milieu sélectif Chapman ? Quelle est sa \cle{forme} microscopique et son \cle{mode de groupement} caractéristique (rappel de cours) ? Ce germe est-il un virus, un champignon ou une bactérie ? Justifiez.
    
    \item \textbf{Interprétation du témoin :} Pourquoi le beignet témoin (fraîchement cuit) ne présente-t-il \textbf{aucune colonie} sur les deux milieux ? Quelle propriété de la \cle{cuisson} (température) explique ce résultat ? Quel lien faites-vous avec le mode de nutrition des bactéries ?
    
    \item \textbf{Analyse de la multiplication bactérienne :} 
    \begin{enumerate}
        \item En exploitant la courbe du Doc 3, précisez la \cle{vitesse de croissance relative} des bactéries mésophiles à \SI{32}{\celsius} (température ambiante de la cuisine). 
        \item Sachant qu'une division bactérienne (génération) dure environ \SI{30}{\minute} dans des conditions optimales (\SI{37}{\celsius}) et que la vitesse relative à \SI{32}{\celsius} est d'environ \num{0,95}, calculez le \textbf{nombre de générations} survenues pendant les \SI{17}{\hour} de conservation de la pâte (de 16 h dimanche à 09 h lendemain) et pendant les \SI{1,5}{\hour} de conservation des beignets cuits (de 09 h 30 à 11 h 00). 
        \item En déduire le \textbf{facteur de multiplication} total (nombre de bactéries finales / nombre initial) pour l'ensemble de la conservation (pâte crue + beignets cuits). (On suppose une contamination initiale faible, ex. \num{10} bactéries/g). Concluez sur le risque.
    \end{enumerate}
    
    \item \textbf{Voies de contamination :} En croisant le témoignage (Doc 1) et les résultats (Doc 2), identifiez \textbf{deux voies de contamination} distinctes ayant conduit à la présence massive de \textit{Staphylococcus aureus} dans les beignets incriminés. Pour chaque voie, citez l'élément du témoignage qui la prouve et le \cle{mode de vie} du microbe concerné (résistance, sporulation, production de toxines).
    
    \item \textbf{Mécanisme pathogène :} Les symptômes (vomissements brutaux, pas de fièvre, incubation courte) sont-ils dus à la bactérie elle-même (infection) ou à une \cle{toxine} préformée (intoxication) ? Justifiez votre réponse par la physiologie bactérienne (conditions de production de l'entérotoxine staphylococcique) et le profil clinique.
    
    \item \textbf{Conclusion et Prévention :} Rédigez une \textbf{conclusion synthétique} (5-6 lignes) désignant la cause probable de l'épidémie. Proposez ensuite \textbf{trois mesures d'hygiène concrètes, applicables immédiatement} dans cette cantine, en les reliant chacune à une étape de la chaîne de contamination identifiée (préparation, cuisson, conservation, service).
\end{enumerate}

\courssub{Ressources à mobiliser}
\begin{itemize}
    \item \textbf{Groupes de microorganismes :} Bactéries (procaryotes, unicellulaires, fission binaire), Virus (acellulaires, parasites obligatoires), Champignons microscopiques (eucaryotes, levures/moisissures, bourgeonnement/spores).
    \item \textbf{Nutrition/Respiration :} Bactéries hétérotrophes (saprophytes ou parasites), aérobies / anaérobies / facultatives. \textit{S. aureus} : anaérobie facultatif, hémolysant (sur gélose sang), coagulase +.
    \item \textbf{Reproduction :} Division binaire (scissiparité). Temps de génération $g$. $N_t = N_0 \times 2^n$ avec $n = \frac{t}{g}$.
    \item \textbf{Facteurs de croissance :} Température (cardinales : min, opt, max), pH, $a_w$ (activité de l'eau), nutriments. Mésophiles : opt \SI{37}{\celsius}, croissance entre \SI{10}{\celsius} et \SI{45}{\celsius}.
    \item \textbf{Contamination :} Voies : digestive (ingestion), respiratoire, cutanée-muqueuse. Portage sain (\textit{S. aureus} : narines, peau, mains). Toxi-infection vs Intoxication (toxines thermostables).
    \item \textbf{Lutte :} Stérilisation (chaleur humide \SI{121}{\celsius}), Pasteurisation, Réfrigération (\textless \SI{4}{\celsius} bloque croissance), Hygiène des mains (lavage savon/eau / solution hydroalcoolique), Chaîne du froid, Cuisson à cœur (\SI{70}{\celsius} cœur).
\end{itemize}

\courssub{Démarche et corrigé commenté}
\begin{exoresolu}[Corrigé détaillé de l'enquête]
\textbf{1. Identification de l'agent pathogène}
\begin{itemize}
    \item \textbf{Genre :} \textit{Staphylococcus} (identifié par les colonies \textbf{jaunes dorées} sur gélose Chapman, milieu sélectif pour \textit{Staphylococcus aureus} qui fermente le mannitol $\rightarrow$ acidification $\rightarrow$ virage du rouge au jaune).
    \item \textbf{Morphologie / Groupement :} Bactéries à \textbf{forme sphérique (cocques)} Gram +, groupées en \textbf{grappes} irrégulières (observation au microscope optique, TP N°5).
    \item \textbf{Nature :} C'est une \textbf{bactérie} (procaryote, unicellulaire, se reproduit par division binaire, cultivable sur milieu inerte). Ce n'est pas un virus (non cultivable sur gélose, parasite intracellulaire strict) ni un champignon (eucaryote, colonies duveteuses, reproduction par bourgeonnement/spores).
\end{itemize}

\textbf{2. Interprétation du témoin}
\begin{itemize}
    \item \textbf{Absence de colonies :} La cuisson à cœur du beignet témoin a atteint une température \textbf{supérieure à \SI{70}{\celsius}} (température de l'huile de friture \textasciitilde \SI{180}{\celsius}). 
    \item \textbf{Explication :} La chaleur détruit les structures protéiques (enzymes, membrane, ADN) des bactéries mésophiles \textit{Staphylococcus aureus} (thermosensibles, $T_{\text{max}} \approx \SI{45}{\celsius}$ d'après Doc 3). C'est une \textbf{stérilisation de surface et à cœur}.
    \item \textbf{Lien nutrition :} Les bactéries sont des êtres vivants \cle{hétérotrophes aérobies/anaérobies facultatifs} ; elles ont besoin d'un métabolisme actif (respiration/fermentation) pour se multiplier. La chaleur dénature les enzymes indispensables à ce métabolisme $\rightarrow$ mort cellulaire $\rightarrow$ aucune UFC.
\end{itemize}

\textbf{3. Analyse de la multiplication bactérienne}
\begin{enumerate}
    \item \textbf{Vitesse à \SI{32}{\celsius} :} Lecture graphique (Doc 3) : $v_{\text{rel}} \approx \num{0,95}$ (proche de l'optimum \SI{37}{\celsius}).
    \item \textbf{Calcul des générations :}
    \begin{itemize}
        \item Temps de génération réel à \SI{32}{\celsius} : $g_{32} = \frac{g_{37}}{v_{\text{rel}}} = \frac{\SI{30}{\minute}}{\num{0,95}} \approx \SI{31,6}{\minute} \approx \SI{0,527}{\hour}$.
        \item Période 1 (Pâte crue) : $\Delta t_1 = \SI{17}{\hour}$ (16 h $\rightarrow$ 09 h). 
        $n_1 = \frac{\Delta t_1}{g_{32}} = \frac{17}{0,527} \approx \textbf{32,3 générations}$.
        \item Période 2 (Beignets cuits à l'air libre) : $\Delta t_2 = \SI{1,5}{\hour}$ (09 h 30 $\rightarrow$ 11 h 00).
        $n_2 = \frac{1,5}{0,527} \approx \textbf{2,8 générations}$.
    \end{itemize}
    \item \textbf{Facteur de multiplication total (Période 1 + 2) :}
    $n_{\text{total}} = n_1 + n_2 \approx 35,1$ générations.
    Facteur $F = 2^{n_{\text{total}}} = 2^{35,1} \approx \mathbf{3,5 \times 10^{10}}$.
    \textbf{Population finale estimée :} $N_f = N_0 \times F = 10 \times 3,5 \times 10^{10} = \mathbf{3,5 \times 10^{11} \text{ UFC/g}}$.
    \textbf{Conclusion :} Ce chiffre (\textgreater \num{10^{11}} UFC/g) est \textbf{très largement supérieur} au seuil d'alerte (\num{10^3} UFC/g) et aux résultats d'analyse (\textgreater \num{10^5} UFC/g). La température ambiante (\SI{32}{\celsius}) a permis une \textbf{prolifération exponentielle explosive} pendant les 17 h de conservation de la pâte crue, puis une recontamination/post-contamination a encore amplifié le nombre sur les beignets cuits.
\end{enumerate}

\textbf{4. Voies de contamination}
\begin{center}
\begin{tabularx}{\linewidth}{|l|X|X|}\hline
\rowcolor{popblueL}
\textbf{Via de contamination} & \textbf{Preuve (Doc 1)} & \textbf{Mode de vie / Mécanisme microbien} \\ \hline
\textbf{1. Contamination primaire (Pâte crue) : Portage sain mains/peau $\rightarrow$ Aliment} & « Je n'ai pas eu le temps de me laver les mains... je les ai juste essuyées sur mon pagne » après manipulation bois/balayage. & \textit{S. aureus} est un \textbf{commensal de la peau et des muqueuses nasales} (portage sain \textasciitilde 30-50\% population). Manipulation à mains sales = ensemencement direct. La pâte riche en protéines/glucides ($a_w$ élevé) est un \textbf{milieu de culture idéal} (nutrition hétérotrophe). \\ \hline
\textbf{2. Contamination secondaire / Post-cuisson (Beignets cuits) : Environnement / Mains $\rightarrow$ Aliment cuit} & « Les derniers beignets sont restés à l'air libre dans une corbeille en osier » + manipulation pour la vente sans lavage mains. & \textit{S. aureus} ne sporule pas (tué par cuisson), mais \textbf{recontamination} possible après cuisson par mains sales ou environnement (poussières, insectes). Production d'\textbf{entérotoxines thermostables} (résistent à \SI{100}{\celsius} 30 min) si croissance post-cuisson. \\ \hline
\end{tabularx}
\end{center}

\textbf{5. Mécanisme pathogène : Intoxication (et non infection invasive)}
\begin{itemize}
    \item \textbf{Arguments cliniques :} Incubation très courte (\SI{2}{\hour} à \SI{4}{\hour}), symptômes \textbf{émétiques majeurs} (vomissements), \textbf{absence de fièvre}, pas de sang dans les selles. Image typique d'une \textbf{intoxication alimentaire par entérotoxine préformée}.
    \item \textbf{Arguments microbiologiques :} \textit{S. aureus} produit des \cle{entérotoxines} (protéines) lors de sa croissance (phase exponentielle/stationnaire) dans l'aliment (pâte crue puis beignets). Ces toxines sont \textbf{thermostables} (ne sont pas détruites par la friture courte). L'ingestion de la toxine préformée déclenche le réflexe vomitif via le système nerveux central (récepteurs sérotoninergiques 5-HT3 du chimiorécepteur trigger zone) sans nécessiter de multiplication bactérienne dans l'intestin ni invasion tissulaire $\rightarrow$ pas de fièvre, pas de diarrhée inflammatoire (bien qu'une diarrhée osmotique puisse survenir secondairement).
\end{itemize}

\textbf{6. Conclusion et Prévention}
\textbf{Conclusion :} Cette toxi-infection alimentaire collective est due à l'ingestion d'\textbf{entérotoxines staphylococciques préformées} dans les beignets de haricots. La cause racine est une \textbf{hygiène défaillante} (mains non lavées, portage sain de \textit{S. aureus} par la vendeuse) couplée à une \textbf{mauvaise maîtrise des températures} : conservation de la pâte crue \SI{17}{\hour} à \SI{32}{\celsius} (zone de croissance optimale, Doc 3) permettant une multiplication massive (\textgreater \num{10^{11}} UFC/g) et synthèse de toxines, suivie d'une recontamination post-cuisson. La cuisson a tué les bactéries mais pas les toxines thermostables.

\textbf{Trois mesures d'hygiène concrètes et immédiates :}
\begin{enumerate}
    \item \textbf{Lavage obligatoire des mains à l'eau et au savon (ou solution hydroalcoolique)} \textbf{avant toute manipulation} de la pâte, après chaque interruption (balayage, bois, toilettes) et avant le service/vente. \textit{(Cible : Voie 1 - Portage sain mains $\rightarrow$ Aliment)}.
    \item \textbf{Respect de la chaîne du froid / Limitation du temps à température ambiante :} Préparer la pâte \textbf{le jour même} (pas la veille) ou la conserver au réfrigérateur (\textless \SI{4}{\celsius}) si préparation anticipée. Ne pas laisser la pâte crue \textgreater \SI{2}{\hour} à T° ambiante. \textit{(Cible : Multiplication exponentielle dans la pâte - Doc 3)}.
    \item \textbf{Protection des aliments cuits et service hygiénique :} Déposer les beignets cuits dans un \textbf{présentoir fermé (vitrine chauffée \textgreater \SI{63}{\celsius} ou réfrigérée \textless \SI{4}{\celsius})} et non dans une corbeille à l'air libre. Utiliser une \textbf{pince ou des gants jetables} pour le service (pas de contact mains nues). \textit{(Cible : Voie 2 - Recontamination post-cuisson + croissance résiduelle)}.
\end{enumerate}
\tcblower
\textbf{Commentaire pédagogique :} Cette activité mobilise la compétence « \emph{Mobiliser les connaissances sur les microorganismes pour expliquer une situation de la vie courante} ». Elle exige le passage du qualitatif (identification germe) au quantitatif (calcul croissance exponentielle), l'articulation entre cours (morphologie, reproduction, thermosensibilité) et documents (courbe, résultats culture, témoignage), et la rédaction d'une conclusion argumentée. L'erreur fréquente est de confondre infection et intoxication, ou d'oublier que la cuisson ne détruit pas les toxines déjà produites. Le calcul de $n$ (nombre de générations) doit utiliser le temps de génération \textbf{corrigé par la vitesse relative} lue sur la courbe, et non le temps de génération optimal à \SI{37}{\celsius} seul.
\end{exoresolu}

\courssub{Bilan de l'activité}
Cette activité d'intégration a permis de mettre en évidence, par une démarche d'investigation complète, que :
\begin{itemize}
    \item Les \cle{bactéries} (ici \textit{Staphylococcus aureus}, cocques Gram + en grappes) sont des microorganismes ubiquitaires, contaminants fréquents des aliments par \cle{portage sain} humain (peau, narines) ou environnement.
    \item Leur \cle{multiplication exponentielle} par division binaire est conditionnée par la température : la zone mésophile (\SI{10}{\celsius}-\SI{45}{\celsius}, optimum \SI{37}{\celsius}) rend la conservation à température ambiante tropicale (\SI{32}{\celsius}) \textbf{extrêmement favorable} à une pullulation rapide (calcul de $N_t = N_0 2^{t/g}$).
    \item La \cle{cuisson} (chaleur) détruit les formes végétatives (stérilisation) mais \textbf{ne détruit pas les entérotoxines thermostables} déjà sécrétées dans l'aliment pendant la phase de croissance.
    \item La \cle{contamination} peut être \textbf{primaire} (matière première/manipulation crue) ou \textbf{secondaire/post-cuisson} (réintroduction sur aliment cuit), d'où la nécessité de barrières successives : hygiène des mains, chaîne du froid, protection des aliments cuits, cuisson à cœur.
    \item La distinction clinique \textbf{infection (fièvre, incubation longue, germe vivant)} / \textbf{intoxication (pas fièvre, incubation courte, toxine préformée)} est essentielle pour l'orientation diagnostique et la prévention.
\end{itemize}
L'élève a ainsi consolidé les savoirs (groupes, nutrition, reproduction, contamination) et les savoir-faire (interprétation culture, exploitation courbe, calcul, rédaction argumentée) de la leçon dans un contexte réel de \cle{santé publique} au Cameroun.

\newpage

\coursec{Méthodes et exercices résolus}

\coursec{Microorganismes dans notre environnement}

\courssub{Découverte des microorganismes}

\begin{definition}[Microorganisme (Microbe)]
    Un \textbf{microorganisme} (ou \textbf{microbe}) est un organisme vivant de taille microscopique, généralement invisible à l'œil nu (taille < 0,1 mm), qui ne peut être observé qu'au microscope. Ils sont omniprésents dans la biosphère : air, eau, sol, peau, tube digestif, aliments.
\end{definition}

\begin{savaistu}[Histoire des sciences : La naissance de la microbiologie]
    \begin{itemize}
        \item \textbf{1665 :} Robert \textbf{Hooke} observe des cellules de liège et des moisissures au microscope (\emph{Micrographia}).
        \item \textbf{1676 :} Antonie van \textbf{Leeuwenhoek} (lentilles simples $\times 200-\times 300$) découvre les « animalcules » (bactéries, protistes, spermatozoïdes) dans l'eau de pluie, la salive, le tartre dentaire.
        \item \textbf{1857-1864 :} Louis \textbf{Pasteur} démontre le rôle des microbes dans la fermentation (levure) et la putréfaction, réfute la génération spontanée (expérience du col de cygne) et pose les bases de l'asepsie et de la vaccination.
        \item \textbf{1876 :} Robert \textbf{Koch} isole \emph{Bacillus anthracis} (charbon), formule les \textbf{postulats de Koch} (preuve de pathogénicité) et développe les cultures sur gélose.
    \end{itemize}
    Au Cameroun, l'Institut Pasteur de Yaoundé (créé 1959) et le Centre Pasteur du Cameroun (Garoua) sont des acteurs majeurs de la surveillance des épidémies (choléra, fièvre jaune, COVID-19, poliomyélite) et de la recherche sur les maladies tropicales.
\end{savaistu}

\begin{propriete}[Rôles écologiques et économiques majeurs]
    Les microorganismes ne sont pas seulement des agents pathogènes. Ils sont \textbf{indispensables à la vie sur Terre} :
    \begin{itemize}
        \item \textbf{Décomposeurs (Saprotrophes) :} Recyclage de la matière organique morte $\rightarrow$ minéralisation (cycle du C, N, P). Ex: Bactéries du sol, champignons (moisissures).
        \item \textbf{Producteurs primaires :} Photosynthèse (Cyanobactéries, Euglènes) $\rightarrow$ production d'$O_2$ et de matière organique (base des chaînes alimentaires aquatiques).
        \item \textbf{Symbiotes mutualistes :} Flore intestinale (vitamines K, B12, digestion cellulose, barrière pathogènes), Nodules de légumineuses (\emph{Rhizobium} : fixation azote atmosphérique $N_2 \rightarrow NH_3$), Lichens (champignon + algue).
        \item \textbf{Biotechnologies / Agro-alimentaire :} Pain, vin, bière, yaourt, fromage, choucroute, antibiotiques (Pénicilline de \emph{Penicillium}), enzymes (lessives), vaccins, bioéthanol, dépollution (bioremédiation).
        \item \textbf{Pathogènes :} Maladies infectieuses (humaines, animales, végétales). Représentent une minorité des espèces microbiennes.
    \end{itemize}
\end{propriete}

\courssub{Les différents groupes de microorganismes}

\begin{definition}[Classification en 4 groupes majeurs (Programme 3ème)]
    On distingue quatre grands groupes selon l'organisation cellulaire (procaryote/eucaryote/acellulaire), la taille, le mode de nutrition et la locomotion :
\end{definition}

\begin{center}
\begin{tabularx}{\linewidth}{|l|X|X|X|X|}\hline
\rowcolor{popblueL}
\textbf{Critère} & \textbf{Bactéries} & \textbf{Virus} & \textbf{Champignons microscopiques} & \textbf{Organismes unicellulaires (Protistes)} \\ \hline
\textbf{Statut cellulaire} & \textbf{Procaryote} (cellule complète) & \textbf{Acellulaire} (particule) & \textbf{Eucaryote} (cellule complète) & \textbf{Eucaryote} (cellule complète) \\ \hline
\textbf{Taille typique} & 0,5 -- 5 $\mu$m (visible MO $\times 400$) & 20 -- 300 nm (invisible MO, visible ME) & Levures : 5 -- 10 $\mu$m ; Moisissures : filaments (hyphes) & 10 -- 200 $\mu$m (visibles MO $\times 100-\times 400$) \\ \hline
\textbf{Structure clé} & Paroi \textbf{peptidoglycane}, pas de noyau (nucléoïde), pas d'organites membranaires, flagelles (certains) & \textbf{Capside} (protéines) + \textbf{Acide nucléique} (ADN \textbf{ou} ARN), pas de métabolisme, pas de ribosomes & Paroi \textbf{chitine} (levures) ou cellulose, \textbf{noyau} vrai, organites (mitochondries, vacuoles), pas de chloroplaste (hétérotrophes) & Noyau vrai, organites, \textbf{organites de locomotion} (cils, flagelles, pseudopodes), parfois chloroplastes \\ \hline
\textbf{Reproduction} & Division binaire (scissiparité) & Multiplication intracellulaire (cycle lytique/lysogénique) & Bourgeonnement (levures), Sporulation (moisissures), Sexuée (rare) & Division transversale (mitose), Conjugaison (Paramecium) \\ \hline
\textbf{Nutrition} & Hétérotrophe (absorption) ou Autotrophe (photosynthèse/chimiosynthèse) & \textbf{Aucune} (parasite intracellulaire strict) & Hétérotrophe \textbf{par absorption} (enzymes extracellulaires) & Hétérotrophe \textbf{par ingestion} (phagocytose) ou Autotrophe (photosynthèse) \\ \hline
\textbf{Exemples} & \emph{E. coli}, \emph{Staphylococcus}, \emph{Vibrio cholerae}, \emph{Lactobacillus}, Cyanobactéries & Grippe (Influenza), VIH, COVID-19 (SARS-CoV-2), Rougeole, Bactériophages & \emph{Saccharomyces} (levure), \emph{Penicillium}, \emph{Aspergillus}, \emph{Candida} & \emph{Paramecium}, \emph{Euglena}, \emph{Amoeba}, \emph{Plasmodium} (paludisme) \\ \hline
\end{tabularx}
\end{center}

\begin{attention}[Pièges de vocabulaire BEPC]
    \begin{itemize}
        \item \textbf{Virus $\neq$ Microbe cellulaire.} Ne jamais écrire « cellule virale », « noyau du virus », « respiration du virus ». C'est une \textbf{particule infectieuse}.
        \item \textbf{Bactéries $\neq$ Animaux/Plantes.} Ce sont des \textbf{procaryotes} (règne \emph{Bacteria} ou \emph{Monera} selon classification).
        \item \textbf{Levures = Champignons unicellulaires.} Ne pas les confondre avec les bactéries (elles ont un noyau !).
        \item \textbf{Protistes = Eucaryotes unicellulaires libres.} \emph{Paramecium} se déplace par cils, \emph{Euglena} par flagelle (et photosynthétise), \emph{Amoeba} par pseudopodes.
    \end{itemize}
\end{attention}

\courssub{Observation des microbes au microscope optique (TP N°5)}

\begin{methode}[Réaliser une observation au microscope optique (TP N°5)]
    \textbf{Objectif :} Observer des microorganismes (levures, bactéries du yaourt, eau de mare) et réaliser un dessin d'observation scientifique conforme aux attendus du BEPC.
    \begin{enumerate}
        \item \textbf{Préparer la lame :} Déposer une goutte de la préparation (ex: eau de mare, dilution de levure de boulanger dans l'eau, yaourt dilué dans l'eau physiologique) sur une lame porte-objet propre. Recouvrir délicatement d'une lamelle en l'inclinant à \textbf{45°} pour chasser les bulles d'air. Essuyer l'excédent avec du papier buvard.
        \item \textbf{Mettre au point (Règle d'or : toujours commencer au plus petit grossissement) :}
              \begin{enumerate}
                  \item Placer la lame sur le plateau, la fixer avec les pinces.
                  \item Sélectionner l'objectif de \textbf{plus faible grossissement} (ex: $\times 4$ ou $\times 10$).
                  \item \textbf{Regarder par le côté} (pas dans l'oculaire) et descendre le tube (ou monter le plateau) au maximum \textbf{sans toucher la lame}.
                  \item Regarder dans l'oculaire et \textbf{remonter le tube lentement} (mise au point grossière puis fine) jusqu'à obtenir une image nette.
              \end{enumerate}
        \item \textbf{Centrer et grossir :} Centrer la zone d'intérêt au centre du champ. Passer à l'objectif supérieur (ex: $\times 40$) en \textbf{vérifiant par le côté} qu'il ne touche pas la lamelle. Refaire la mise au point \textbf{uniquement avec la vis micrométrique (fine)}. Ne jamais toucher à la vis macrométrique aux forts grossissements (risque de casser la lamelle).
        \item \textbf{Observer et dessiner :} Au grossissement le plus pertinent (souvent $\times 400$ pour bactéries/levures ; $\times 100$ pour protistes), réaliser le dessin d'observation au \textbf{crayon à papier (HB ou 2H)} sur \textbf{papier blanc non quadrillé}.
        \item \textbf{Légender obligatoirement :}
              \begin{itemize}
                  \item Titre : Nom de l'échantillon + coloration si utilisée (ex: « Levure de boulanger, frais, $\times 400$ »).
                  \item \textbf{Grossissement total :} $G_t = G_{objectif} \times G_{oculaire}$ (ex: $40 \times 10 = \times 400$).
                  \item \textbf{Barre d'échelle} (représentant 10 $\mu$m ou 50 $\mu$m selon grossissement) \textbf{OU} mention « Taille réelle : X $\mu$m ».
                  \item Légendes des structures visibles (paroi, noyau, vacuole, bourgeon, cils, flagelle, chloroplaste...) par traits de rappel fins et fléchés.
              \end{itemize}
    \end{enumerate}
    \textbf{Consignes de sécurité et entretien :} Nettoyer les lentilles \textbf{uniquement} avec du papier à lentilles (jamais le doigt, ni papier ordinaire, ni tissu). Ne jamais forcer sur les vis. Jeter les lames et lamelles usagées dans le \textbf{bac à déchets biologiques tranchants} (contenant jaune à aiguilles/lames). Se laver les mains après manipulation.
\end{methode}

\begin{exemplebox}[Dessin d'observation scientifique : Levure de boulanger ($\times 400$)]
    \begin{popfigure}
        \begin{tikzpicture}[scale=1.2]
            % Cellule mère
            \draw[popdark, thick] (0,0) ellipse (1.2cm and 0.8cm);
            % Noyau
            \draw[popblue, thick, fill=popblue!20] (-0.3,0.1) circle (0.15cm);
            \node[popblue] at (-0.3,0.1) {N};
            % Vacuole
            \draw[popgreen, thick, fill=popgreen!10] (0.5,-0.2) ellipse (0.25cm and 0.15cm);
            \node[popgreen, font=\tiny] at (0.5,-0.2) {V};
            % Bourgeon
            \draw[popdark, thick, fill=popcream] (1.2,0.3) .. controls (1.6,0.5) and (1.6,-0.1) .. (1.2,-0.2) .. controls (1.0,-0.3) and (1.0,0.4) .. cycle;
            % Noyau dans bourgeon
            \draw[popblue, thick, fill=popblue!20] (1.4,0.1) circle (0.08cm);
            % Paroi
            \node[popdark, align=left] at (-2.5, 1.2) {\textbf{Légendes :}};
            \node[popdark, align=left] at (-2.5, 0.9) {1. Paroi cellulaire (chitine)};
            \node[popdark, align=left] at (-2.5, 0.6) {2. Noyau (N)};
            \node[popdark, align=left] at (-2.5, 0.3) {3. Vacuole (V)};
            \node[popdark, align=left] at (-2.5, 0.0) {4. Cytoplasme};
            \node[popdark, align=left] at (-2.5, -0.3) {5. Bourgeon (reproduction)}; 
            % Barre d'echelle
            \draw[popdark, very thick] (0.5, -1.2) -- (1.5, -1.2);
            \node[popdark, below] at (1.0, -1.2) {10 $\mu$m};
            % Grossissement
            \node[popdark, right] at (2.5, -1.2) {$G_t = \times 400$};
        \end{tikzpicture}
        \legende{Dessin d'observation au crayon à papier : cellule de levure (\emph{Saccharomyces cerevisiae}) en phase de bourgeonnement. Traits fins, continus, pas de hachures, pas de coloriage. Barre d'échelle et grossissement indiqués.}
    \end{popfigure}
\end{exemplebox}

\courssub{Nutrition et respiration des microbes}

\begin{definition}[Types de nutrition]
    \begin{itemize}
        \item \textbf{Autotrophe :} L'organisme utilise le \textbf{CO$_2$} comme source de carbone. Énergie fournie par la lumière (\textbf{photosynthèse} : cyanobactéries, euglènes) ou par l'oxydation de composés inorganiques (\textbf{chimiosynthèse} : bactéries nitrifiantes, bactéries sulfuriques).
        \item \textbf{Hétérotrophe :} L'organisme utilise de la \textbf{matière organique préformée} comme source de carbone et d'énergie.
              \begin{itemize}
                  \item \textbf{Par absorption (osmotrophie) :} Sécrétion d'enzymes extracellulaires $\rightarrow$ digestion externe $\rightarrow$ absorption des molécules simples. \textbf{Champignons (levures, moisissures), Bactéries décomposeuses.}
                  \item \textbf{Par ingestion (phagocytose) :} Englobement de particules solides dans une vacuole digestive. \textbf{Protistes (Paramecium, Amibe).}
              \end{itemize}
    \end{itemize}
\end{definition}

\begin{definition}[Types de respiration (Devenir du pyruvate)]
    La respiration permet de produire de l'ATP. Elle dépend de la présence d'oxygène ($O_2$) et de l'équipement enzymatique du microbe.
    \begin{itemize}
        \item \textbf{Respiration aérobie (stricte ou facultative en présence d'$O_2$) :} $O_2$ est l'accepteur final d'électrons. Dégradation complète du glucose $\rightarrow$ $CO_2 + H_2O$. Rendement élevé ($\sim 38$ ATP/glucose). \textbf{Ex:} Bactéries aérobies strictes (\emph{Bacillus subtilis}), Levures (si $O_2$), Protistes.
        \item \textbf{Fermentation (Anaérobiose) :} Pas d'$O_2$ (ou absence de chaîne respiratoire). Le pyruvate est réduit pour régénérer le $NAD^+$ (indispensable à la glycolyse). Rendement faible ($2$ ATP/glucose). Deux types majeurs :
              \begin{itemize}
                  \item \textbf{Fermentation alcoolique :} Pyruvate $\rightarrow$ Acétaldéhyde + $CO_2$ $\rightarrow$ \textbf{Éthanol}. \textbf{Levures (\emph{Saccharomyces})}. Bilan : $\ce{C6H12O6 -> 2 C2H5OH + 2 CO2 + 2 ATP}$.
                  \item \textbf{Fermentation lactique :} Pyruvate $\rightarrow$ \textbf{Acide lactique}. \textbf{Bactéries lactiques (\emph{Lactobacillus, Streptococcus})}. Bilan : $\ce{C6H12O6 -> 2 CH3-CHOH-COOH + 2 ATP}$. (Pas de $CO_2$ pour les homofermentaires).
              \end{itemize}
        \item \textbf{Respiration anaérobie (stricte) :} Utilisation d'un accepteur final \textbf{autre que $O_2$} (nitrates $NO_3^-$, sulfates $SO_4^{2-}$, $CO_2$). Propres à certaines bactéries (dénitrifiantes, sulfato-réductrices, méthanogènes). $O_2$ souvent toxique.
    \end{itemize}
\end{definition}

\begin{propriete}[Classification respiratoire simplifiée (Niveau 3ème)]
    \begin{center}
    \begin{tabular}{|l|l|l|}\hline
    \rowcolor{popblueL}
    \textbf{Groupe} & \textbf{Comportement vis-à-vis de l'$O_2$} & \textbf{Exemples} \\ \hline
    \textbf{Aérobies strictes} & $O_2$ \textbf{indispensable} (respiration only) & \emph{Bacillus subtilis}, \emph{Pseudomonas}, Mycobactéries \\ \hline
    \textbf{Anaérobies facultatives} & Respirent \textbf{si $O_2$}, fermentent \textbf{si pas $O_2$} & \emph{E. coli}, \emph{Salmonella}, \textbf{Levures} \\ \hline
    \textbf{Aérobies tolérantes (Anaérobies aérobies tolérantes)} & Ne respirent \textbf{jamais} (pas de cytochromes), $O_2$ \textbf{non toxique} (enzymes de protection), fermentation permanente & \textbf{Bactéries lactiques} (\emph{Lactobacillus, Streptococcus}) \\ \hline
    \textbf{Anaérobies strictes} & $O_2$ \textbf{toxique} (pas d'enzymes de protection), fermentation ou respiration anaérobie & \emph{Clostridium} (tétanos, botulisme), Méthanogènes \\ \hline
    \end{tabular}
    \end{center}
\end{propriete}

\begin{experience}[Mise en évidence des deux types de fermentation (TP type)]
    \textbf{Matériel :} 2 tubes à essai, glucose, levure de boulanger (\emph{Saccharomyces}), bactéries lactiques (\emph{Lactobacillus} du yaourt), huile de paraffine, eau de chaux, papier pH.
    \textbf{Protocole :}
    \begin{enumerate}
        \item Tube A : Glucose + Levure. Tube B : Glucose + Yaourt (source lactobacilles).
        \item Boucher les deux tubes avec un bouchon percé muni d'un tube à essai renversé rempli d'eau de chaux (piège à $CO_2$) OU simplement verser une couche d'\textbf{huile de paraffine} (2-3 cm) pour chasser l'air et créer l'anaérobiose.
        \item Incuber 24h à 30-35°C.
        \item Tester : Eau de chaux (trouble = $CO_2$), pH (acide = acide lactique), odeur (éthanol).
    \end{enumerate}
    \textbf{Résultats attendus :}
    \begin{center}
    \begin{tabular}{|c|c|c|c|}\hline
    \rowcolor{popblueL}
    \textbf{Tube} & \textbf{Microbe} & \textbf{Eau de chaux} & \textbf{pH / Odeur} \\ \hline
    A & Levure & \textbf{Trouble} ($CO_2$) & Neutre / \textbf{Odeur éthanol} (fermentation alcoolique) \\ \hline
    B & Lactobacille & \textbf{Limpide} (pas $CO_2$) & \textbf{Acide (pH $\sim$ 4)} / Odeur aigre (fermentation lactique) \\ \hline
    \end{tabular}
    \end{center}
    \textbf{Interprétation :} L'huile de paraffine assure l'anaérobiose stricte. La levure (anaérobie facultative) fermente $\rightarrow$ $CO_2$ + Éthanol. Le lactobacille (aérobie tolérant) fermente $\rightarrow$ Acide lactique (pas de $CO_2$ gazeux).
    \textbf{Applications :} Levure $\rightarrow$ Pain ($CO_2$ lève la pâte), Alcool. Lactobacilles $\rightarrow$ Yaourt, fromage, choucroute, saucisson (acidification = conservation + texture + goût).
\end{experience}

\courssub{Reproduction des microbes}

\begin{definition}[Reproduction asexuée (clonale) vs Sexuée]
    \begin{itemize}
        \item \textbf{Asexuée :} Un seul parent $\rightarrow$ descendants \textbf{génétiquement identiques} (clones). Mécanisme : Mitose (eucaryotes) ou Division binaire (procaryotes). Rapide, efficace en milieu favorable.
        \item \textbf{Sexuée :} Fusion de deux gamètes (haploïdes) $\rightarrow$ Zygote (diploïde) $\rightarrow$ Méiose $\rightarrow$ Descendants \textbf{génétiquement différents} (brassage). Survient souvent en conditions défavorables (manque nutriments, stress) pour former des formes de résistance (spores, zygotes).
    \end{itemize}
\end{definition}

\begin{propriete}[Modes de reproduction par groupe]
    \begin{itemize}
        \item \textbf{Bactéries (Procaryotes) :} \textbf{Division binaire (Scissiparité).}
              \begin{enumerate}
                  \item Réplication de l'ADN circulaire (chromosome unique) $\rightarrow$ 2 chromosomes attachés à la membrane.
                  \item Allongement de la cellule, ségrégation des chromosomes aux pôles.
                  \item Formation du \textbf{septum} (nouvelle paroi transversale) au centre.
                  \item Séparation en \textbf{2 cellules filles identiques (clones)}.
              \end{enumerate}
              Temps de génération : $\sim 20$ min (\emph{E. coli} optimal, 37°C, milieu riche). \textbf{Transferts horizontaux} (Conjugaison, Transformation, Transduction) = \textbf{brassage génétique}, \textbf{pas reproduction} (pas d'augmentation du nombre d'individus).
        \item \textbf{Levures (Champignons unicellulaires) :} \textbf{Bourgeonnement} (asexué dominant). Un bourgeon se forme, noyau se divise par mitose, un noyau migre dans le bourgeon $\rightarrow$ croissance $\rightarrow$ séparation (cicatrice de bourgeon visible). \textbf{Reproduction sexuée} possible : Fusion de 2 cellules haploïdes de types $MAT\alpha$ et $MATa$ $\rightarrow$ Zygote diploïde $\rightarrow$ Méiose $\rightarrow$ Asque à 4-8 ascospores haploïdes (résistance).
        \item \textbf{Moisissures (Champignons filamenteux) :} \textbf{Sporulation asexuée} : Formation de \textbf{conidies} (spores asexuées) au bout de \textbf{conidiophores} (pieds porteurs) $\rightarrow$ dispersion aérienne massive (clonage). \textbf{Reproduction sexuée} : Fusion de hyphes compatibles $\rightarrow$ formation de structures résistantes (thécas, zygospores) $\rightarrow$ méiose.
        \item \textbf{Protistes (ex: \emph{Paramecium}) :} \textbf{Division transversale} (asexuée, mitose, plan de division transverse). \textbf{Conjugaison} (phénomène sexué) : 2 paramecia s'unissent par la face ventrale (pont cytoplasmique) $\rightarrow$ méiose des micronoyaux $\rightarrow$ échange de noyaux haploïdes migrateurs $\rightarrow$ fécondation croisée $\rightarrow$ nouveau macronoyau. But : \textbf{rajeunissement} (macronoyau neuf) et \textbf{brassage génétique}.
        \item \textbf{Virus :} \textbf{Pas de reproduction autonome.} \textbf{Multiplication parasitaire} (cycle lytique ou lysogénique). Étapes : \textbf{Fixation} (récepteurs) $\rightarrow$ \textbf{Pénétration/Injection} (acide nucléique seul) $\rightarrow$ \textbf{Usine virale} (réplication génome + synthèse protéines capsides/enzymes sur ribosomes hôtes) $\rightarrow$ \textbf{Assemblage} (auto-assemblage) $\rightarrow$ \textbf{Lyse/Libération} (lyse cellulaire) ou \textbf{Budding} (bourgeonnement enveloppé).
    \end{itemize}
\end{propriete}

\begin{attention}[Distinction cruciale BEPC]
    \begin{itemize}
        \item \textbf{Division binaire (Bactéries)} $\neq$ \textbf{Mitose (Eucaryotes)}. Résultat identique (2 filles clones), mais mécanisme différent : pas de fuseau mitotique, pas de condensation chromosomique visible, pas de membrane nucléaire à dissoudre chez les bactéries.
        \item \textbf{Conjugaison bactérienne} $\neq$ \textbf{Reproduction}. C'est un transfert d'ADN (plasmide ou chromosome) d'un donneur ($F^+$ ou $Hfr$) vers un receveur ($F^-$). Le nombre de bactéries ne double pas pendant la conjugaison.
        \item \textbf{Virus} : On dit « \textbf{multiplication} » ou « \textbf{réplication} », jamais « reproduction » ni « division ».
    \end{itemize}
\end{attention}

\courssub{La contamination par les microorganismes}

\begin{definition}[Voies de pénétration (Portails d'entrée) et Vecteurs/Véhicules]
    Pour provoquer une infection, le microbe doit franchir les barrières de l'hôte. On distingue la \textbf{voie} (porte d'entrée anatomique), le \textbf{véhicule} (milieu transportant le microbe) et le \textbf{vecteur} (être vivant transportant activement le microbe).
    \begin{center}
    \begin{tabularx}{\linewidth}{|l|X|X|X|}\hline
    \rowcolor{popblueL}
    \textbf{Voie de contamination} & \textbf{Portail d'entrée} & \textbf{Véhicules / Vecteurs principaux} & \textbf{Exemples de maladies (Cameroun / Monde)} \\ \hline
    \textbf{Respiratoire} & Muqueuses nasales, trachée, bronches, alvéoles (inhalation gouttelettes \textgreater 5 $\mu$m, aérosols \textless 5 $\mu$m, poussières) & Air (gouttelettes Flügge), poussières contaminées & Grippe, COVID-19, Tuberculose (\emph{Mycobacterium tuberculosis}), Méningite à méningocoque, Pneumonie, Coqueluche \\ \hline
    \textbf{Digestive} & Bouche, œsophage, estomac, \textbf{intestin} (iléon, côlon) (ingestion) & \textbf{Eau} souillée (fèces), \textbf{Aliments} mal lavés/cuits, \textbf{Mains sales} (véhicules) & Choléra (\emph{Vibrio cholerae}), Fièvre typhoïde (\emph{Salmonella Typhi}), Hépatite A, Amibiase (\emph{Entamoeba histolytica}), Toxi-infections (\emph{Salmonella}, \emph{Staph. aureus}) \\ \hline
    \textbf{Cutanée-Muqueuse (Percée)} & Peau lésée (plaie, brûlure, égratignure), muqueuses génitales, oculaires, voie sanguine directe & \textbf{Vecteurs} : Moustiques (\emph{Anopheles} $\rightarrow$ Paludisme), Tiques ; \textbf{Objets} contaminés (aiguilles, rasoirs) ; \textbf{Rapports sexuels} (MST) ; Transfusion ; Accouchement & \textbf{Paludisme} (\emph{Plasmodium}), \textbf{VIH/Sida}, Hépatites B/C, Tétanos (\emph{Clostridium tetani}), Syphilis, Fièvre jaune \\ \hline
    \textbf{Verticale (Mère $\rightarrow$ Enfant)} & Barrière placentaire, voies génitales (accouchement), lait maternel & Sang maternel, sécrétions vaginales, lait & VIH, Syphilis (\emph{Treponema pallidum}), Rubéole, CMV, Toxoplasmose, Hépatite B \\ \hline
    \end{tabularx}
    \end{center}
\end{definition}

\begin{propriete}[Barrières naturelles non spécifiques (Innées) franchies par le pathogène]
    L'organisme possède des défenses de \textbf{première ligne} (barrières) et de \textbf{deuxième ligne} (inflammation, phagocytose, fièvre) qui agissent \textbf{sans spécificité} antigénique.
    \begin{itemize}
        \item \textbf{Barrières physiques et chimiques (1ère ligne) :}
              \begin{itemize}
                  \item \textbf{Peau :} Kératine (imperméable), pH acide ($\sim 5,5$), flore commensale (concurrence, bactériocines), sébum (antiseptique). \textit{Franchie par : plaie, piqûre, brûlure.}
                  \item \textbf{Muqueuses respiratoires :} Cils vibratiles + mucus (piégeage $\rightarrow$ expectoration/mucociliary clearance), éternuement, toux. \textit{Franchie par : inhalation massive, virus ciliotoxiques, tabac (paralysie cils).}
                  \item \textbf{Estomac :} \textbf{Acidité chlorhydrique (pH 1-2)} : tue la majorité des bactéries ingérées. \textit{Franchie par : dose infectante élevée, protection par aliments (tampon), résistance acide (V. cholerae, H. pylori).}
                  \item \textbf{Autres :} Lysozyme (larmes, salive, mucus - hydrolyse paroi bactérienne), Flore intestinale/vaginale (concurrence, pH acide), Urine (flux mécanique), Spermine (sperme).
              \end{itemize}
        \item \textbf{Réaction inflammatoire (2ème ligne) :} Si barrières franchies $\rightarrow$ Vasodilatation, perméabilité vasculaire $\rightarrow$ arrivée \textbf{phagocytes} (Polynucléaires neutrophiles, Macrophages) $\rightarrow$ \textbf{Phagocytose} + enzymes lysosomales. \textbf{Fièvre} (ralentit multiplication bactérienne, stimule immunité).
    \end{itemize}
\end{propriete}

\begin{exemplebox}[Cas d'école : Choléra vs Tétanos au Cameroun]
    \begin{itemize}
        \item \textbf{Choléra (\emph{Vibrio cholerae}) :} Voie \textbf{Digestive}. Portail : \textbf{Intestin grêle}. Barrière critique : \textbf{Acidité gastrique}. Stratégie du microbe : Dose infectante élevée ($10^6-10^8$), protection par aliments (bouillie de maïs/riz tiède tamponne pH), réponse acide (systèmes tampons). Toxine cholérique $\rightarrow$ sécrétion massive $Cl^-/H_2O$ $\rightarrow$ diarrhée « eau de riz » $\rightarrow$ déshydratation mortelle. Prévention : \textbf{Hygiène eau/aliments/mains} + Vaccin oral (bactéries inactivées + sous-unité B toxine).
        \item \textbf{Tétanos (\emph{Clostridium tetani}) :} Voie \textbf{Cutanée (percée)}. Portail : \textbf{Plaie souillée (terre/poussière)}. Barrière critique : \textbf{Peau intacte}. Stratégie du microbe : \textbf{Spores} ultra-résistantes (chaleur, désinfectants, dessiccation). Germent \textbf{uniquement en anaérobiose stricte} (plaie profonde, nécrosée, mal vascularisée, corps étranger). Production \textbf{tétanospasmine} (neurotoxine bloquant inhibition glycine/GABA $\rightarrow$ contractures tétaniques, trismus, opisthotonos). Prévention : \textbf{Nettoyage/désinfection plaies} + \textbf{Vaccination Anatoxine Tétanique} (Toxoïde = toxine détotoxiquée par formol). Inclus dans \textbf{Pentavalent (DTC-HepB-Hib)} au PEV Cameroun. Rappels \textbf{tous les 10-20 ans}.
    \end{itemize}
\end{exemplebox}

\courssub{Méthodes et exercices résolus}

\begin{methode}[Réaliser une observation au microscope optique (TP N°5)]
    \textbf{Objectif :} Observer des microorganismes (levures, bactéries du yaourt, eau de mare) et réaliser un dessin d'observation scientifique.
    \begin{enumerate}
        \item \textbf{Préparer la lame :} Déposer une goutte de la préparation (ex: eau de mare, dilution de levure, yaourt dilué) sur une lame porte-objet. Recouvrir délicatement d'une lamelle (angle 45°) pour chasser les bulles d'air.
        \item \textbf{Mettre au point :} Placer la lame sur le plateau. Commencer \textbf{toujours} par l'objectif de plus faible grossissement (ex: $\times 4$ ou $\times 10$). Regarder \textbf{par le côté} et descendre le tube (ou monter le plateau) au maximum sans toucher la lame. Regarder dans l'oculaire et remonter le tube (mise au point grossière puis fine) jusqu'à obtenir une image nette.
        \item \textbf{Centrer et grossir :} Centrer la zone d'intérêt. Passer à l'objectif supérieur (ex: $\times 40$) en vérifiant par le côté qu'il ne touche pas la lamelle. Refaire la mise au point \textbf{uniquement avec la vis micrométrique (fine)}.
        \item \textbf{Observer et dessiner :} Au grossissement le plus pertinent (souvent $\times 400$ pour les bactéries/levures), réaliser le dessin d'observation au crayon à papier sur papier blanc non quadrillé.
        \item \textbf{Légender :} Indiquer le grossissement total ($G_t = G_{objectif} \times G_{oculaire}$), le nom de l'échantillon, la coloration si utilisée, et légender les structures visibles (paroi, noyau, vacuole, bourgeon...).
    \end{enumerate}
    \textbf{Reflexes :} Nettoyer les lentilles avec du papier à lentilles (jamais le doigt, ni papier ordinaire). Ne jamais forcer sur les vis. Jeter les lames usagées dans le bac à déchets biologiques (contenant à aiguilles/lames).
\end{methode}

\begin{exoresolu}[Observation de levures de boulanger au microscope]
    \textbf{Énoncé :} Un élève observe une goutte de levure de boulanger diluée dans l'eau, au microscope optique (oculaire $\times 10$, objectifs $\times 4$, $\times 10$, $\times 40$). Il utilise le grossissement $\times 400$ pour son dessin final.
    \begin{enumerate}
        \item Calculer le grossissement total utilisé pour le dessin.
        \item L'élève dessine une cellule de levure de 5 cm de long sur sa feuille. Sachant qu'une levure mesure environ 8 $\mu$m en réalité, vérifier si le dessin est à l'échelle du grossissement utilisé.
        \item Citer trois règles à respecter pour réaliser un dessin d'observation scientifique conforme.
        \item L'élève observe des petites excroissances sur certaines cellules. Nommer ce mode de reproduction et l'expliquer brièvement.
    \end{enumerate}
    \tcblower
    \textbf{Solution détaillée :}
    \begin{enumerate}
        \item \textbf{Calcul du grossissement total :} $G_t = G_{objectif} \times G_{oculaire} = 40 \times 10 = \mathbf{400}$. Le grossissement total est $\times 400$.
        \item \textbf{Vérification de l'échelle :}
            \begin{itemize}
                \item Taille réelle de la levure : $8 \,\mu\text{m} = 8 \times 10^{-3} \,\text{mm}$.
                \item Taille théorique du dessin au grossissement 400 : $8 \times 10^{-3} \,\text{mm} \times 400 = 3,2 \,\text{mm} = \mathbf{0,32 \,\text{cm}}$.
                \item Taille du dessin de l'élève : $\mathbf{5 \,\text{cm}}$.
            \end{itemize}
            Le dessin de l'élève (5 cm) est \textbf{beaucoup plus grand} que la taille théorique (0,32 cm). Il n'est \textbf{pas à l'échelle} du grossissement $\times 400$. C'est une erreur fréquente : on dessine « grand » pour voir les détails, mais il faut indiquer l'échelle ou la barre d'échelle.
        \item \textbf{Trois règles du dessin scientifique :}
            \begin{enumerate}
                \item Dessiner au \textbf{crayon à papier} (mine HB ou 2H) sur \textbf{papier blanc uni} (pas de quadrillage).
                \item Faire un \textbf{trait continu, fin et net} (pas de hachures, pas de traits de construction visibles, pas de coloriage).
                \item Respecter les \textbf{proportions} des structures observées (forme, taille relative) et ajouter une \textbf{barre d'échelle} ou la mention du grossissement.
            \end{enumerate}
        \item \textbf{Reproduction :} Il s'agit de la \textbf{reproduction asexuée par bourgeonnement}. Un bourgeon (petite excroissance) se forme sur la cellule mère, le noyau se divise par mitose, un noyau migre dans le bourgeon qui grossit, puis se sépare pour devenir une nouvelle cellule fille indépendante.
    \end{enumerate}
    \textbf{Conclusion :} L'élève a correctement calculé le grossissement ($\times 400$) mais son dessin n'est pas à l'échelle. Il a observé le bourgeonnement, mode de reproduction asexuée caractéristique des levures (champignons microscopiques).
\end{exoresolu}

\begin{methode}[Identifier et classer un microorganisme dans son groupe majeur]
    \textbf{Objectif :} À partir de critères observables (microscope, culture, documents), ranger un microbe dans l'un des 4 groupes : Bactéries, Virus, Champignons microscopiques (Levures/Moisissures), Protistes (Organismes unicellulaires libres).
    \begin{enumerate}
        \item \textbf{Observer la taille :} Bactéries et levures : quelques $\mu$m (visibles au MO $\times 400$). Virus : nanomètres (invisibles au MO, nécessitent ME). Protistes : souvent plus grands (10 à 200 $\mu$m).
        \item \textbf{Observer la structure cellulaire (au MO ou sur document) :}
            \begin{itemize}
                \item \textbf{Bactérie :} Cellule \textbf{procaryote} : pas de noyau individualisé (nucléoïde), pas d'organites membranaires (pas de mitochondries, pas de chloroplastes). Paroi en peptidoglycane. Formes : bacille, coccus, spirille.
                \item \textbf{Champignon (Levure/Moisissure) :} Cellule \textbf{eucaryote} : \textbf{noyau} visible, organites (vacuole, mitochondries). Paroi en \textbf{chitine} (levure) ou cellulose. Levure = unicellulaire (ovale, bourgeonnement). Moisissure = filaments (hyphes) formant un mycélium.
                \item \textbf{Protiste (ex: Paramecium, Euglène, Amibe) :} Cellule \textbf{eucaryote} complexe : noyau, vacuoles (contractile, digestive), cils/flagelles/pseudopodes, parfois chloroplastes (Euglène).
                \item \textbf{Virus :} \textbf{Non cellulaire} (particule). Capside protéique + acide nucléique (ADN ou ARN). Pas de métabolisme propre.
            \end{itemize}
        \item \textbf{Analyser le mode de nutrition :} Autotrophe (photosynthèse : certaines bactéries, Euglène) / Hétérotrophe (absorption : champignons, bactéries décomposeuses ; ingestion : protistes phagocytoses).
        \item \textbf{Conclure :} Remplir un tableau de détermination ou écrire une phrase justifiée : « Cet organisme est une bactérie car c'est une cellule procaryote sans noyau... »
    \end{enumerate}
\end{methode}

\begin{exoresolu}[Identification de microorganismes à partir de documents]
    \textbf{Énoncé :} On observe trois documents microscopiques (grossissement $\times 400$).
    \begin{center}
        \begin{tabular}{|c|p{5cm}|p{5cm}|}\hline
        \textbf{Doc} & \textbf{Description} & \textbf{Indice biochimique} \\ \hline
        1 & Cellules sphériques (cocques) groupées en grappes irrégulières. Pas de noyau visible. & Catalase + (bulles avec $H_2O_2$) \\ \hline
        2 & Cellules ovales, isolées ou en chaînes courtes. Noyau visible. Une cellule porte une excroissance. & Fermentation alcoolique ($CO_2$ dégagé) \\ \hline
        3 & Organisme unique, forme de savate, recouvert de cils. Deux noyaux (macro et micro). Vacuole contractile visible. & Hétérotrophe phagocyte (ingestion bactéries) \\ \hline
        \end{tabular}
    \end{center}
    \begin{enumerate}
        \item Pour chaque document, identifier le groupe d'appartenance (Bactérie, Levure, Protiste, Virus) et l'espèce probable si possible.
        \item Justifier chaque identification en citant \textbf{deux arguments morphologiques} et \textbf{un argument physiologique}.
        \item Préciser le mode de reproduction observé sur le document 2.
    \end{enumerate}
    \tcblower
    \textbf{Solution détaillée :}
    \begin{enumerate}
        \item \textbf{Document 1 :} \textbf{Bactérie} (probablement \emph{Staphylococcus} sp. / \emph{Staphylococcus aureus}).
              \textbf{Document 2 :} \textbf{Champignon microscopique : Levure} (probablement \emph{Saccharomyces cerevisiae}).
              \textbf{Document 3 :} \textbf{Protiste (Cilié)} : \emph{Paramecium} sp.
        \item \textbf{Justifications :}
              \begin{itemize}
                  \item \textbf{Doc 1 (Bactérie) :}
                        \begin{itemize}
                            \item \textit{Morpho 1 :} Cellules procaryotes \textbf{sans noyau individualisé} (matériel génétique diffus).
                            \item \textit{Morpho 2 :} Taille microscopique ($\sim 1 \mu$m), forme sphérique (cocques) en \textbf{grappes}.
                            \item \textit{Physio :} Test \textbf{catalase positif} (décomposition $H_2O_2 \rightarrow H_2O + O_2$), typique des staphylocoques (aérobies/anaérobies facultatifs).
                        \end{itemize}
                  \item \textbf{Doc 2 (Levure) :}
                        \begin{itemize}
                            \item \textit{Morpho 1 :} Cellules \textbf{eucaryotes} avec \textbf{noyau visible}.
                            \item \textit{Morpho 2 :} Forme ovale, présence d'un \textbf{bourgeon} (reproduction asexuée).
                            \item \textit{Physio :} \textbf{Fermentation alcoolique} (production de $CO_2$ en anaérobiose), métabolisme caractéristique des levures.
                        \end{itemize}
                  \item \textbf{Doc 3 (Paramecium) :}
                        \begin{itemize}
                            \item \textit{Morpho 1 :} Cellule \textbf{eucaryote} complexe, forme de savate, \textbf{ciliaires} (cils vibratiles).
                            \item \textit{Morpho 2 :} Deux types de noyaux : \textbf{macronoyau} (végétatif) et \textbf{micronoyau} (reproduction), \textbf{vacuole contractile} (osmorégulation).
                            \item \textit{Physio :} Nutrition \textbf{hétérotrophe par phagocytose} (ingestion de bactéries dans vacuole digestive).
                        \end{itemize}
              \end{itemize}
        \item \textbf{Reproduction Doc 2 :} \textbf{Bourgeonnement} (reproduction asexuée). La cellule mère forme une excroissance (bourgeon) qui reçoit une copie du noyau (mitose) et se détache.
    \end{enumerate}
    \textbf{Conclusion :} L'analyse conjointe de la morphologie cellulaire (procaryote/eucaryote, organites, structures de locomotion) et du métabolisme (respiration, fermentation, nutrition) permet une classification rigoureuse dans les 3 grands groupes cellulaires (Bactéries, Eucaryotes unicellulaires, Virus exclus car non cellulaires).
\end{exoresolu}

\begin{methode}[Expliquer la nutrition et la respiration des microorganismes]
    \textbf{Objectif :} Relier le type de nutrition (source de carbone/énergie) et le type de respiration (acceptateur d'électrons final) au métabolisme du microbe et à son rôle écologique.
    \begin{enumerate}
        \item \textbf{Identifier la source de carbone :}
            \begin{itemize}
                \item \textbf{Autotrophe :} $CO_2$ (photosynthèse : cyanobactéries, euglènes ; chimiosynthèse : bactéries nitrifiantes).
                \item \textbf{Hétérotrophe :} Matière organique préformée.
            \end{itemize}
        \item \textbf{Identifier le mode respiratoire (devenir du pyruvate en présence/absence d'$O_2$) :}
            \begin{itemize}
                \item \textbf{Aérobie (stricte ou facultative) :} $O_2$ nécessaire pour respiration complète (Glycolyse + Krebs + Chaîne respiratoire). Rendement élevé ($\sim 38$ ATP).
                \item \textbf{Anaérobie stricte :} $O_2$ toxique. Fermentation seulement (Glycolyse + régénération $NAD^+$). Rendement $2$ ATP. Ex: \emph{Clostridium}, méthanogènes.
                \item \textbf{Anaérobie facultative :} Respire si $O_2$, fermente sinon. Ex: \emph{E. coli}, \emph{Saccharomyces} (levure).
                \item \textbf{Aérobie tolérant (Anaérobie aérobie tolérant) :} Ne respire pas (pas de chaîne respiratoire), mais $O_2$ non toxique (fermentation permanente). Ex: \emph{Lactobacillus}, \emph{Streptococcus}.
            \end{itemize}
        \item \textbf{Écrire le bilan global des fermentations :}
            \begin{itemize}
                \item Fermentation lactique : $\ce{C6H12O6 -> 2 CH3-CHOH-COOH + 2 ATP}$ (Acide lactique).
                \item Fermentation alcoolique : $\ce{C6H12O6 -> 2 C2H5OH + 2 CO2 + 2 ATP}$ (Éthanol + $CO_2$).
            \end{itemize}
        \item \textbf{Conclure sur le rôle :} Décomposeurs (recyclage $C, N$), Producteurs ($O_2$, Matière Organique), Pathogènes, Symbiotes (flore intestinale, nodules soja), Biotechnologie (pain, vin, fromage, antibiotiques, yaourt).
    \end{enumerate}
\end{methode}

\begin{exoresolu}[Métabolisme de la levure et du lactobacille]
    \textbf{Énoncé :} On réalise deux montages anaérobiose (bouchon à huile de paraffine) avec du glucose :
    \begin{itemize}
        \item Tube A : Levure de boulanger (\emph{Saccharomyces cerevisiae}).
        \item Tube B : Bactérie lactique (\emph{Lactobacillus bulgaricus}).
    \end{itemize}
    Après 24h à 30°C, on teste le gaz produit (eau de chaux) et le pH (papier pH).
    \begin{enumerate}
        \item Prédire les résultats des tests (eau de chaux, pH) pour chaque tube. Justifier par les équations-bilans.
        \item Classer ces deux microbes selon leur type de respiration (aérobie/anaérobie stricte/facultative/tolérante).
        \item Expliquer pourquoi l'huile de paraffine est indispensable.
        \item Citer une application industrielle pour chaque fermentation.
    \end{enumerate}
    \tcblower
    \textbf{Solution détaillée :}
    \begin{enumerate}
        \item \textbf{Tube A (Levure - Fermentation alcoolique) :}
              \begin{itemize}
                  \item \textbf{Eau de chaux :} \textbf{Trouble} (précipité de $CaCO_3$). Justification : $\ce{C6H12O6 -> 2 C2H5OH + 2 CO2+ 2 ATP}$. Le $CO_2$ gazeux traverse l'huile et trouble l'eau de chaux : $\ce{CO2 + Ca(OH)2 -> CaCO3+ H2O}$.
                  \item \textbf{pH :} \textbf{Légèrement acide ou neutre} (pH $\sim 4-5$ à cause de l'éthanol/acides mineurs, mais pas de forte acidification). Pas de production d'acide fort.
              \end{itemize}
              \textbf{Tube B (Lactobacille - Fermentation lactique) :}
              \begin{itemize}
                  \item \textbf{Eau de chaux :} \textbf{Reste limpide}. Justification : $\ce{C6H12O6 -> 2 CH3-CHOH-COOH + 2 ATP}$ (Acide lactique). \textbf{Pas de $CO_2$ gazeux produit} (fermentation homofermentaire).
                  \item \textbf{pH :} \textbf{Acide} (pH $\sim 4$ ou moins). L'accumulation d'acide lactique abaisse le pH.
              \end{itemize}
        \item \textbf{Classification respiratoire :}
              \begin{itemize}
                  \item \emph{Saccharomyces cerevisiae} : \textbf{Anaérobie facultative} (respire en présence d'$O_2$, fermente en son absence).
                  \item \emph{Lactobacillus bulgaricus} : \textbf{Aérobie tolérant} (ne possède pas de chaîne respiratoire, ne peut pas respirer, mais supporte l'$O_2$ ; il fermente en permanence).
              \end{itemize}
        \item \textbf{Rôle de l'huile de paraffine :} Créer des conditions \textbf{anaérobies strictes} (absence d'$O_2$ dissous dans le milieu) pour \textbf{forcer la fermentation} et empêcher la respiration aérobie (qui ne produirait pas d'éthanol ni d'acide lactique en grande quantité, mais du $CO_2$ et de l'eau). Elle isole le milieu de l'air atmosphérique.
        \item \textbf{Applications :}
              \begin{itemize}
                  \item Levure : \textbf{Boulangerie} (levée de la pâte par $CO_2$), \textbf{Brasserie/Œnologie} (production éthanol), \textbf{Bioéthanol} (carburant).
                  \item Lactobacille : \textbf{Yaourt / Fromage / Choucroute / Saucisson} (acidification $\rightarrow$ coagulation protéines, conservation, saveur), \textbf{Probiotiques}.
              \end{itemize}
    \end{enumerate}
    \textbf{Conclusion :} Bien que toutes deux soient des fermentations (rendement 2 ATP/glucose), la nature du produit final ($CO_2$ + éthanol vs Acide lactique) et la tolérance à l'oxygène distinguent fondamentalement ces deux groupes d'intérêt agro-alimentaire majeur.
\end{exoresolu}

\begin{methode}[Décrire les modes de reproduction des microorganismes]
    \textbf{Objectif :} Différencier reproduction asexuée (clonale) et sexuée, et décrire les cycles spécifiques à chaque groupe.
    \begin{enumerate}
        \item \textbf{Bactéries (Procaryotes) :} \textbf{Division binaire (scissiparité)}. Réplication de l'ADN circulaire $\rightarrow$ ségrégation aux pôles $\rightarrow$ formation du septum (paroi) $\rightarrow$ 2 cellules filles \textbf{identiques} (clones). Rapidité : temps de génération $\sim 20$ min (\emph{E. coli} optimal). Pas de méiose, pas de fécondation. \textit{Transferts horizontaux} (conjugaison, transformation, transduction) = brassage génétique $\neq$ reproduction.
        \item \textbf{Levures (Champignons unicellulaires) :} \textbf{Bourgeonnement} (asexué, mitose) dominant. \textbf{Reproduction sexuée} possible en conditions défavorables : deux cellules haploïdes de types sexuels différents ($MAT\alpha$ et $MATa$) fusionnent (plasmogamie + caryogamie) $\rightarrow$ zygote diploïde $\rightarrow$ méiose $\rightarrow$ asque à 4-8 spores haploïdes (ascospores).
        \item \textbf{Moisissures (Champignons filamenteux) :} \textbf{Sporulation asexuée} (conidies sur conidiophores) = dispersion clonale. \textbf{Reproduction sexuée} : formation de structures complexes (thécas, zygospores) après fusion de souches compatibles $\rightarrow$ méiose.
        \item \textbf{Protistes (ex: Paramecium) :} \textbf{Division transversale} (asexuée, mitose). \textbf{Conjugaison} (phénomène sexué) : 2 paramecia s'unissent par la face ventrale $\rightarrow$ échange de micronoyaux haploïdes (issus de méiose) $\rightarrow$ fécondation croisée $\rightarrow$ nouveau macronoyau. \textbf{But :} Rajeunissement et brassage génétique.
        \item \textbf{Virus :} \textbf{Pas de reproduction autonome}. \textbf{Cycle lytique} (multiplication + lyse) ou \textbf{lysogénique} (intégration génome viral = prophage/provirus). Étapes : Fixation $\rightarrow$ Pénétration/Injection $\rightarrow$ Réplication (usine virale) $\rightarrow$ Assemblage $\rightarrow$ Libération.
    \end{enumerate}
    \textbf{Point clé BEPC :} Ne pas confondre \textbf{division binaire} (bactéries, procaryotes, mitose simplifiée) et \textbf{mitose} (eucaryotes : levures, protistes, cellules somatiques). Le résultat est similaire (2 filles identiques), le mécanisme chromosomique diffère (pas de fuseau mitotique, pas de condensation chromosomique visible chez les bactéries).
\end{methode}

\begin{exoresolu}[Cycle de multiplication du bactériophage (Virus à ADN)]
    \textbf{Énoncé :} On étudie le cycle lytique d'un bactériophage T4 infectant \emph{Escherichia coli}.
    \begin{enumerate}
        \item Ordonner les étapes suivantes : Assemblage, Fixation, Pénétration/Injection, Réplication/Transcription/Traduction (Usine virale), Lyse/Libération.
        \item Sur un schéma de cycle (fourni mentalement), indiquer le moment où le génome viral est injecté et où les capsides sont synthétisées.
        \item Expliquer pourquoi les antibiotiques (ex: pénicilline) sont inefficaces sur ce virus.
        \item Calculer le nombre de virions produits après 3 cycles de lyse complets si une cellule libère 200 virions et que tous infectent une nouvelle cellule (conditions idéales).
    \end{enumerate}
    \tcblower
    \textbf{Solution détaillée :}
    \begin{enumerate}
        \item \textbf{Ordre chronologique du cycle lytique :}
              1. \textbf{Fixation} (reconnaissance récepteurs par fibres de queue).
              2. \textbf{Pénétration / Injection} (queue se contracte, ADN injecté dans le cytoplasme, capside vide reste dehors = « coque »).
              3. \textbf{Réplication / Transcription / Traduction (Usine virale)} : L'ADN viral détourne la machinerie de la bactérie (ARN polymérase, ribosomes, nucléotides, AA, ATP) pour synthétiser : \textbf{enzymes virales} (ex: lysozyme), \textbf{protéines structurales} (capside, queue), \textbf{répliquer l'ADN viral}.
              4. \textbf{Assemblage (Maturation)} : Auto-assemblage spontané des capsides autour de l'ADN viral $\rightarrow$ virions complets intracellulaires.
              5. \textbf{Lyse / Libération} : La lysozyme virale détruit la paroi bactérienne (peptidoglycane) $\rightarrow$ éclatement (lyse) $\rightarrow$ libération de $\sim 100-300$ virions.
        \item \textbf{Localisation :}
              \begin{itemize}
                  \item Injection ADN : Étape 2 (Pénétration). L'ADN entre seul dans le cytoplasme.
                  \item Synthèse capsides : Étape 3 (Usine virale), sur les ribosomes de l'hôte (protéines structurales « tardives »).
              \end{itemize}
        \item \textbf{Inefficacité des antibiotiques :} Les antibiotiques (comme la pénicilline) ciblent des structures ou métabolismes \textbf{propres aux bactéries} (synthèse paroi peptidoglycane, ribosomes 70S, ADN gyrase, métabolisme folates). Le virus \textbf{n'a pas de métabolisme propre}, \textbf{pas de paroi}, \textbf{pas de ribosomes}. Il utilise \textbf{exclusivement} la machinerie de la cellule hôte. Tuer le virus sans tuer la cellule hôte est impossible avec ces cibles.
        \item \textbf{Calcul de la multiplication :}
              \begin{itemize}
                  \item Cycle 1 : 1 cellule $\rightarrow$ 200 virions.
                  \item Cycle 2 : 200 cellules infectées $\rightarrow$ $200 \times 200 = 40\,000$ virions.
                  \item Cycle 3 : $40\,000$ cellules infectées $\rightarrow$ $40\,000 \times 200 = \mathbf{8\,000\,000}$ virions ($8 \times 10^6$).
              \end{itemize}
              Formule générale : $N = N_0 \times b^n$ ($N_0=1$, $b=200$, $n=3$).
    \end{enumerate}
    \textbf{Conclusion :} Le cycle lytique est une exploitation parasitaire totale de la machinerie cellulaire bactérienne aboutissant à sa destruction. L'absence de cible métabolique propre explique l'absence d'antibiotiques antiviraux classiques (on utilise des antiviraux ciblant des enzymes virales spécifiques : transcriptase inverse, protéase, neuraminidase).
\end{exoresolu}

\begin{methode}[Analyser les voies de contamination et les barrières de l'organisme]
    \textbf{Objectif :} Relier le portail d'entrée du microbe (voie de contamination) aux barrières naturelles (non spécifiques) franchies et aux maladies associées.
    \begin{enumerate}
        \item \textbf{Lister les 4 grandes voies de pénétration :}
            \begin{enumerate}
                \item \textbf{Vie respiratoire :} Inhalation (gouttelettes, aérosols, poussières). Maladies : Grippe, Covid-19, Tuberculose, Méningite, Pneumonie.
                \item \textbf{Vie digestive :} Ingestion (eau, aliments souillés, mains sales). Maladies : Choléra, Fièvre typhoïde, Hépatite A, Amibiase, Toxi-infections (Salmonelle).
                \item \textbf{Vie cutanée-muqueuse (percée) :} Plaies, piqûres (moustiques, tiques), rapports sexuels (MST), transfusion, accouchement. Maladies : Paludisme, VIH/Sida, Tétanos, Hépatite B/C, Syphilis.
                \item \textbf{Vie verticale (mère-enfant) :} Transplacentaire, accouchement, allaitement. Maladies : VIH, Syphilis, Rubéole, CMV, Toxoplasmose.
            \end{enumerate}
        \item \textbf{Identifier les barrières naturelles franchies (Non-spécifiques / Innées) :}
            \begin{itemize}
                \item \textbf{Barrières physiques/chimiques :} Peau (kératine, pH acide, flore commensale), Cils vibratiles + mucus (piégeage/expectoration), Acidité gastrique (pH 1-2), Lysozyme (larmes, salive), Flore intestinale/vaginale (concurrence, pH).
                \item \textbf{Réaction inflammatoire :} Vasodilatation, phagocytes (polynucléaires, macrophages), fièvre.
            \end{itemize}
        \item \textbf{Structurer la réponse :} « Le microbe X pénètre par la voie [respiratoire/digestive/cutanée]. Il franchit la barrière [muqueuse cillée / acidité gastrique / peau intacte] grâce à [facteur de virulence : capsule, toxine, mobilité, enzyme]. Il atteint le sang (septicémie) ou un organe cible. »
    \end{enumerate}
\end{methode}

\begin{exoresolu}[Contamination et barrières : Cas du Choléra et du Tétanos]
    \textbf{Énoncé :} Deux maladies bactériennes au Cameroun : le Choléra (\emph{Vibrio cholerae}) et le Tétanos (\emph{Clostridium tetani}).
    \begin{enumerate}
        \item Pour chaque maladie, indiquer la voie de contamination principale et le portail d'entrée.
        \item Décrire les barrières naturelles non spécifiques que le microbe doit franchir pour provoquer la maladie.
        \item \emph{V. cholerae} est sensible à l'acidité. Expliquer comment il survit à l'estomac pour atteindre l'intestin.
        \item \emph{C. tetani} est une anaérobie stricte. Expliquer pourquoi une plaie profonde et nécrosée favorise son développement.
        \item Citer la prévention vaccinale efficace pour chacune (nom du vaccin, type d'antigène).
    \end{enumerate}
    \tcblower
    \textbf{Solution détaillée :}
    \begin{enumerate}
        \item \textbf{Choléra :} Voie \textbf{digestive} (ingestion d'eau/aliments souillés par des fèces). Portail d'entrée : \textbf{Muqueuse intestinale} (iléum/côlon).
              \textbf{Tétanos :} Voie \textbf{cutanée-muqueuse} (percée). Portail d'entrée : \textbf{Plaie} (souillée par de la terre/poussière contenant des spores).
        \item \textbf{Barrières franchies :}
              \begin{itemize}
                  \item \textbf{Choléra :} 1. Barrière buccale (salive, lysozyme). 2. \textbf{Barrière gastrique (acidité pH 1-2)} $\leftarrow$ \textbf{Principale}. 3. Barrière intestinale (mucus, flore commensale, péristaltisme, IgA sécrétoires). 4. Pénétration dans les cryptes / fixation sur entérocytes.
                  \item \textbf{Tétanos :} 1. \textbf{Barrière cutanée (peau intacte, kératine, pH acide, flore)} $\leftarrow$ \textbf{Principale franchie par la plaie}. 2. Défense inflammatoire locale (phagocytes, $O_2$). 3. Diffusion lymphatique/sanguine de la toxine (tétanospasmine).
              \end{itemize}
        \item \textbf{Survie de V. cholerae à l'estomac :}
              \begin{itemize}
                  \item \textbf{Dose infectante élevée} ($10^6$ à $10^8$ bactéries) : « effet de nombre », certaines survivent.
                  \item \textbf{Protection par le véhicule alimentaire :} Les aliments (riz, mil, maïs bouillis tièdes) tamponnent l'acidité gastrique, protégeant les bactéries pendant le transit gastrique.
                  \item \textbf{Adaptation physiologique :} Synthèse de systèmes tampons intracellulaires (consommation de $H^+$), modification des lipides de membrane pour résister au pH bas.
              \end{itemize}
        \item \textbf{Plaie profonde + nécrose = Anaérobiose :}
              \begin{itemize}
                  \item Plaie profonde $\rightarrow$ mauvaise vascularisation $\rightarrow$ faible apport en $O_2$.
                  \item Nécrose (tissus morts) $\rightarrow$ consommation d'$O_2$ résiduel par la flore aérobie / fermentation $\rightarrow$ $Eh$ (potentiel redox) très bas.
                  \item \emph{Clostridium tetani} forme des \textbf{spores} (résistantes) qui \textbf{germent} uniquement en \textbf{anaérobiose stricte} $\rightarrow$ multiplication $\rightarrow$ production de \textbf{tétanospasmine} (neurotoxine bloquant la glycine/GABA $\rightarrow$ contractures tétaniques).
              \end{itemize}
        \item \textbf{Prévention vaccinale :}
              \begin{itemize}
                  \item \textbf{Choléra :} Vaccin oral (ex: \emph{Shanchol}, \emph{Euvichol}) : \textbf{Bactéries entières inactivées (tuées)} + sous-unité B de la toxine cholérique (WC-rBS). Immunité muqueuse (IgA).
                  \item \textbf{Tétanos :} Vaccin \textbf{Anatoxine tétanique} (Toxoïde = toxine détotoxiquée par formol, conservant l'immunogénicité). Inclus dans le vaccin combiné \textbf{DTC/HepB/Hib} (Pentavalent) au PEV (Programme Élargi de Vaccination) au Cameroun. Rappels tous les 10-20 ans.
              \end{itemize}
    \end{enumerate}
    \textbf{Conclusion :} La connaissance de la voie de pénétration dicte la stratégie de prévention : \textbf{hygiène eau/aliments/mains} (choléra, fécal-oral) vs \textbf{propreté plaies + vaccination anatoxine} (tétanos, cutané). La virulence repose sur l'adaptation aux barrières (résistance acide, sporulation anaérobie).
\end{exoresolu}

\begin{methode}[Rédiger un compte-rendu d'expérience / Argumenter une conclusion (Savoir-faire transversal)]
    \textbf{Objectif :} Structurer une réponse écrite scientifique rigoureuse (format TP ou question de cours type BEPC).
    \begin{enumerate}
        \item \textbf{Titre :} Précis (ex: « Influence de la température sur la croissance de \emph{E. coli} »).
        \item \textbf{Problématique / Question :} Une phrase interrogative (ex: « La température influence-t-elle la vitesse de division bactérienne ? »).
        \item \textbf{Hypothèse :} Réponse anticipée justifiée (ex: « Oui, car les enzymes bactériennes ont un optimum thermique... »).
        \item \textbf{Matériel \& Méthode :} Liste + Protocole numéroté (passif ou infinitif). Indiquer \textbf{témoin}, \textbf{variables} (indépendante, dépendante, contrôlées), \textbf{répétitions}.
        \item \textbf{Résultats :} \textbf{Tableau} (unités en en-tête) + \textbf{Graphique} (courbe, axes légendés, unités, titre). Description textuelle des tendances (croissance, phase de latence, stationnaire, mort).
        \item \textbf{Interprétation / Discussion :} Expliquer les résultats par les connaissances (cinétique enzymatique, dénaturation, substrats limitants). Comparer au témoin.
        \item \textbf{Conclusion :} Réponse \textbf{affirmative et définitive} à la problématique. « L'hypothèse est validée/invalidée. La température optimale est de 37°C... »
        \item \textbf{Critères de rédaction :} Vocabulaire précis (pas « ça bout » mais « dégagement gazeux »), phrases complètes, passé composé pour les faits observés, présent pour les lois scientifiques, pas de « je » (style impersonnel).
    \end{enumerate}
\end{methode}

\begin{exoresolu}[Rédaction scientifique : Expérience pâte à pain]
    \textbf{Énoncé :} On prépare deux béchers identiques contenant : 100 g farine, 60 mL eau, 10 g levure boulangère.
    \begin{itemize}
        \item Bécher A : Incubé à 30°C.
        \item Bécher B : Incubé à 5°C (réfrigérateur).
    \end{itemize}
    On mesure la hauteur de la pâte toutes les 30 min pendant 3h.
    \textbf{Rédigez le compte-rendu structuré de cette expérience} en insistant sur l'interprétation biochimique des résultats attendus.
    \tcblower
    \textbf{Compte-rendu type :}
    \textbf{Titre :} Influence de la température sur l'activité fermentaire de \emph{Saccharomyces cerevisiae} (levée de la pâte).
    \textbf{Problématique :} La température influence-t-elle la vitesse de fermentation alcoolique (production de $CO_2$) par la levure ?
    \textbf{Hypothèse :} La fermentation est plus rapide à 30°C qu'à 5°C car l'activité enzymatique de la glycolyse et de la fermentation est température-dépendante (optimum vers 30-35°C pour cette levure mésophile).
    \textbf{Matériel :} 2 béchers (500 mL), balance, farine, eau, levure fraîche, règle graduée, étuve 30°C, réfrigérateur 5°C, chronomètre.
    \textbf{Variables :}
    \begin{itemize}
        \item Indépendante : Température (30°C vs 5°C).
        \item Dépendante : Hauteur de la pâte (cm) / Volume de $CO_2$ produit.
        \item Contrôlées : Masse farine, volume eau, masse levure, temps, récipient.
    \end{itemize}
    \textbf{Protocole :}
    1. Mélanger ingrédients dans chaque bécher jusqu'à pâte homogène.
    2. Niveler la surface, noter hauteur initiale $h_0$ (t=0).
    3. Placer Bécher A à 30°C, Bécher B à 5°C.
    4. Mesurer hauteur $h$ toutes les 30 min pendant 180 min.
    \textbf{Résultats attendus (Tableau) :}
    \begin{center}
    \begin{tabular}{|c|c|c|}\hline
    \textbf{Temps (min)} & \textbf{Hauteur A (30°C) (cm)} & \textbf{Hauteur B (5°C) (cm)} \\ \hline
    0 & 5,0 & 5,0 \\ \hline
    30 & 6,2 & 5,1 \\ \hline
    60 & 8,5 & 5,2 \\ \hline
    90 & 11,0 & 5,3 \\ \hline
    120 & 13,5 & 5,4 \\ \hline
    150 & 14,0 (plateau) & 5,5 \\ \hline
    180 & 14,0 & 5,5 \\ \hline
    \end{tabular}
    \end{center}
    \textbf{Graphique :} Courbe A : phase de latence courte, croissance exponentielle forte, phase stationnaire (épuisement substrat/éthanol toxique). Courbe B : quasi plate (légère hausse due dilatation thermique/activité résiduelle).
    \textbf{Interprétation :}
    À 30°C (optimum mésophile), les enzymes de la glycolyse (hexokinase, PFK, pyruvate kinase) et de la fermentation (pyruvate décarboxylase, alcool déshydrogénase) ont une \textbf{vitesse catalytique maximale}. Le glucose est rapidement converti en pyruvate puis en éthanol + $CO_2$. Le $CO_2$ (gaz) se dissout puis forme des bulles piégées dans le réseau de gluten $\rightarrow$ augmentation volume/hauteur.
    À 5°C, l'énergie cinétique des molécules est faible, les collisions enzyme-substrat sont rares $\rightarrow$ \textbf{vitesse réactionnelle très faible} (règle de Van't Hoff : vitesse divisée par 2 à 3 tous les 10°C). La levure est viable mais métaboliquement « endormie ».
    \textbf{Conclusion :} L'hypothèse est \textbf{validée}. La température de 30°C favorise une fermentation alcoolique rapide et massive (levée de la pâte), tandis que 5°C la ralentit drastiquement (conservation de la pâte/levure). Cela explique les pratiques boulangères (pousse à température ambiante/étuve, rétention au froid pour ralentir).
\end{exoresolu}

\begin{attention}[Les 5 erreurs qui coûtent des points au BEPC]
    \begin{enumerate}
        \item \textbf{Confondre « Reproduction » et « Nutrition/Respiration » :} Dire « la levure se reproduit par fermentation » (\textbf{Faux}). La fermentation est son mode de \textbf{respiration anaérobie} (production d'énergie). Elle se \textbf{reproduit par bourgeonnement} (mitose). De même, « les bactéries se reproduisent par conjugaison » (\textbf{Faux}) : la conjugaison est un \textbf{transfert génétique}, la reproduction est la \textbf{division binaire}.
        \item \textbf{Oublier les unités ou l'échelle au dessin microscopique :} Un dessin sans grossissement ($G_t = \times 400$) ni barre d'échelle vaut 0 point pour la partie observation. Dessiner au stylo, colorier, faire des traits épais/hachurés = pénalité.
        \item \textbf{Dire « Virus = Microbe vivant / Cellule » :} Le virus est une \textbf{particule acellulaire}, \textbf{inerte} hors cellule hôte. Il ne se « nourrit » pas, ne « respire » pas, ne « se divise » pas. Il se \textbf{multiplie} (réplication) en détournant la machinerie cellulaire. Écrire « cellule virale » = erreur majeure.
        \item \textbf{Mélanger voies de contamination et vecteurs :} Le moustique n'est pas une « voie de contamination », c'est un \textbf{vecteur} (voie cutanée par piqûre). L'eau sale n'est pas une voie, c'est un \textbf{véhicule} (voie digestive par ingestion). Voies = Respiratoire, Digestive, Cutanée-muqueuse, Verticale.
        \item \textbf{Confondre Anaérobie stricte et Anaérobie facultative / Aérobie tolérant :} \emph{Lactobacillus} ne respire \textbf{jamais} (pas de cytochromes), c'est un \textbf{aérobie tolérant} (fermentation permanente). \emph{E. coli} / \emph{Levure} respirent \textbf{si $O_2$ présent}, fermentent \textbf{si $O_2$ absent} = \textbf{Anaérobies facultatives}. \emph{Clostridium} / Méthanogènes : $O_2$ toxique = \textbf{Anaérobies strictes}. Ne pas écrire « les bactéries anaérobies meurent à l'air » sans préciser le groupe.
    \end{enumerate}
\end{attention}

\newpage

\coursec{Exercices}

\courssub{Je vérifie mes connaissances}

\begin{exercice}[QCM : Groupes de microorganismes]
Parmi les propositions suivantes, coche la ou les bonne(s) réponse(s) :
\begin{enumerate}
\item Les microorganismes sont : \textbf{a)} tous pathogènes \textbf{b)} visibles à l'œil nu \textbf{c)} des êtres vivants microscopiques \textbf{d)} uniquement des bactéries.
\item Le groupe de microorganismes dont les cellules ne possèdent pas de noyau individualisé est : \textbf{a)} les champignons microscopiques \textbf{b)} les bactéries \textbf{c)} les protistes \textbf{d)} les virus.
\item Un virus est constitué de : \textbf{a)} une membrane, du cytoplasme et un noyau \textbf{b)} une capside protéique et un acide nucléique (ADN ou ARN) \textbf{c)} une paroi cellulaire et des ribosomes \textbf{d)} des mitochondries et un ADN circulaire.
\item La levure de boulanger (\emph{Saccharomyces cerevisiae}) appartient au groupe : \textbf{a)} bactéries \textbf{b)} virus \textbf{c)} champignons microscopiques \textbf{d)} protistes.
\end{enumerate}
\end{exercice}

\begin{exercice}[Vrai / Faux justifié]
Indique si chaque affirmation est vraie (V) ou fausse (F) et justifie brièvement ta réponse.
\begin{enumerate}
\item Tous les microorganismes se reproduisent par division cellulaire (scissiparité).
\item La nutrition des bactéries est toujours autotrophe.
\item Les virus sont capables de se multiplier en dehors d'une cellule hôte.
\item La contamination par voie digestive se fait par l'ingestion d'aliments ou d'eau souillés.
\item Le microscope optique permet d'observer les virus.
\end{enumerate}
\end{exercice}

\begin{exercice}[Définitions et vocabulaire]
Donne la définition précise des termes suivants :
\begin{enumerate}
\item Microorganisme
\item Pathogène
\item Scissiparité
\item Fermentation alcoolique
\item Voie de pénétration (ou porte d'entrée) d'un microorganisme
\item Milieu de culture
\end{enumerate}
\end{exercice}

\begin{exercice}[Calcul direct : Multiplication bactérienne]
Une bactérie se divise toutes les 20 minutes dans des conditions optimales. On place une seule bactérie dans un milieu de culture favorable à \SI{8}{h} du matin.
\begin{enumerate}
\item Combien de divisions auront lieu à \SI{10}{h} ?
\item Calcule le nombre de bactéries présentes à \SI{10}{h}.
\item Si la division se poursuit au même rythme, combien y aura-t-il de bactéries à \SI{12}{h} ? (Donne le résultat sous forme de puissance de 2).
\end{enumerate}
\end{exercice}

\courssub{Je m'entraîne}

\begin{exercice}[Identification de microorganismes à partir de schémas]
On te présente quatre descriptions (A, B, C, D) de microorganismes observés au microscope.
\begin{center}
\begin{tabularx}{\linewidth}{|c|X|X|}\hline
\textbf{Schéma} & \textbf{Description} & \textbf{Observations notables} \\ \hline
A & Cellule unique, forme sphérique, pas de noyau visible, paroi cellulaire fine & Groupes de deux (diplocoques) ou chaînes \\ \hline
B & Cellule unique, forme ovale, noyau visible, vacuole, paroi cellulaire épaisse & Bourgeonnement visible sur certaines cellules \\ \hline
C & Cellule unique de forme variable, noyau visible, pseudopodes & Déplacement par pseudopodes \\ \hline
D & Particules très petites (environ \SI{100}{nm}), structure géométrique régulière & Visibles seulement au microscope électronique \\ \hline
\end{tabularx}
\end{center}
\begin{enumerate}
\item Pour chaque schéma, nomme le groupe de microorganismes correspondant (bactérie, champignon microscopique, protiste, virus).
\item Justifie ton classement pour chaque schéma en citant \textbf{deux} caractères observables.
\item Lequel de ces microorganismes ne peut pas être observé au microscope optique scolaire ? Justifie.
\end{enumerate}
\end{exercice}

\begin{exercice}[Expérience de Pasteur : génération spontanée ou contamination aérienne ?]
Louis Pasteur a réalisé l'expérience suivante : il a mis du bouillon nutritif dans deux flacons à col de cygne et l'a fait bouillir dans les deux flacons. Le flacon 1 est resté intact. Le col du flacon 2 a été cassé après ébullition. Après quelques jours, le bouillon du flacon 2 est trouble, celui du flacon 1 reste clair.
\begin{enumerate}
\item Quelle hypothèse Pasteur testait-il avec cette expérience ?
\item Pourquoi le col de cygne du flacon 1 a-t-il été laissé intact ?
\item Interprète la turbidité (trouble) du bouillon dans le flacon 2.
\item Quelle conclusion tire-t-on sur l'origine des microorganismes dans le bouillon ?
\item Relie cette expérience à la stérilisation des conserves alimentaires (ex. : confitures, sauces tomate au Cameroun).
\end{enumerate}
\end{exercice}

\begin{exercice}[Fermentation alcoolique : levure et eau de chaux]
On introduit une suspension de levures (\emph{Saccharomyces cerevisiae}) dans un tube à essai contenant de l'eau sucrée (sirop de canne dilué). On bouche le tube avec un bouchon percé muni d'un tube plongeant dans un second tube contenant de l'eau de chaux. Au bout de \SI{30}{min}, l'eau de chaux devient trouble.
\begin{enumerate}
\item Nomme le gaz responsable du trouble de l'eau de chaux. Écris l'équation chimique de la réaction avec l'eau de chaux (solution de \ce{Ca(OH)2}).
\item Quel est le nom du phénomène métabolique réalisé par les levures en l'absence de dioxygène ?
\item Écris l'équation bilan de ce phénomène (substrat $\rightarrow$ produits + énergie).
\item Cite \textbf{deux} applications alimentaires traditionnelles qui utilisent ce phénomène.
\item Que se passerait-il si on répétait l'expérience en milieu bien aéré ? Justifie.
\end{enumerate}
\end{exercice}

\begin{exercice}[Compte-rendu de TP : observation de levures au microscope]
Lors du TP N°5, tu observes une goutte de suspension de levure vivante entre lame et lamelle au microscope optique (objectif $\times 10$, oculaire $\times 10$).
\begin{enumerate}
\item Calcule le grossissement total utilisé.
\item Décris ce que tu observes : forme, structures visibles (noyau, vacuole, paroi, bourgeonnement), mouvement.
\item Dessine et légende \textbf{une} cellule de levure en phase de bourgeonnement (traits nets, légendes alignées à droite).
\item Quelle précaution prends-tu pour éviter que la préparation ne sèche trop vite ?
\item Conclus sur le mode de reproduction des levures.
\end{enumerate}
\end{exercice}

\courssub{Je me prépare au BEPC}

\begin{exercice}[Cas concret : épidémie de diarrhée dans un village]
Dans un village de la région de l'Extrême-Nord, plusieurs enfants de moins de 5 ans présentent des diarrhées aqueuses abondantes, des vomissements et une déshydratation rapide après avoir bu l'eau d'un puits non protégé, situé à proximité d'une zone de défécation à ciel ouvert. Le centre de santé signale une augmentation brutale des cas en 3 jours.
\begin{enumerate}
\item Identifie la \textbf{voie de contamination} la plus probable dans ce cas. Justifie par deux arguments tirés du texte.
\item Parmi les groupes de microorganismes (bactéries, virus, protistes), lesquels sont le plus souvent responsables de telles épidémies d'origine hydrique ? Donne \textbf{deux} exemples d'agents pathogènes précis (nom de genre ou d'espèce).
\item Le personnel de santé préconise le traitement de l'eau par ébullition ou par pastilles de chlore. Explique le \textbf{principe} de chaque méthode sur les microorganismes.
\item Rédige un \textbf{message de prévention} court (5 à 6 lignes) à destination des familles du village, en utilisant un langage simple mais scientifiquement juste. Intègre les notions de : voie de pénétration, traitement de l'eau, hygiène des mains, consultation précoce.
\item Pourquoi la déshydratation est-elle le danger principal chez l'enfant ? Cite la solution recommandée par l'OMS pour la combattre.
\end{enumerate}
\end{exercice}

\begin{exercice}[Exploitation d'une courbe de croissance bactérienne]
On a ensemencé un milieu de culture liquide stérile avec \num{100} bactéries \emph{Escherichia coli} et on a compté le nombre de bactéries vivantes (UFC/mL) toutes les \SI{20}{min} à \SI{37}{\celsius}. Résultats :
\begin{center}
\begin{tabular}{|c|c|c|c|c|c|c|c|}\hline
Temps (min) & 0 & 20 & 40 & 60 & 80 & 100 & 120 \\ \hline
Nombre de bactéries (UFC/mL) & \num{100} & \num{200} & \num{400} & \num{800} & \num{1600} & \num{3200} & \num{6400} \\ \hline
\end{tabular}
\end{center}
\begin{enumerate}
\item Représente graphiquement cette courbe (abscisses : temps en min ; ordonnées : nombre de bactéries). Indique les unités.
\item Que remarques-tu sur l'allure de la courbe ? Comment appelle-t-on ce type de croissance ?
\item Calcule le \textbf{temps de génération} (durée d'une division). Montre ton raisonnement.
\item Combien de bactéries y aura-t-il théoriquement après \SI{3}{h} (\SI{180}{min}) si la croissance se poursuit au même rythme ? Exprime le résultat en notation scientifique.
\item Dans la réalité, la population ne croît pas indéfiniment. Cite \textbf{deux} facteurs limitants qui ralentissent puis arrêtent la croissance.
\item Un plat de \emph{ndolé} cuisiné le matin est laissé à température ambiante (environ \SI{30}{\celsius}) jusqu'au soir (\SI{10}{h} plus tard). En te basant sur le temps de génération calculé, explique pourquoi ce plat présente un risque sanitaire élevé.
\end{enumerate}
\end{exercice}

\begin{exercice}[Synthèse : microorganismes, santé et alimentation au Cameroun]
Le document ci-dessous résume trois situations courantes :
\begin{quote}
\textbf{Situation 1 :} Une mère prépare du \emph{foléré} (jus d'hibiscus) maison. Elle laisse la décoction refroidir toute la nuit dans une casserole non couverte. Le lendemain, le liquide est trouble, mousseux et dégage une odeur alcoolisée.

\textbf{Situation 2 :} Un agriculteur de la région du Centre utilise du compost bien décomposé pour fertiliser son champ de maïs. Il observe la présence de nombreux vers de terre et une bonne structure du sol.

\textbf{Situation 3 :} Un élève se blesse au pied en jouant pieds nus. La plaie n'est pas lavée ni désinfectée. Deux jours plus tard, le pied est rouge, chaud, gonflé et douloureux. Du pus s'écoule.
\end{quote}
\begin{enumerate}
\item \textbf{Situation 1 :}
      \begin{enumerate}
      \item Quels microorganismes sont responsables de la transformation du jus ? Nomme le groupe et l'espèce la plus probable.
      \item Quel est le nom du phénomène métabolique ? Écris son équation bilan.
      \item Cette boisson est-elle propre à la consommation ? Justifie.
      \item Comment éviter cette contamination ? (Deux mesures concrètes).
      \end{enumerate}
\item \textbf{Situation 2 :}
      \begin{enumerate}
      \item Quel rôle les microorganismes (bactéries, champignons) jouent-ils dans la formation du compost ? Cite le terme scientifique.
      \item Pourquoi la présence de vers de terre est-elle un indicateur de bonne qualité du sol ?
      \end{enumerate}
\item \textbf{Situation 3 :}
      \begin{enumerate}
      \item Identifie la voie de pénétration des microorganismes.
      \item Les symptômes (rougeur, chaleur, gonflement, douleur, pus) correspondent à quelle réaction de l'organisme ? Nomme-la.
      \item Cite \textbf{deux} genres de bactéries fréquemment responsables d'infections de plaies sales.
      \item Quels gestes de premiers secours auraient dû être faits immédiatement après la blessure ? (Trois actions dans l'ordre).
      \end{enumerate}
\item \textbf{Bilan transversal :} Complète le tableau récapitulatif suivant :
\begin{center}
\begin{tabularx}{\linewidth}{|l|X|X|X|}\hline
\textbf{Domaine} & \textbf{Microorganismes utiles (exemple + rôle)} & \textbf{Microorganismes nuisibles (exemple + pathologie)} & \textbf{Moyen de lutte / prévention} \\ \hline
Alimentation & & & \\ \hline
Santé & & & \\ \hline
Agriculture / Environnement & & & \\ \hline
\end{tabularx}
\end{center}
\end{enumerate}
\end{exercice}



\newpage

\coursec{Corrigés des exercices}

\courssub{Je vérifie mes connaissances}

\begin{corrige}[Exercice 1 — QCM : Groupes de microorganismes]
\begin{enumerate}
\item \textbf{c)} : les microorganismes sont des êtres vivants microscopiques, invisibles à l'œil nu. Ils ne sont pas tous pathogènes (beaucoup sont utiles ou inoffensifs) et ne se limitent pas aux bactéries.
\item \textbf{b)} : les bactéries sont des cellules sans noyau individualisé (leur ADN est libre dans le cytoplasme). Les champignons et les protistes possèdent un noyau ; les virus ne sont pas des cellules.
\item \textbf{b)} : un virus est formé d'une capside (enveloppe protéique) renfermant un acide nucléique, ADN ou ARN. Il n'a ni cytoplasme, ni organites.
\item \textbf{c)} : la levure de boulanger est un champignon microscopique unicellulaire.
\end{enumerate}
\textbf{Points de vigilance :} ne pas confondre « virus » et « cellule » : un virus n'a ni noyau, ni cytoplasme, il n'est donc pas une cellule.
\end{corrige}

\begin{corrige}[Exercice 2 — Vrai / Faux justifié]
\begin{enumerate}
\item \textbf{F} : la scissiparité est le mode de reproduction des bactéries ; les levures se reproduisent par bourgeonnement et les virus en forçant une cellule hôte à fabriquer de nouveaux virus.
\item \textbf{F} : la plupart des bactéries sont hétérotrophes (elles se nourrissent de matière organique) ; seules quelques-unes sont autotrophes.
\item \textbf{F} : un virus est un parasite obligatoire : il ne peut se multiplier qu'à l'intérieur d'une cellule vivante.
\item \textbf{V} : la voie digestive est la pénétration des microbes par la bouche, avec de l'eau ou des aliments contaminés.
\item \textbf{F} : les virus (environ \SI{100}{nm}) sont trop petits pour le microscope optique ; ils ne sont visibles qu'au microscope électronique.
\end{enumerate}
\end{corrige}

\begin{corrige}[Exercice 3 — Définitions et vocabulaire]
\begin{enumerate}
\item \textbf{Microorganisme} : être vivant de taille microscopique, invisible à l'œil nu, observable seulement au microscope.
\item \textbf{Pathogène} : qualifie un microorganisme capable de provoquer une maladie chez un hôte.
\item \textbf{Scissiparité} : mode de reproduction asexuée dans lequel une cellule mère se divise en deux cellules filles identiques (cas des bactéries).
\item \textbf{Fermentation alcoolique} : transformation du glucose, en l'absence de dioxygène, en éthanol et en dioxyde de carbone, avec libération d'énergie ; elle est réalisée par les levures.
\item \textbf{Voie de pénétration} : chemin par lequel un microorganisme entre dans l'organisme (voie digestive, respiratoire, cutanée par une plaie, etc.).
\item \textbf{Milieu de culture} : substance nutritive (liquide ou gélifiée) sur laquelle on fait se développer des microorganismes au laboratoire.
\end{enumerate}
\end{corrige}

\begin{corrige}[Exercice 4 — Calcul direct : Multiplication bactérienne]
\begin{enumerate}
\item De \SI{8}{h} à \SI{10}{h}, il s'écoule \SI{120}{min}. Nombre de divisions : $\dfrac{120}{20} = 6$ divisions.
\item En partant d'une seule bactérie à \SI{8}{h}, chaque division double le nombre de bactéries : $N = 2^{6} = 64$ bactéries à \SI{10}{h}.
\item De \SI{8}{h} à \SI{12}{h}, il s'écoule \SI{240}{min}, soit $\dfrac{240}{20} = 12$ divisions. Le nombre de bactéries est $N = 2^{12} = \num{4096}$ bactéries.
\end{enumerate}
\textbf{Points de vigilance :} bien compter le temps écoulé depuis le début (\SI{8}{h}) et non depuis \SI{10}{h} pour la question 3 ; vérifier que $2^{12} = 2 \times 2^{11} = \num{4096}$.
\end{corrige}

\courssub{Je m'entraîne}

\begin{corrige}[Exercice 5 — Identification de microorganismes à partir de schémas]
\begin{enumerate}
\item A : bactérie ; B : champignon microscopique (levure) ; C : protiste (amibe) ; D : virus.
\item A : absence de noyau visible et cellules sphériques groupées en paires ou en chaînes. B : noyau visible et bourgeonnement (reproduction typique des levures). C : noyau visible et présence de pseudopodes servant au déplacement. D : taille d'environ \SI{100}{nm} et invisibilité au microscope optique.
\item Le microorganisme D (virus) ne peut pas être observé au microscope optique scolaire : sa taille (environ \SI{100}{nm}) est inférieure au pouvoir de résolution du microscope optique ; il faut un microscope électronique.
\end{enumerate}
\end{corrige}

\begin{corrige}[Exercice 6 — Expérience de Pasteur]
\begin{enumerate}
\item Pasteur testait l'hypothèse de la \cle{génération spontanée}, selon laquelle les microorganismes naîtraient spontanément dans le bouillon.
\item Le col de cygne intact laisse passer l'air mais piège les poussières et les microbes dans ses coudes : le bouillon reste en contact avec l'air sans être contaminé. C'est le témoin qui montre que l'air seul ne trouble pas le bouillon.
\item Le trouble du bouillon du flacon 2 traduit le développement d'une population de microorganismes qui s'y sont multipliés.
\item Les microorganismes ne naissent pas spontanément dans le bouillon : ils proviennent de l'air et des poussières qui pénètrent dans le flacon. La génération spontanée est réfutée.
\item Faire bouillir les aliments tue les microorganismes (stérilisation) ; fermer hermétiquement les bocaux empêche toute nouvelle contamination venue de l'air. C'est le principe des conserves (confitures, sauces tomate) qui se conservent longtemps sans se gâter.
\end{enumerate}
\textbf{Points de vigilance :} bien distinguer le rôle du témoin (flacon 1) du dispositif expérimental (flacon 2) ; la conclusion doit porter sur l'origine des microbes (l'air), pas seulement sur le trouble.
\end{corrige}

\begin{corrige}[Exercice 7 — Fermentation alcoolique : levure et eau de chaux]
\begin{enumerate}
\item Le gaz est le \cle{dioxyde de carbone} \ce{CO2}. Il trouble l'eau de chaux en formant un précipité blanc de carbonate de calcium :
\[ \ce{CO2 + Ca(OH)2 -> CaCO3 + H2O} \]
\item Il s'agit de la \cle{fermentation alcoolique}, une respiration anaérobie (sans dioxygène) réalisée par les levures.
\item Équation bilan :
\[ \ce{C6H12O6 -> 2 C2H5OH + 2 CO2} + \text{énergie} \]
\item La fabrication de boissons fermentées traditionnelles (vin de palme, \emph{bil-bil} à base de mil) et la panification (la pâte lève grâce au \ce{CO2} dégagé).
\item En milieu aéré, les levures réalisent la \cle{respiration aérobie} : le glucose est entièrement dégradé en \ce{CO2} et en eau, sans production d'alcool. L'eau de chaux se trouble encore (du \ce{CO2} est toujours produit), mais il n'y a pas d'odeur alcoolisée.
\end{enumerate}
\textbf{Points de vigilance :} ne pas écrire que le trouble disparaît en milieu aéré : le \ce{CO2} est produit dans les deux cas ; seule la production d'alcool disparaît.
\end{corrige}

\begin{corrige}[Exercice 8 — Compte-rendu de TP : observation de levures au microscope]
\begin{enumerate}
\item Grossissement total : $G = 10 \times 10 = 100$, soit un grossissement de $\times 100$.
\item On observe des cellules ovales, isolées ou groupées, entourées d'une paroi, avec un cytoplasme granuleux, une vacuole claire et parfois un noyau. Certaines cellules portent un petit bourgeon. Les levures ne se déplacent pas (pas de mouvement propre).
\item Schéma attendu : cellule ovale avec un bourgeon plus petit accolé ; légendes : paroi, cytoplasme, noyau, vacuole, bourgeon.
\begin{popfigure}
\begin{tikzpicture}[scale=1.8, every node/.style={font=\small}]
% Paroi cellule mère
\draw[popdark, thick] (0,0) ellipse (1.2 and 0.7);
% Cytoplasme mère
\fill[popcream] (0,0) ellipse (1.15 and 0.65);
% Noyau mère
\draw[popblue, thick] (-0.4,0.15) circle (0.12);
\fill[popblue!30] (-0.4,0.15) circle (0.12);
% Vacuole mère
\draw[popblue, thick] (0.5,-0.15) circle (0.18);
% Bourgeon
\draw[popdark, thick] (1.15,0.35) ellipse (0.45 and 0.35);
\fill[popcream] (1.15,0.35) ellipse (0.4 and 0.3);
% Noyau bourgeon
\draw[popblue, thick] (1.25,0.45) circle (0.07);
\fill[popblue!30] (1.25,0.45) circle (0.07);
% Légendes
\node[align=left, anchor=west, popdark] at (1.8, 0.9) {1. Paroi cellulaire};
\node[align=left, anchor=west, popdark] at (1.8, 0.6) {2. Cytoplasme granuleux};
\node[align=left, anchor=west, popdark] at (1.8, 0.3) {3. Noyau};
\node[align=left, anchor=west, popdark] at (1.8, 0.0) {4. Vacuole};
\node[align=left, anchor=west, popdark] at (1.8, -0.3) {5. Bourgeon};
\node[align=left, anchor=west, popdark] at (1.8, -0.6) {6. Noyau du bourgeon};
\end{tikzpicture}
\legende{Schéma d'une levure en phase de bourgeonnement (observation idéalisée, grossissement $\times 400$ recommandé pour les détails).}
\end{popfigure}
\item On dépose une goutte suffisante de suspension et on pose la lamelle en évitant les bulles ; on peut aussi sceller les bords de la lamelle ou ré-ajouter une goutte de liquide si nécessaire.
\item Les levures se reproduisent par \cle{bourgeonnement} : un petit bourgeon se forme sur la cellule mère, grossit, puis se détache pour donner une nouvelle cellule.
\end{enumerate}
\textbf{Points de vigilance :} le grossissement total est le produit des grossissements de l'oculaire et de l'objectif, pas leur somme ; le dessin doit être fait au crayon, traits nets, légendes alignées à droite reliées par des traits horizontaux.
\end{corrige}

\courssub{Je me prépare au BEPC}

\begin{corrige}[Exercice 9 — Cas concret : épidémie de diarrhée dans un village]
\begin{enumerate}
\item Il s'agit de la \cle{voie digestive} (eau et aliments souillés). Arguments : les malades ont bu l'eau d'un puits non protégé ; ce puits est proche d'une zone de défécation à ciel ouvert, donc l'eau est contaminée par des microbes d'origine fécale.
\item Ce sont surtout des \textbf{bactéries} et des \textbf{virus}. Exemples : \emph{Vibrio cholerae} (choléra) et \emph{Salmonella typhi} (fièvre typhoïde) ; on peut aussi citer les amibes (protistes) responsables de l'amibiase.
\item \textbf{Ébullition} : la chaleur (\SI{100}{\celsius}) détruit les microorganismes présents dans l'eau. \textbf{Chloration} : le chlore est un désinfectant qui tue les microbes au contact de l'eau.
\item Exemple de message : « Les microbes de la diarrhée entrent dans le corps par la bouche, avec l'eau et les aliments souillés. Faites bouillir l'eau de boisson ou traitez-la avec des pastilles de chlore. Lavez-vous les mains avec du savon avant de manger et après être allé aux toilettes. Couvrez les aliments. Dès les premiers signes de diarrhée, donnez une solution de réhydratation et amenez vite l'enfant au centre de santé. »
\item La diarrhée et les vomissements font perdre beaucoup d'eau et de sels minéraux ; chez le jeune enfant, cette perte est rapide et peut entraîner la mort. L'OMS recommande la \cle{solution de réhydratation orale} (SRO), un mélange d'eau, de sel et de sucre.
\end{enumerate}
\textbf{Barème indicatif (10 points) :} voie de contamination justifiée par 2 arguments (2 pts) ; groupes + 2 exemples d'agents (2 pts) ; principe des deux traitements de l'eau (2 pts) ; message de prévention complet et correct (2 pts) ; déshydratation + SRO (2 pts).\\
\textbf{Points de vigilance :} les arguments doivent être tirés du texte (puits non protégé, défécation à proximité) ; le message doit intégrer les quatre notions demandées (voie de pénétration, traitement de l'eau, hygiène des mains, consultation précoce).
\end{corrige}

\begin{corrige}[Exercice 10 — Exploitation d'une courbe de croissance bactérienne]
\begin{enumerate}
\item Courbe : points $(0\,;\,100)$, $(20\,;\,200)$, $(40\,;\,400)$, $(60\,;\,800)$, $(80\,;\,1600)$, $(100\,;\,3200)$, $(120\,;\,6400)$, reliés par une courbe qui monte de plus en plus vite. Axes : temps (min) en abscisse, nombre de bactéries (UFC/mL) en ordonnée.
\begin{popfigure}
\begin{tikzpicture}
\begin{axis}[
    width=\linewidth,
    height=0.55\linewidth,
    xlabel={Temps (min)},
    ylabel={Nombre de bactéries (UFC/mL)},
    xmin=0, xmax=130,
    ymin=0, ymax=7000,
    xtick={0,20,40,60,80,100,120},
    ytick={0,1000,2000,3000,4000,5000,6000,7000},
    grid=major,
    grid style={popline},
    tick label style={font=\small},
    label style={font=\small},
    clip=false
    ]
\addplot[popcurve, mark=*, mark size=2pt, popblue] coordinates {
    (0,100) (20,200) (40,400) (60,800) (80,1600) (100,3200) (120,6400)
};
\end{axis}
\end{tikzpicture}
\legende{Croissance exponentielle d'une population bactérienne (temps de génération = 20 min).}
\end{popfigure}
\item La courbe monte de plus en plus vite : le nombre de bactéries double à intervalles réguliers. C'est une \cle{croissance exponentielle}.
\item Entre deux mesures successives, le nombre double : $200/100 = 2$, $400/200 = 2$, etc. L'intervalle entre deux mesures est \SI{20}{min}. Le temps de génération est donc de \SI{20}{min}.
\item En \SI{180}{min}, il y a $\dfrac{180}{20} = 9$ divisions. $N = 100 \times 2^{9} = 100 \times 512 = \num{51200} = \num{5,12 \times 10^{4}}$ bactéries/mL.
\item L'épuisement des substances nutritives du milieu et l'accumulation de déchets toxiques (on peut aussi citer le manque d'espace ou de dioxygène).
\item À environ \SI{30}{\celsius}, les bactéries se divisent rapidement (de l'ordre de \SI{20}{min} à \SI{30}{min}). En \SI{10}{h}, soit au moins 20 générations, une seule bactérie peut donner plus de $2^{20} \approx 10^{6}$ bactéries. Le plat contient alors des millions de microbes (et leurs toxines), d'où un risque élevé d'intoxication alimentaire. Il faut conserver les plats au frais ou les réchauffer fortement.
\end{enumerate}
\textbf{Barème indicatif (10 points) :} courbe correcte avec axes légendés et unités (2 pts) ; allure et nom de la croissance (1 pt) ; temps de génération justifié (2 pts) ; calcul à \SI{180}{min} en notation scientifique (2 pts) ; deux facteurs limitants (1 pt) ; application au \emph{ndolé} argumentée avec un calcul (2 pts).\\
\textbf{Points de vigilance :} ne pas oublier de multiplier par l'effectif initial (100) : $N = N_0 \times 2^{n}$ ; vérifier $2^{9} = 512$ et $100 \times 512 = \num{51200}$.
\end{corrige}

\begin{corrige}[Exercice 11 — Synthèse : microorganismes, santé et alimentation au Cameroun]
\begin{enumerate}
\item \textbf{Situation 1 :}
      \begin{enumerate}
      \item Des \textbf{champignons microscopiques} : des levures, très probablement \emph{Saccharomyces cerevisiae}, apportées par l'air et les poussières pendant la nuit.
      \item La \cle{fermentation alcoolique} : $\ce{C6H12O6 -> 2 C2H5OH + 2 CO2} + \text{énergie}$. Le \ce{CO2} explique la mousse, l'éthanol l'odeur alcoolisée.
      \item Non : la présence de mousse, de trouble et d'odeur alcoolisée montre que le jus est contaminé et fermenté ; d'autres microbes indésirables ont pu s'y développer.
      \item Couvrir la casserole pendant le refroidissement ; conserver le jus au frais (réfrigérateur) et le consommer rapidement ; faire bouillir à nouveau avant consommation.
      \end{enumerate}
\item \textbf{Situation 2 :}
      \begin{enumerate}
      \item Les bactéries et champignons du sol décomposent la matière organique morte en matière minérale assimilable par les plantes : ce sont des \cle{décomposeurs} (dégradation de la matière organique, humification).
      \item Les vers de terre brassent et aèrent le sol, mélangent l'humus et facilitent le travail des microorganismes décomposeurs : leur présence signale un sol vivant, riche et fertile.
      \end{enumerate}
\item \textbf{Situation 3 :}
      \begin{enumerate}
      \item La \cle{voie cutanée}, par la plaie (porte d'entrée traumatique).
      \item Il s'agit de la \cle{réaction inflammatoire} (inflammation), réponse de défense de l'organisme contre l'infection ; le pus contient des microbes et des cellules de défense mortes.
      \item \emph{Staphylococcus} (staphylocoques) et \emph{Streptococcus} (streptocoques).
      \item (1) Laver la plaie à l'eau propre et au savon ; (2) désinfecter avec un antiseptique ; (3) protéger par un pansement propre et consulter un agent de santé si la plaie est profonde ou sale.
      \end{enumerate}
\item \textbf{Bilan transversal (exemple de tableau complété) :}
\begin{center}
\begin{tabularx}{\linewidth}{|l|X|X|X|}\hline
\textbf{Domaine} & \textbf{Microorganismes utiles} & \textbf{Microorganismes nuisibles} & \textbf{Moyen de lutte / prévention} \\ \hline
Alimentation & Levures : fermentation du pain et des boissons & Bactéries de la pourriture et toxines : intoxications alimentaires & Ébullition, cuisson, conservation au frais, couverture des aliments \\ \hline
Santé & Bactéries de la flore intestinale : aide à la digestion & \emph{Vibrio cholerae} : choléra ; virus de la grippe & Hygiène des mains, eau traitée, désinfection des plaies, vaccination \\ \hline
Agriculture / Environnement & Bactéries et champignons décomposeurs : formation de l'humus et du compost & Champignons parasites des plantes : pourriture des récoltes & Rotation des cultures, compostage, traitement des semences \\ \hline
\end{tabularx}
\end{center}
\end{enumerate}
\textbf{Barème indicatif (20 points) :} situation 1 (5 pts) ; situation 2 (3 pts) ; situation 3 (6 pts) ; tableau bilan complet et cohérent (6 pts).\\
\textbf{Points de vigilance :} relier chaque observation (mousse, odeur, pus) à son interprétation ; dans le tableau, chaque case doit contenir un exemple ET son rôle ou sa pathologie, pas seulement un nom de microbe.
\end{corrige}

\newpage

\coursec{Fiche bilan}

\begin{popfigure}
\centering
\begin{tikzpicture}[
    mindmap,
    grow cyclic,
    every node/.style={concept, execute at begin node=\hskip0pt},
    root concept/.append style={concept color=popblue, minimum size=3.5cm, text=white, font=\bfseries\large},
    level 1 concept/.append style={level distance=4.5cm, sibling angle=360/6, minimum size=2.8cm, text=white, font=\bfseries\normalsize},
    level 2 concept/.append style={level distance=3.5cm, sibling angle=45, minimum size=2cm, text=white, font=\small},
    concept color=popblue!70,
    font=\small
]
\node[root concept, align=center]{Microorganismes\\dans notre\\environnement}
    child[concept color=popgreen] { node[level 1 concept, align=center]{1. Découverte\\et définition} }
    child[concept color=poporange] { node[level 1 concept, align=center]{2. Groupes\\majeurs} }
    child[concept color=poppurple] { node[level 1 concept, align=center]{3. TP 5 : Observation\\au microscope} }
    child[concept color=poppink] { node[level 1 concept, align=center]{4. Nutrition\\et respiration} }
    child[concept color=popgold] { node[level 1 concept] {5. Reproduction} }
    child[concept color=popdark] { node[level 1 concept, align=center]{6. Contamination\\et pénétration} };
\end{tikzpicture}
\legende{Carte mentale récapitulative de la leçon 04 : les 6 grands axes du programme officiel 3ème SVTEEHB.}
\end{popfigure}

\begin{bilan}
\courssub{Essentiel du cours à maîtriser}

\textbf{1. Découverte et définition}
\begin{itemize}
    \item \cle{Microorganisme (microbe)} : organisme vivant de taille microscopique (généralement $< 0,1$ mm), invisible à l'œil nu, observable au microscope optique ou électronique.
    \item Présents partout : air, eau, sol, peau, tube digestif, aliments.
    \item Rôle majeur : décomposition (recyclage matière), fermentation (aliments), pathogènes (maladies), biotechnologies.
\end{itemize}

\textbf{2. Les différents groupes de microorganismes}
\begin{center}
\begin{tabularx}{\linewidth}{|l|X|X|X|}\hline
\rowcolor{popblueL}
\textbf{Groupe} & \textbf{Caractéristiques cellulaires} & \textbf{Exemples} & \textbf{Taille approx.} \\ \hline
\textbf{Bactéries} & Procaryotes (pas de noyau, pas d'organites membranaires), paroi en peptidoglycane, unicellulaires. & \emph{Escherichia coli}, \emph{Streptocoque}, \emph{Lactobacille}, \emph{Salmonelle}. & $1$ à $10\ \mu m$ \\ \hline
\textbf{Virus} & Acellulaires (pas de cellule), ADN \textbf{ou} ARN + capside protéique (+ enveloppe parfois), parasites intracellulaires obligatoires. & VIH, Grippes, Rougeole, COVID-19, Bactériophages. & $20$ à $300\ nm$ ($0,02$ à $0,3\ \mu m$) \\ \hline
\textbf{Champignons microscopiques} & Eucaryotes (noyau, organites), paroi en chitine, hétérotrophes, unicellulaires (levures) ou filamenteux (moisissures). & \emph{Saccharomyces cerevisiae} (levure), \emph{Penicillium}, \emph{Aspergillus}, \emph{Candida}. & $3$ à $50\ \mu m$ (levures) \\ \hline
\textbf{Protistes (unicellulaires)} & Eucaryotes, mobiles (flagelles, cils, pseudopodes), autotrophes ou hétérotrophes. & \emph{Paramecium}, \emph{Euglena}, \emph{Amoeba}, \emph{Plasmodium} (paludisme). & $10$ à $200\ \mu m$ \\ \hline
\end{tabularx}
\end{center}

\textbf{3. TP N°5 : Observation au microscope optique (Savoir-faire)}
\begin{methode}[Protocole d'observation d'une goutte d'eau de mare / culture de levures]
\begin{enumerate}
    \item \textbf{Préparation} : Déposer une goutte de l'échantillon sur une lame (pipette). Poser délicatement une lamelle (angle $45^\circ$) pour chasser les bulles d'air.
    \item \textbf{Mise au point} : Objectif $\times 4$ ou $\times 10$ (grossissement faible). Regarder \textbf{par le côté} pour descendre l'objectif près de la lamelle sans la toucher. Regarder dans l'oculaire et monter lentement (vis macrométrique puis micrométrique) jusqu'à l'image nette.
    \item \textbf{Observation} : Passer à l'objectif $\times 40$ (si nécessaire, seulement vis micrométrique). Dessiner au crayon à papier, traits nets, sans hachures.
    \item \textbf{Dessin scientifique} : Cercle = champ microscopique. Échelle indiquée. Légendes alignées à droite (trait de rappel, pas de flèche). Titre informatif (nom, grossissement, coloration).
\end{enumerate}
\end{methode}
\begin{attention}[Pièges fréquents TP]
\begin{enumerate}
\item Ne jamais descendre l'objectif en regardant dans l'oculaire (risque casser lamelle/objectif).
\item Ne pas confondre bulles d'air (rondes, bord net, vide au centre) et microorganismes.
\item Grossissement total $G = G_{objectif} \times G_{oculaire}$ (ex: $40 \times 10 = \times 400$). Taille réelle $= \frac{\text{Taille sur dessin}}{G}$.
\end{enumerate}
\end{attention}

\textbf{4. Nutrition et respiration des microbes}
\begin{itemize}
    \item \textbf{Nutrition} : Mode d'approvisionnement en carbone et énergie.
    \begin{itemize}
        \item \cle{Autotrophes} : Synthétisent leur matière organique à partir de $CO_2$ (photosynthèse : cyanobactéries, euglènes ; chimiosynthèse : certaines bactéries).
        \item \cle{Hétérotrophes} : Prélèvent matière organique toute faite.
        \begin{itemize}
            \item \cle{Saprophytes} (décomposeurs) : Sur matière morte (champignons, bactéries du sol). Rôle écologique majeur : recyclage.
            \item \cle{Parasites} : Sur hôte vivant (pathogènes : virus, \emph{Plasmodium}, bactéries pathogènes).
            \item \cle{Symbiotes} : Association durable bénéfique (flore intestinale, lichens).
        \end{itemize}
    \end{itemize}
    \item \textbf{Respiration} : Oxydation de matière organique pour produire ATP.
    \begin{itemize}
        \item \cle{Aérobie} : $O_2$ nécessaire ($\ce{C6H12O6 + 6O2 -> 6CO2 + 6H2O + ATP}$). Ex: levures en présence $O_2$, bactéries aérobies.
        \item \cle{Anaérobie} (fermentation) : Sans $O_2$. Rendement ATP faible.
        \begin{itemize}
            \item Fermentation lactique : $\ce{C6H12O6 -> 2 CH3-CHOH-COOH + 2 ATP}$ (bactéries lactiques, yaourt, muscles).
            \item Fermentation alcoolique : $\ce{C6H12O6 -> 2 C2H5OH + 2 CO2 + 2 ATP}$ (levures, pain, vin, bière).
        \end{itemize}
    \end{itemize}
\end{itemize}

\textbf{5. Reproduction des microbes}
\begin{itemize}
    \item \textbf{Bactéries} : \cle{Division binaire (scissiparité)} asexuée. Une cellule mère $\to$ 2 filles identiques (clones). Temps de génération court ($\approx 20$ min pour \emph{E. coli} à $37^\circ C$). Croissance exponentielle : $N = N_0 \times 2^n$.
    \item \textbf{Levures (champignons)} : \cle{Budding (bourgeonnement)} asexué. Un bourgeon se forme, noyau se divise (mitose), bourgeon se détache. Parfois reproduction sexuée (fusion de deux cellules haploïdes $\to$ zygote diploïde $\to$ méiose $\to$ asques/ascospores).
    \item \textbf{Virus} : \cle{Cycle lytique} (réplication virale dans cellule hôte $\to$ lyse et libération). Pas de division cellulaire. Étapes : Fixation $\to$ Pénétration (ADN/ARN) $\to$ Réplication/Transcription/Traduction (usine à virus) $\to$ Assemblage $\to$ Libération. Cycle lysogénique possible (intégration génome hôte).
    \item \textbf{Protistes} : Division cellulaire (mitose) souvent multiple (schizogonie pour \emph{Plasmodium}). Conjugaison possible (\emph{Paramecium} : échange matériel génétique).
\end{itemize}

\textbf{6. Contamination par les microorganismes (Voies de pénétration)}
\begin{center}
\begin{tabularx}{\linewidth}{|l|X|l|}\hline
\rowcolor{popblueL}
\textbf{Voie de pénétration} & \textbf{Description / Mécanisme} & \textbf{Exemples maladies (Cameroun)} \\ \hline
\textbf{Voie respiratoire} & Inhalation gouttelettes (postillons, aérosols) ou poussières contaminées. & Grippe, COVID-19, Tuberculose, Méningite, Rougeole, Coqueluche. \\ \hline
\textbf{Voie digestive} & Ingestion eau/aliments souillés (mains sales, mouches, fèces). & Choléra, Fièvre typhoïde, Amibiase, Hépatite A, Poliomyélite, Ver intestinal. \\ \hline
\textbf{Voie cutanée/muqueuse} & Plaie, piqûre (vecteur), rapport sexuel, transfusion, mère-enfant. & Paludisme (moustique), VIH/Sida, Tétanos, Bilharziose, IST (syphilis, gonorrhée). \\ \hline
\end{tabularx}
\end{center}
\begin{propriete}[Prévention (Barrières)]
Hygiène des mains (eau + savon / gel hydroalcoolique), eau potable (filtration, ébullition, chlore), cuisson complète aliments, vaccination (PEV : BCG, DTC, Rougeole, Fièvre jaune, Hépatite B, Pneumo, Rotavirus, HPV), moustiquaires imprégnées (MILD), préservatifs, dépistage.
\end{propriete}
\end{bilan}

\begin{aretenir}[Checklist avant le BEPC]
\begin{itemize}
    \item $\square$ Je sais définir un microorganisme et donner sa taille ordre de grandeur ($\mu m$ vs $nm$).
    \item $\square$ Je distingue les 4 groupes majeurs (Bactéries, Virus, Champignons, Protistes) par leur structure cellulaire (procaryote/eucaryote/acellulaire) et donne un exemple de chaque.
    \item $\square$ Je sais réaliser et schématiser une observation au microscope (préparation lame/lamelle, mise au point par le côté, dessin légendé, calcul grossissement $G$ et taille réelle).
    \item $\square$ Je classe les modes de nutrition : autotrophe / hétérotrophe (saprophyte, parasite, symbiote) avec exemples.
    \item $\square$ J'écris les équations bilans de la respiration aérobie, fermentation lactique et fermentation alcoolique (réactifs, produits, ATP).
    \item $\square$ J'explique la division binaire des bactéries (croissance exponentielle) et le bourgeonnement des levures.
    \item $\square$ Je décris le cycle lytique d'un virus (fixation $\to$ libération) et pourquoi c'est un parasite intracellulaire obligatoire.
    \item $\square$ Je cite les 3 grandes voies de pénétration des microbes (respiratoire, digestive, cutanée/muqueuse) et une maladie exemple par voie.
    \item $\square$ Je propose 3 mesures de prévention adaptées à chaque voie de contamination (hygiène, vaccination, vectorielle, sexuelle).
    \item $\square$ Je sais exploiter un document (texte, tableau, courbe croissance bactérienne) pour en extraire l'information demandée.
    \item $\square$ Je rédige une conclusion argumentée en français correct, en utilisant le vocabulaire précis du cours (ex: "scissiparité" pas "division en deux", "bourgeonnement" pas "petit bout qui pousse").
\end{itemize}
\end{aretenir}

\findecours

\end{document}
